\documentclass[showtrims,		% Mostra as marcas de corte em cruz
			   %trimframe,		% Mostra as marcas de corte em linha, para conferência
			   11pt				% 8pt, 9pt, 10pt, 11pt, 12pt, 14pt, 17pt, 20pt 
			   ]{memoir}
\usepackage[brazilian,
			% english,
			% italian,
			% ngerman,
			% french,
			% russian,
			% polutonikogreek
			]{babel}
\usepackage{anyfontsize}			    % para tamanhos de fontes maiores que \Huge 
\usepackage{relsize}					% para aumentar ou diminuir fonte por pontos. Ex. \smaller[1]
\usepackage{fontspec}					% para rodar fontes do sistema
\usepackage[switch]{lineno} 			% para numerar linhas
\usepackage{lipsum}						% para colocar textos lipsum
\usepackage{alltt}						% para colocar espaços duplos. Ex: verso livre
\usepackage{graphicx}					% para colocar imagens
\usepackage{float}						% para flutuar imagens e tabelas 		
\usepackage{lettrine}					% para capsulares
\usepackage{comment}					% para comentar o código em bloco \begin{comment}...
\usepackage{adforn}						% para adornos & glyphs
\usepackage{xcolor}					 	% para texto colorido
\usepackage{array}
\usepackage[babel]{microtype}			% para ajustes finos na mancha
\usepackage{enumerate,enumitem}			% para tipos diferentes de enumeração/formatação ver `edlab-extra.sty`
\usepackage{url}					% para citar sites \url
\usepackage{marginnote}					% para notas laterais
\usepackage{titlesec}					% para produzir os distanciamentos entre pontos no \dotfill
\usepackage{textpos}

%\usepackage{makeidx} 					% para índice remissivo

\usepackage{edlab-penalties}
\usepackage{edlab-git}
\usepackage{edlab-toc}					% define sumário
\usepackage{edlab-extra}				% define epígrafe, quote
\usepackage[largepost]{edlab-margins}
\usepackage[semcabeco, 				% para remover cabeço, sobe mancha e mantem estilos
			]{edlab-sections}			% define pagestyle (cabeço, rodapé e seções)
\usepackage[%
			% notasemlinha 			
			% notalinhalonga
			chicagofootnotes			% para notas com número e ponto cf. man. de Chicago
			]{edlab-footnotes}


% Medidas (ver: memoir p.11 fig.2.3)
\parindent=3ex			% Tamanho da indentação
\parskip=0pt			% Entre parágrafos
\marginparsep=1em		% Entre mancha e nota lateral 
\marginparwidth=4em		% Tamanho da caixa de texto da nota laterial

% Fontes
\newfontfamily\formular{Formular}
\newfontfamily\formularlight{Formular Light}
\setmainfont[Ligatures=TeX,Numbers=OldStyle]{Linux Libertine O}


% Estilos
\makeoddhead{baruch}{}{\scshape\MakeLowercase{\scshape\MakeTextLowercase{}}}{}
\makeevenhead{baruch}{}{\scshape\MakeLowercase{}}{}
\makeevenfoot{baruch}{}{\footnotesize\thepage}{}
\makeoddfoot{baruch}{}{\footnotesize\thepage}{}
\pagestyle{baruch}		
\headstyles{baruch}

% Testes


\begin{document}


%\input{MUSSUMIPSUM}  		% Teste da classe
% Este documento tem a ver com as partes do LIVRO. 

\thispagestyle{empty}

\begin{comment}
\begin{textblock*}{2.625in}(0pt,0pt)%
\vspace*{-2cm}
\hspace*{-2.65cm}\includegraphics[width=138mm]{./ABERTURA.png}  
\end{textblock*}

\pagebreak
\blankpage

\thispagestyle{empty}
\end{comment}

%\blankpage

% Tamanhos
% \tiny
% \scriptsize
% \footnotesize
% \small 
% \normalsize
% \large 
% \Large 
% \LARGE 
% \huge
% \Huge

% Posicionamento
% \centering 
% \raggedright
% \raggedleft
% \vfill 
% \hfill 
% \vspace{Xcm}   % Colocar * caso esteja no começo de uma página. Ex: \vspace*{...}
% \hspace{Xcm}

% Estilo de página
% \thispagestyle{<<nosso>>}
% \thispagestyle{empty}
% \thispagestyle{plain}  (só número, sem cabeço)
% https://www.overleaf.com/learn/latex/Headers_and_footers

% Compilador que permite usar fonte de sistema: xelatex, lualatex
% Compilador que não permite usar fonte de sistema: latex, pdflatex

% Definindo fontes
% \setmainfont{Times New Roman}  % Todo o texto
% \newfontfamily\avenir{Avenir}  % Contexto
\begingroup\thispagestyle{empty}\vspace*{.05\textheight} 

              \formular
              \Huge
              \noindent
              \textbf{A árvore dos cantos}
              
              \vspace{0.3em}

              \noindent\large\textit{Ou o livro das transformações contadas pelos yanomami do grupo Parahiteri}
                    
\endgroup
\vfill
\pagebreak       % [Frontistício]
\newcommand{\linha}[2]{\ifdef{#2}{\linhalayout{#1}{#2}}{}}

\begingroup\tiny
\parindent=0cm
\thispagestyle{empty}

\textbf{edição brasileira©}\quad			 {Hedra \the\year}\\
\textbf{organização e tradução\,©}\quad	 	 {Anne Ballester Soares}\\

\textbf{coordenação da coleção}\quad		 {Luísa Valentini}\\
\textbf{edição}\quad			 			 {Jorge Sallum}\\
\textbf{coedição}\quad			 			 {Suzana Salama}\\
\textbf{revisão}\quad			 		 	 {Luísa Valentini e Vicente Sampaio}\\
\textbf{capa}\quad			 				 {Lucas Kröeff}\\
%\textbf{assistência editorial}\quad		 {Julia Murachovsky}\\

\textbf{\textsc{isbn}}\quad			 				 {978-65-89705-69-7}\smallskip

\vfill

\begin{minipage}{7cm}
\textbf{Dados Internacionais de Catalogação na Publicação (\textsc{cip})\\
(Câmara Brasileira do Livro, \textsc{sp}, Brasil)}

\textbf{\hrule}\smallskip

Pajés Parahiteri\\

\textit{A árvore dos cantos}. Pajés Parahiteri. Organizado e traduzido por Anne Ballester.\\ 
2.\,ed. São Paulo, \textsc{sp}: Hedra, 2025.\\

\textsc{isbn} 978-65-89705-69-7\\

1.\,Conto. 2.\,Literatura brasileira. \textsc{i}.\,Pajés Parahiteri. \textsc{ii.} Ballester, Anne. \textsc{iii.}\,Título.\\

\hfill \textsc{cdd}: 869.93

\textbf{\hrule}\smallskip

\textbf{Elaborado por Janaina Ramos (\textsc{crb} 8/\,9166)}\\

\textbf{Índices para catálogo sistemático:}\\
\textsc{i.} Conto: Literatura brasileira


\end{minipage}

\vfill

\textit{Grafia atualizada segundo o Acordo Ortográfico da Língua\\
Portuguesa de 1990, em vigor no Brasil desde 2009.}\\

\textit{Direitos reservados em língua\\
portuguesa somente para o Brasil}\\

\textsc{editora hedra ltda.}\\
Av.~São Luís, 187, Piso 3, Loja 8 (Galeria Metrópole)\\
01046--912 São Paulo \textsc{sp} Brasil\\
Telefone/Fax +55 11 3097 8304\\\smallskip
editora@hedra.com.br\\
www.hedra.com.br\\
\bigskip

Foi feito o depósito legal.

\endgroup
\pagebreak     % [Créditos]
% Tamanhos
% \tiny
% \scriptsize
% \footnotesize
% \small 
% \normalsize
% \large 
% \Large 
% \LARGE 
% \huge
% \Huge

% Posicionamento
% \centering 
% \raggedright
% \raggedleft
% \vfill 
% \hfill 
% \vspace{Xcm}   % Colocar * caso esteja no começo de uma página. Ex: \vspace*{...}
% \hspace{Xcm}

% Estilo de página
% \thispagestyle{<<nosso>>}
% \thispagestyle{empty}
% \thispagestyle{plain}  (só número, sem cabeço)
% https://www.overleaf.com/learn/latex/Headers_and_footers

% Compilador que permite usar fonte de sistema: xelatex, lualatex
% Compilador que não permite usar fonte de sistema: latex, pdflatex

% Definindo fontes
% \setmainfont{Times New Roman}  % Todo o texto
% \newfontfamily\avenir{Avenir}  % Contexto

\begingroup\thispagestyle{empty}\vspace*{.05\textheight} 

              \formular
              \Huge
              \noindent
              \textbf{A árvore dos cantos}

              \bigskip  
              
              \large
              \noindent
              \textit{Ou o livro das transformações contadas\\pelos Yanomami do grupo Parahiteri}
              \vspace{11.8em}
              
              \newfontfamily\minion{Minion Pro}
              {\selectfont\minion\small\noindent Anne Ballester (\textit{organização})}

              \bigskip

              \noindent
              {\selectfont\minion\small\noindent 2ª edição}

              \vfill

              \newfontfamily\timesnewroman{Times New Roman}
              {\noindent\fontsize{30}{40}\selectfont \timesnewroman hedra}

              \noindent{\selectfont\minion\small
              \noindent São Paulo \quad\the\year}

\endgroup
\pagebreak
	       % [folha de rosto]
% nothing			is level -3
% \book				is level -2
% \part				is level -1
% \chapter 			is level 0
% \section 			is level 1
% \subsection 		is level 2
% \subsubsection 	is level 3
% \paragraph 		is level 4
% \subparagraph 	is level 5
\setcounter{secnumdepth}{-2}
\setcounter{tocdepth}{0}

% \renewcommand{\contentsname}{Índex} 	% Trocar nome do sumário para 'Índex'
%\ifodd\thepage\relax\else\blankpage\fi 	% Verifica se página é par e coloca página branca
%\tableofcontents*

\pagebreak
\begingroup \footnotesize \parindent0pt \parskip 5pt \thispagestyle{empty} \vspace*{-0.5\textheight}\mbox{} \vfill
\baselineskip=.92\baselineskip
A Coleção Mundo Indígena reúne materiais produzidos com pensadores de
diferentes povos indígenas e pessoas que pesquisam, trabalham
ou lutam pela garantia de seus direitos. Os livros
foram feitos para serem utilizados pelas comunidades envolvidas na sua
produção, e por isso uma parte significativa das obras é bilíngue. 
Isso garante também a divulgação da imensa diversidade linguística dos povos
indígenas no Brasil, que compreende mais de 150 línguas pertencentes a mais de
20 troncos linguísticos.

Caso você tenha gostado do que aprendeu neste livro e queira usar alguma
história ou conhecimento, por favor entre em contato com a editora 
para pensarmos juntos com as comunidades. Lembre-se por favor
que mitologia, neste caso, não é questão autoral tampouco domínio público.\par

\endgroup

\pagebreak\thispagestyle{empty}\movetooddpage
{\begingroup\mbox{}\pagestyle{empty}
\pagestyle{empty} 
% \renewcommand{\contentsname}{Índex} 	% Trocar nome do sumário para 'Índex'
%\ifodd\thepage\relax\else\blankpage\fi 	% Verifica se página é par e coloca página branca
\addtocontents{toc}{\protect\thispagestyle{empty}}
\tableofcontents*\clearpage\endgroup}

\chapter{Nota da organizadora}

\textls[-5]{Este livro reúne histórias contadas por pajés yanomami do rio Demini sobre os tempos antigos, quando seres que hoje são animais e espíritos eram gente como os Yanomami de hoje. Estas histórias contam como o mundo veio a ser como ele é agora.}

\textls[15]{Trata-se de um saber sobre a origem do mundo e dos conhecimentos dos Yanomami que as pessoas aprendem e amadurecem ao longo da vida, por isto este é um livro para adultos. As crianças yanomami também conhecem estas histórias, mas sugerimos que os pais das crianças de outros lugares as leiam antes de compartilhá-las com seus filhos.}
\chapter{Como foi feito este livro}

\begin{flushright}
\textsc{anne ballester soares}
\end{flushright}

\noindent{}Os Yanomami habitam uma grande extensão da floresta amazônica, que cobre
parte dos estados de Roraima e do Amazonas, e também uma parte da
Venezuela. Sua população está estimada em 35 mil pessoas, que falam
quatro línguas diferentes, todas pertencentes a um pequeno tronco
linguístico isolado. Essas línguas são chamadas yanomae, ninam, sanuma e
xamatari.


As comunidades de onde veio este livro são falantes da língua xamatari
ocidental, e ficam no município de Barcelos, no estado do Amazonas, na
região conhecida como Médio Rio Negro, em torno do rio Demini. 
 

Em 2008, as comunidades Ajuricaba, do rio Demini, Komixipɨwei, do rio
Jutaí, e Cachoeira Aracá, do rio Aracá --- todas situadas no município
de Barcelos, estado do Amazonas --- decidiram gravar e transcrever todas
as histórias contadas por seus pajés. Elas conseguiram fazer essas
gravações e transcrições com o apoio do Prêmio Culturas Indígenas de
2008, promovido pelo Ministério da Cultura e pela Associação Guarani
Tenonde Porã.


No mês de junho de 2009, o pajé Moraes, da comunidade de Komixipɨwei,
contou todas as histórias, auxiliado pelos pajés Mauricio, Romário e
Lauro. Os professores yanomami Tancredo e Maciel, da comunidade de
Ajuricaba, ajudaram nas viagens entre Ajuricaba e Barcelos durante a
realização do projeto. Depois, no mês de julho, Tancredo e outro
professor, Simão, me ajudaram a fazer a transcrição das gravações, e
Tancredo e Carlos, professores respectivamente de Ajuricaba e
Komixipɨwei, me ajudaram a fazer uma primeira tradução para a língua
portuguesa.  

Fomos melhorando essa tradução com a ajuda de muita gente: Otávio
Ironasiteri, que é professor yanomami na comunidade Bicho"-Açu, no rio
Marauiá, o linguista Henri Ramirez, e minha amiga Ieda Akselrude de
Seixas. Esse trabalho deu origem ao livro\emph{ Nohi patama Parahiteri
pë rë kuonowei të ã} --- \emph{História mitológica do grupo Parahiteri},
editado em 2010 para circulação nas aldeias yanomami do Amazonas onde se
fala o xamatari, especialmente os rios Demini, Padauiri e Marauiá. 
 

Em 2013, a editora Hedra propôs a essas mesmas comunidades e a mim que
fizéssemos uma reedição dos textos, retraduzindo, anotando e ordenando as
narrativas para apresentar essas histórias para adultos e para crianças de todo
o Brasil. Assim, o livro original deu origem a diversos livros com as muitas
histórias contadas pelos pajés yanomami.  E com a ajuda do \textsc{proac},
programa de apoio da Secult-\textsc{sp} e da antropóloga Luísa Valentini, que organiza a
série Mundo Indígena, publicamos agora uma versão bilíngue das principais
narrativas coletadas, com o digno propósito de fazer circular um livro que
seja, ao mesmo tempo, de uso dos Yanomami e dos \textit{napë} – como eles nos chamam. 


 

 

 

 

 

 

\chapter{Para ler as palavras yanomami}


Foi adotada neste livro a ortografia elaborada pelo lingüista Henri Ramirez, que é a mais utilizada no Brasil e, em particular, nos programas de alfabetização de comunidades yanomami. 

\begin{itemize}
\item[/ɨ/] vogal alta, emitida do céu da boca, e que soa próximo a I e U
\item[/ë/] vogal entre o E e o O do português
\item[/w/] U curto, como em “língua”
\item[/y/] I curto, como em “Mário”
\item[/e/] vogal E, como em português
\item[/o/] O, como em português
\item[/u/] U, como em português
\item[/i/] U, como em português
\item[/a/] A, como em português
\item[/p/] como P ou B em português
\item[/t/] como T ou D em português
\item[/k/] como C de “casa”
\item[/h/] como o RR em “carro”, aspirado e suave
\item[/x/] como X em “xaxim”
\item[/s/] como S em “sapo”
\item[/m/] como M em “mamãe”
\item[/n/] como N em “nada”
\item[/r/] como R em “puro”
\end{itemize}


\part{A árvore dos cantos}

%\chapterspecial{O {surgimento} {dos} {pássaros}}{}{}
 

\letra{O}{grande tuxaua} que dividiu a terra onde nós moramos era um verdadeiro
líder, e morava além da parte central da terra dos
Këpropë.\footnote{  Trata"-se de outro nome dos próprios Parahiteri, narradores da história.}  Esse nome, um nome importante, também é o nome
da região onde ele nos ensinou a morar, é o nome da região onde eles
moravam antigamente.

Nessa terra, elefazia o \emph{kawaamou} e transmitiu esse ritual; ele os
mandava se reunirem e nos ensinou, assim, a nos
reunirmos.\footnote{  Em seu xapono.} 

Como se chamava aquele que dividiu a terra, quando, no início, os
Yanomami ainda não praticavam esse ritual; que, em seguida, se tornou
líder e deu a terra?Ele tem nome. Aquele que mostrou
o \emph{kawaamou} se chamava Gavião.

Ele existiu desde o início, esse verdadeiro líder que, depois, chegou
para morar com os poucos líderes que sobreviveram ao dilúvio. Gavião
morava em um xapono e o cunhado dele era o líder do grupo dos
Wãhaawëteri, ao mesmo tempo em outro xapono vizinho. Wãhaawë era o nome
dele. Os dois sobreviveram.

A região, da qual se apossou se chamava Wãha, por isso o nome dele era
Wãhaawë. Foram os ancestrais mesmo. Foi ele que dividiu e deu a terra
para nossos antepassados morarem. Aquele que deu a terra aos nossos
antepassados se chamava Wãhaawë. Os moços iguais aos nossos, que moravam
com ele, se chamavam também Wãhaawë. 

Moravam os dois, Gavião e Irara. Chamava"-se Irara aquele que ensinou os
Yanomami a atirar com zarabatana.

Nessa terra, quando não havia Yanomami, moravam somente três pessoas.
Era o xapono de Gavião, que bebia a água do rio Porena.Oxapono dele se
localizava na beira do rio Porena e bebia sua água. É o nome do rio onde
os dois moravam. Tomavam banho nessa água. A região Porena era a região
deles. Moravam lá, perto daquele rio. 

Nós, Yanomami, quase não aprendemos a soprar veneno nos outros. Já
mencionei esses nomes, Irara e Gavião, nomes terríveis que vocês já
ouviram. 

Agora é a história daquele que experimentou soprar. Foram esses dois que
ensinaram àquele feio;Irara conhecia o veneno, pois ele o possuía. 

Depois de ter observado Irara soprar veneno, o feio aprendeu e logo
soprou em Mel.Experimentou nele; aquele feio experimentou em Mel pela
primeira vez, quando o ato de matar não existia e quando ninguém
morria. 

Aquele antepassado chamado Irara recebeu o nome de Nokahorateri, o que
entende de veneno; o grupo dele também passou a se chamar assim. 

O feioquis logo atirar em Mel, pois pediu veneno a Irara. Ele o pegou,
quebrou um pedaço de semente para experimentá"-lo imediatamente. Esta é a
história do feio que soprou pela primeira vez o veneno. Não são os
nossos antepassados, são outros! 

 

Havia um líder muito bonito chamado Mel. Era nobre e muito cheiroso:
onde ele passava, deixava seu perfume, pois o cheiro dele era como o dos
cabelos das mulheres que os Yanomami acham tão bonitas que desejam
voltar ao xapono delas para se casarem.

Duas mulheres chegaram e viram Melroçando, as duas ficaram sentadas. Ele
ensinou a roçar, ele abriu a roça. Saíram do xaponodelase andaram uma
atrás da outra, até o xapono dele, só para vê"-lo. 

--- Olha só, esse homem é muito bonito! 

As duas pensavam isso por causa dos seus cabelos bem macios e
cheirosos. \emph{Hɨ̃a!,} fizeram as duas cheirando. Toda a floresta
estava cheirosa, o perfume dele se espalhava. 

No mesmo xapono, morava Mucura,\footnote{  Mucura é outro nome para o gambá, usado na região amazônica.}  o feio, que não
cheirava bem assim. Ao avistar as duas mulheres, cheias de desejo e se
dirigindoapressadas à casa de Mel, Mucura se zangou. Ele era mesquinho
e, assim, ensinará os feios a matar os Yanomami bonitos.Indiretamente, as
duas mulheres entregaram Mel à morte. 

Quando apareceram as duas mulheres, passando na frente da casa do feio,
ele estava deitado no chão. \emph{Xiri}! \emph{Kuxuha}! \emph{Krɨhɨ}!
Mucura bateu o chão com o pé para chamar a atenção das duas, e quando
elas olharam aquele feio fingiu cair no chão para elas cuidarem dele.

Apesar de ser fedorento, de estar sempre com olhos purgando e cheio de
feridas, quando ele viu as duas mulheres se abraçando, ele bestamente
pensou:

--- Elas chegaram para mim, minhas duas mulheres apareceram, gosto
delas. 

Ao perceber que o olhar das duas mulheres se dirigia à rede bonita de
Mel, ele se zangou mais ainda. 

--- Estou aqui! --- disse ele, mas sem efeito, pois elas não vieram por
sua causa.

A rede estava amarrada lá do outro lado, em um espaço livre, e as
franjas de sua rede vermelha, cor de sangue, se balançavam levemente, a
rede estava lá mesmo. 

--- É essa daí, lá mesmo, lá está a rede dele! --- disse uma das duas
mulheres, em pé de frente para a casa de Mucura. 

Enquanto isso, Mucura fingia fazer o \emph{wayamou}, apesar de não saber
fazê"-lo: 

--- \emph{Aë, aë, aë}! Minha mãe, a comida que eu guardo no moqueador,
pegue a carne e dê às minhas duas esposas que chegaram, dê a elas
agora. 

Apesar de não serem esposas dele, ele falou assim. 

Ele deu um pedaço de carne da bunda dele às duas mulheres, que olharam
de soslaio e fugiram, por causa do mau cheiro. 

\emph{Kuxu!}, elas fizeram, cuspindo, \emph{kuxu!}

O mau cheiro afastou as duas, que foram em direção da rede de Mel, cuja
casa estava um pouco longe. Elas foram lá, deitar"-se juntas na rede
dele. Ao chegar em casa, Mel pensou,\emph{hɨ̃ɨ!}, e perguntou:

--- O que vocês duas estão fazendo? Não havia ninguém como vocês antes,
como vocês duas chegaram?

--- Viemos por sua causa! Venha! Venha! 

Agarraram"-no e o deitaram em cima delas. Deitaram uma de cada lado para
cheirá"-lo. O feio, que fazia o ritual de \emph{wayamou}, viu isso e
chorou. 

Ele se abaixou e as moscas seguiam seu mau cheiro. Ele estava cheio de
ura\footnote{  Isto é, bernes.}, as pálpebras estavam cheias de ovos de mosca.
Quase cego pela ira, ele rapidamente saiu, pois aqueles que guardavam o
veneno, Irara e Gavião, moravam perto. Ele pediu: 

--- Dêem o veneno! Vou experimentar! Quero fazer sofrer as duas mulheres
que chegaram e me fizeram sofrer. É verdade o que estou dizendo! Agora
vou experimentar! Dêem um pedaço! --- disse ele. 

Eles logo responderam:

--- Então pegue o meu! Experimente com este! Experimente! Tente! 

--- Mas com quem? --- perguntaram os dois. 

--- Não digam que não sabem! Vocês já conhecem aquele feio Mel! Sabem
como ele não presta! É o nome daquele que não presta! 

Apesar de ser menos que ordinário, de ser careca, apesar de ter um rosto
feíssimo, ele disse isso! O nariz dele era horrível, e mesmo assim ele
disse isso. O olhar dele era horroroso, e mesmo assim, ele disse isso. 

--- Verdade, mesmo? 

--- Sim! Verdade! Hoje vou experimentar!

--- Se você diz isso, então tome! 

Quando voltou, ele rapidamente se preparou. 

Querendo impedi"-lo, a mãe de Mucura falou: 

--- Querido! O que você está querendo fazer? Deixe para lá! Deixe para
lá! --- disse ela --- Cuidado, querido! O que você está querendo? O que
você está pretendendo fazer? Cuidado! Não vai, não, não vai! Fique
quieto! --- disse ela --- Não vá lá de novo, com Irara e Gavião! São as
mulheres de Mel! São mulheres dele! Não pense que elas vieram para você!
--- disse a mãe --- Elas fecharam o nariz por causa do seu mau cheiro!
--- disse.

Pois seu mau cheiro chegava longe.

--- Não pergunte o porquê! Eu vou! 

Melhavia voltado para derrubar as árvores grandes. Mucurajá queria matar
Mel. 

Mucura correu e logo foi, enquanto zoava a queda das grandes
árvores. \emph{Kou! Kou! Kou! Kou! Ãaaaaaõoooo! Aë, aë, aë,
aëëëëëëëëëë!}, estrondava o eco da voz bonita de Mel. 

Para experimentar a força do veneno, Mucura o experimentou primeiro em
Lagartixa. As lagartixas, que ficam grudadas aos jutaís, aquelas que
sempre sobem. Mucura foi à direção de Lagartixa, enquanto Mel
derrubava. 

Onde se erguia um jutaí, Lagartixa subia, fazendo \emph{tararararara}! 

--- \emph{Kuxu! Kuxu! Kuxu}! Vou soprar naquele mesmo! Eu vou
experimentar a força do veneno! --- disse Mucura. 

Primeiro ele soprou Lagartixa. Soprou. Pegou de raspão na garganta de
Lagartixa. \emph{Paha!}E uma última vez, \emph{Paha! Paiii!}

Com isso, ele o feriu, e onde arrancou um pedaço da garganta, deixou
vermelhas as gargantas das lagartixas. Mucura ficou olhando, para ver
Lagartixa passar. 

O veneno não agiu. O veneno era fraco. Somente as folhas de jutaí,
tontas de veneno, caíram. Lagartixa, que Mucura havia soprado, subiu
mais acima. Mesmo tonto, ele não caiu. Ele se recuperou. 

--- Puxa! Não faz isso comigo! 

Mucura não gostou de ele ter resistido ao veneno. Lagartixa sumiu, já
estava em outro lugar. Mucura, depois de soprar o veneno que tocou
levemente a garganta de Lagartixa, em direção ao bonito Mel, que
trabalhava e não agüentaria a força do veneno: o fôlego dele não seria
tão forte. 

Mel estava virado em cima do andaime. \emph{Paha! Paha!} Ele se
desequilibrou. Aquele que foi soprado antes pegou um atalho e foi falar
a Mel: 

--- Tome cuidado! Ele me soprou! Aquele feio fedorento quis primeiro
experimentar a força do veneno comigo! Cuidado, fique atento! Ele vem em
sua direção! --- disse ele. 

--- \emph{Hɨ̃ɨɨ}! A força do veneno não me destruirá! O feio não me
pegará! --- Mel disse. 

O feio já estava perto.

--- Ele está me desafiando! --- disse Mucura. 

Mel virou, enquanto derrubava com o machado. \emph{Paha! Paha! Paha!}

Mucura simplesmente soprou; você não agüentaria muito tempo a força
desse veneno, não dá para agüentar. 

Mucura foi logo embora e Mel caiu, tonto. Ele caiu. \emph{Sarai}! Aquele
que tinha caído voltou para suas duas mulheres. Quando ele chegou, não
demorou…

 --- Estou com muito frio! Estou morrendo de frio! --- chegou em casa,
delirando, tonto. --- Faça fogo para mim! --- disse ele. 

Ele já estava morrendo. Enquanto as duas cuidavam do marido,
esfregando"-o, ele caiu morto. 

Ensinando o choro aos Yanomami, elas duas logo choraram, choraram,
pensando nele. As duas se abraçavam, chorando, se segurando pela mão,
chorando. Onde estava o marido morto, as duas mulheres avançam e
recuavam.\footnote{  Trata"-se de um movimento da dança feita em ocasiões fúnebres.}  Aí o feio se juntou, aquele que havia
soprado fingia chorar, de medo fingia chorar, e disse assim: 

--- Meu grande e querido sobrinho! Meu sobrinho, mesmo! Mataram meu
sobrinho! Mataram meu irmão! Puxa vida! Reapareça para elas! ---
fingia. 

Enquanto falava assim, as duas mulheres fugiram, não queriam escutar.
Ele as seguiu. Enquanto ambas as mulheres rodavam ao redor do xapono,
chorando, ele falava, fingindo chorar; rodeava ao redor do
xaponodissimulando, e as duas mulheres choravam, fugiam dele de novo,
pois estavam bem zangadas. Apesar de ele chorar, as lágrimas não saíam,
ele simplesmente fingia. 

Ele fugiu. As duas mulheres cremaram o corpo de Mel e, enquanto o
cremavam, o feio fugiu para se esconder, como se fosse escapar: fugiu
pensando poder se esconder. 

Depois da fuga dele, ensinando aos Yanomami a matar, a seguir os
rastros, eles seguiram logo os rastros. Procuraram os rastros.
Procuraram os rastros. Procuraram. Fizeram isso logo, e não perderam os
rastros. 

--- Foi o feio que o matou! --- diziam --- Foi mesmo o feio, o feio
fugiu! Agora é a vez dele! --- esbravejavam --- Nunca sobreviverá,
criminoso! Vamos dar o troco! Acabou com nosso líder, então vamos
matá"-lo! --- disseram, e seguiram os rastros. 

As formigas \emph{Tokonari}, as formigas \emph{Xĩriana}, as
formigas \emph{Mamisipreima} e as saúvas se mexeram, seguindo os
rastros. Os {kiritari} também seguiram os rastros. 

Eles chamaram {Resimaritawë}, que seguiu os rastros como cachorro,
pelo cheiro; eles passaram levando Resimaritawë. Não conseguiam
encontrar os rastros onde havia pedras, onde havia a
pedra \emph{maharixitoma}, porque os rastros já tinham sumido. 

Naquele lugar, onde Mucura havia matado Mel, os seus rastros estavam
dando voltas, sumindo no meio de um pântano. Ele fugiu e estava a uma
distância igual àquela que nos separa do rio Maupuuwei, aonde
vamos {caçar em grupo}. 

Mucura se escondia na montanha, ele foi lá em cima, porque queria
escapar. Ele subiu em uma árvore. A montanha era redonda como um jutaí,
ele entrou lá, onde a montanha tinha uma fenda. Pretendia se trancar
ali. Eles derrubariam a montanha para pegá"-lo. Mucura não entrou no
jutaí. Ele entrou nessa montanha porque queria fugir. 

Como criminoso, ele se arranhava, mas não com as
unhas.\footnote{  A pessoa que assassina alguém costuma arranhar"-se com as unhas,
castigando"-se. Mucura é mau, então não chega a se machucar.}  Adormeceu. Criminoso. Ele se admirava, porque
matou. Ele ensinou a ser criminoso. Ele perpetuou o crime. Ele estava
dormindo. 

Resimaritawë seguiu o fio de Mucura, e ficou escutando. 

--- Aqui! Ele está dormindo --- pensou.

Como criminoso, Mucura guardava um tipo de canudo no braço, guardava um
em cima da orelha. \emph{Õooo, õooo, õooo!} Escutava"-se a respiração de
Mucura e Resimaritawë ouvia:

--- \emph{Õooo, õooo, õooo!} Mel! Mel! Você viu como é bom! Foi isso que
fiz para você! --- roncava Mucura. 

Apesar de estar roncando, ele se gabava. 

--- Está vendo? Foi o que fiz pra você! Mel, foi o que fiz para você!
--- dizia, roncando.

Ouvindo isso, Resimaritawëlogo se assustou e gritou: 

--- \emph{Hɨ̃aaaaaaaaaëë, aaaaaaëë}, o Mucura fedorento, aqui, o feio em
pessoa está se gabando na minha frente, \emph{aaaaaë}! --- disse ele. 

Esse abrigo estava na base da montanha, que vocês não conseguiriam
derrubar rapidamente, mas eles trabalharam como loucos e causaram um
impacto incrível no cume da montanha. Não tinham terçados, mas mesmo
assim conseguiram matar Mucura. Apesar de não possuírem terçados como os
dos \emph{napë,} eles conseguiram matá"-lo. 

Chamaram os do grupo {dos tucanos} Parawari, porque o grupo das
Maitacas não conseguiam. Os Araris estavamtendo dificuldades com seus
machados de pedra, que se destruíam. Apesar dos machadinhos
dos \emph{Tokorari}, todos sofriam por causa das ferramentas, que se
quebravam em pedaços e não entravam na pedra. Chamaram os do grupo
Parawari, que viviam agrupados na baixada. 

Devido ao que fez o feio ao seu sogro, o bonito Mel, seu genro Parawari
disse:

--- Vamos! Vão chamar meus pais, que moram com meus avós: eles moram bem
perto e têm verdadeiros facões! 

--- Vai, querido, corra! Vai você! Vai buscar! Vai buscar! --- diziam
assim. 

Eles os buscaram. Eles chegaram e atacaram a montanha. Não queriam
destruir o facão deles. Como outros facões haviam sido destruídos, os
pedaços das lâminas estavam espalhados no chão. Para poupar esforços
inúteis, eles amoleceram a parte interna da montanha, como se fosse uma
árvore, com a força do pensamento. 

\emph{Kraxi! Kraxi! Kraxi!} Depois de amolecer a pedra, derrubaram uma
parte. \emph{Kraxi, kraxi, kraxi!} Atacaram logo de todos os lados para
fazerem voar lascas de pedra. A montanha era do tamanho de uma sumaúma.
Fizeram outro buraco grande, para poderem continuar com a destruição da
pedra. O buraco no tronco se aprofundava; fizeram uma espécie de cratera
na pedra. 

--- Vejam essa montanha! --- disseram os Parawari. 

Conseguiram fazer esse buraco porque os Parawari têm esse bico mais
comprido e, assim, será mais comprido para sempre. Os terçados, sendo
mais compridos, começaram a derrubar a montanha. O bico do tucano
empoleirado é, na verdade, seu terçado. Ele nãonos morde, quando olhamos
para ele?\emph{ Kraxi, kraxi, kraxi!} Foi nesse momento que começaram a
fazer a montanhabalançar. Não demoraram. \emph{Kru tu tu tu tu tu}!
Pedacinhos de pedra voavam e pulavam; os pedaços se espalharam no
chão. \emph{Huãaaaaaa}! Eram muitos, e estavam tristes por estragar seus
terçados, pois o bico deles é curto. Era isso mesmo: o trabalho começava
a ser realizado. \emph{Kru tu tu tu tu}! Continuavam trabalhando.

Só faltava o coração da montanha. Os terçados sendo curtos, esse
pedacinho ainda resistia. Apesar de pequeno, o tronco da montanha não
quebrou rapidamente.

--- \emph{Pei kë! Aaaaaaaaooooo!} Vamos! 

Apesar de o tronco estar quase torado, a montanha não estava se
mexendo. 

--- \emph{Aaaaaëëëë!} Vai você! A ponta da montanha vai nessa direção
--- disseram a Preguiça. 

Queriam que ele puxasse a ponta da montanha, pois a montanha não caía.
Preguiça esticou um fio flexível e puxava a montanha. Os terçados dos
Parawari, cortando a pedra, não atingiam o núcleo. 

Preguiça esticou uma espécie de fio. Elenem pensou: ``Essa montanha vai
me machucar!''.Ele não tem o costume de ficar com medo. Com o fio,
parecido com linha de pesca, ele puxou o cume da montanha. Não havia
nada para conter a pedra. Preguiça estava sozinho, sem apoio, mas também
sem medo. \emph{Ku tu tu tu tu tuuuuuuu!} A pedra começou a estourar,
fazendo um barulho enorme; parecia cair um pedaço grande de céu. 

Mucura estava preocupado e chorava, olhando para baixo com lágrimas
escorrendo. Ele estava desesperado. 

\emph{Tuuuuuuuuuuuu, tẽẽẽrërërërë!} Preguiça caiu mais à frente, para
não acabar embaixo da pedra, que caía. \emph{Ãaaaaaa!} O fio o
impulsionou e ele foi cair bem longe, ele se engatou lá longe com suas
garras. A pedra, quebrando e levando árvores, não alcançou a árvore onde
Preguiça se agarrou; se ficasse pendurado mais perto, ele se
machucaria. 

Eles destruíram Mucura. 

Todos os animais, as araras, os tucanos, os urus, os inambus, os mutuns,
os jacus, os urumutuns, os mutuns"-de"-traseiro"-vermelho e os jacamins
eram como Yanomami. 

Todos juntos pegaram aquele que foi destruído. 

Chegaram até o sangue de Mucura derramado no chão para se pintarem. As
pequenas maitacas amarelas se pintaram com seu sangue cinzento. Todos os
passarinhos são diferentes: uns são vermelhos, outros cinzentos, outros
têm pálpebras cinzentas. As cores dos pássaros vêm daquele momento,
quando se transformaram em animais, naquele mesmo lugar,com o sangue
derramado. Transformaram"-se onde havia o sangue derramado. 

Aqueles que foram destruir Mucura não voltaram para seus xaponos, onde
havia roças. Não voltaram para recuperá"-las. De tanto se pintarem com o
sangue derramado, ele acabou e, depois de terminarem, logo voaram. Logo
se transformaram. Eles ocuparam a floresta toda, não restou nenhum
espaço. 

Outros pintaram de preto o seu peito, passavam um pouco de sangue na
garganta, outros pintaram as pálpebras, outros, os cabelos, outros
pintaram o cume da cabeça; outros se pintaram de cor cinzenta,
derramaram os miolos e os excrementos com os quais se pintaram. Assim
fizeram. 

Preguiça, que puxou a pedra, passou no seu corpo os excrementos, por
isso ele tem cor cinzenta, é a cor dos excrementos de Mucura. Ele também
passou sangue na bunda, levemente. 

Assim, depois de esgotarem o sangue de Mucura, eles logo voaram, e todos
sumiram. 

Termina assim essa história, porque eles voaram, se transformaram e
partiram. Onde aconteceu essa história, outra segue. Contamos o que
aconteceu a Preguiça e Mucura, cujo esconderijo foi derrubado.

A região onde eles moravam tem um nome. Aquele homem bonito, que foi
morto, morava à frente da serra Moyenapɨwei. Ele nasceu, morava à frente
da serra Moyenapɨwei.  Essa serra se chamava Moyena.
Ao pé dessa serra, Mel abriu roças. Por isso, se chamaram assim, pois
ocupavam a região desse nome. 

Apesar de ser uma serra, chamava"-se assim. Tinha esse nome, pois era uma
região bonita; quando as árvores \emph{moyena} floresciam, as exalações
tomavam toda a floresta. Ainda existe o perfume onde ficava sua moradia,
e seus descendentes ainda moram lá. Chamam"-se Moyenapɨweiteri. Ficaram
morando lá, pois os antepassados se chamavam assim. 

Nessa região central, morava também Jacaré. Jacaré moravanessa região
central, onde Mucura matou Mel.

%\chapterspecial{{Naroriwë} {mucura}}{}{}
 

 

Ɨhɨ të pë rë kuaanowei, hei pata a yai rë kui a urihi rë wëyënowei,
urihi a rë kui, urihi kamiyë pëma kɨ rë përɨaɨwei a rë wëyënowei, pata
përɨamɨ yai, përɨamɨ a rë kuonowei, Këpropë kɨ urihi mɨ amo hamɨ he tore
ha a kuoma. Pei a wãha rë kui, pata yai, pata pë rë përɨaɨ rë hira
hërɨɨwei, pë rë përɨanowei, a wãha urihi, a wãha kua. 

Urihi a rë kui, a rë hẽreamonowei, kawãa të rë takenowei, kawaamou pëma
të tapë kurenaha, a hirapë, a rë wëkenowei, hapa Yanomamɨ të pënɨ hẽrea
a pou maopehe tëhë, a pou mao ma makuhei, a patamoa rë notirayonowei,
ɨ̃hɨnɨ pata urihi a rë hipënowei, weti naha a wãha kuoma? A wãha kua:
Yanomamɨ të pë hẽrea rë hipëkenowei Koparikesi wãha kuoma. 

Hapa Koparisi yai përɨkema, përɨamɨ yai, ɨ̃hɨnɨ të pë kãi notikema. Pata
të pë rë pakakumanowei, të pë wãisiapɨ rë hëprarionowei, ɨ̃hɨ Koparisinɨ
të pë kãi përɨkema. Hei Koparisi rë kui hei a përɨa, hei Koparisi, hei
heri e rë kui përɨamɨ Wãhaawëteri e kuoma. Ɨhɨ a wãha rii kuoma
përɨamɨ. 

Pe heri xo kɨ rë hëpɨprarionowei, kama urihi Wãha rë yukenowei Wãhaawë a
wãha kuoma. Pëma kɨ no patama rë përɨamaɨwei a urihi rë hipënowei, hapa
pata kɨ yai. Kamiyë yama kɨ no patapɨ yai, hei kamiyë yama kɨ no patapɨ
iha a urihi rë hipëke hërɨnowei, Wãhaawë a wãha kuoma.Hei ipa huya pë rë
kurenaha pë kãi rë përɨonowei, kama Wãhaawëteri e pë wãha kuoma. 

Hei kama Kopari, Hoari xo kɨ rii përɨpɨoma. Hoariwë a wãha kuoma.
Yanomamɨ të pë horaɨ rë hiranowei ɨ̃ha a kuoma. 

Urihi a hamɨ Yanomamɨ të pë përɨaɨ mao tëhë, hei pë përɨoma, hei, hei,
hei. Ɨnaha pë kuoma. Hei Koparisi yahipɨ, hei kee xapono. Koparisinɨ
Porena u koama. Porena u ha xaponopɨ praoma. Ɨhɨ u koama. Kɨ rë
përɨpɨonowei u wãha. Ɨhɨ u yaruama. Porena urihi, kama urihipɨ Porena a
wãha rë kuonowei a yupoma. Ɨha pë përɨoma, ɨ̃hɨ u ha, hei pe heri pei a
urihi yai rë huwëponowei, mɨ amo yai hamɨ. 

Yanomamɨ të pë horayopë, ɨ̃naha të mori kuaanomi makui, hei ya wãha
yuprarema. Hoarinɨ, Koparisinɨ a kãi rë përɨkenowei, wãritiwë pëma të
wãha rë hiripouwei. 

Hapa totihi të mixiã rë wapanowei, ɨ̃hɨ të wãti yaia, ei kɨpɨnɨ të ha
hirakɨnɨ, a hiraɨ, a hiraa xoakema, ɨ̃hɨnɨ. Të xi wãrihiwë rë taɨwehei,
të xi wãrihiwë xomi rë taɨwehei, hei Hoarinɨ të yai tama, hɨrɨ a yai rë
pore. 

Ɨhɨ iha të kɨ taɨ ha tararɨnɨ, hei wãriti të rë hiranowei, ai a horaɨ
xoaoma, Yamonamowãro a rë horanowei, a rë wapanowei, wãriti tënɨ a
wapama, hapa, ai xëpraɨ kuo mao tëhë, ai të nomaɨkuo mao tëhë. 

Pata Hoari a wãha rë kuonowei, pata Nokahorateri a wãha yai rë
yukenowehei, Nokahorateri pë kuprou xoarayoma. 

Ɨhɨ tënɨ a horaɨ xoaoma, ɨ̃hɨ iha e kɨ nakarema. E yuaɨ xoarema, e mo ko
wai ha huyetirënɨ e wapaɨ xoaoma. Wãriti tënɨ a rë horanowei të rii. Hei
kamiyë yama kɨ no patama rë kui mai! Ai rii. Ai të rii. 

Ɨhɨ katehe a rë përɨkenowei, përɨamɨ a kuoma, Yamonamowãro. A no xi
hiraarema. Hei kamiyë kurenaha mai, mɨ amo hamɨ amoamo a yai kuoma. A rĩ
yai kẽterio totihioma, wa henakɨ hɨprɨo, wa hɨprɨo, wa huɨ ha, wa rĩ
warou rë kurenaha, a rĩ kuoma. Pei henakɨ rĩ, kuwë yaro, ɨ̃ha suwë
Yanomamɨ të pë ha riëhëaprahenɨ, ɨ̃naha të pë napë kõo kuo puhiohe yaro.
Ɨhɨ iha Yanomamo wãro iha të pë hirakemahe. Hirano të taprapë. 

Suwë kɨpɨ waropɨkema, a ha tapɨrarɨnɨ. Ɨhɨ a rë hikarimore ha, ɨ̃ha kɨ rë
ropɨawei, kama kɨ kupɨopë ha, a kuaaɨ yaro. Të pë hikarimou hiraɨ ha,
hikari a tama. Ɨhɨ kama kɨ xaponopɨ prapɨopë ha, kɨ ropɨopë ha,
kɨpɨkororoa xoakema. A tapɨrarema. A mɨpɨma. 

--- Kurahë wãro a riëhëwë rë totihi! 

Kɨpɨ puhi kupɨtarioma, henakɨ si yai ëpëhëoma, ɨ̃hɨ pei henakɨ rĩ ha.
Hĩa! Kɨ puhi kupɨɨ xoaoma, urihi a rĩ pata kẽteri haikiwë ha, a rĩ
hirakaama. 

Ɨnaha a rĩ kuwë ha, kama a rĩ kuwëmi makui. 

Naroriwë a rĩ kuwëmi makui, a ha huxutarunɨ, suwë kɨ no wayuawë kõpei
yaro, kɨ no xi harirapɨwë totihiwë ha, a huxutarioma. Të pë huxutou
hiraɨ ha, a nohi ohotatanomi, wãriti të pë makui Yanomamɨ totihitawë të
pë xëpehe, huxuo hiraɨ ha. A napë rë rurupɨkenowei të rii. 

Suwë kɨpɨ pëpɨtario ha, hei a rë kurenaha wãriti të rëpraoma. Të praopë
ha, suwë kɨpɨ hapɨa katitirayoma. Xiri! Kuxuha! Krɨhɨ! Mamikɨ të marë
tamaɨwei, ɨ̃hɨ të ha krɨhɨmonɨ, suwë kɨpɨ mamo xatiprakema. 

Të rë wãritii, të kãi xomi kea nokarayoma. Kama ihamɨ kɨpɨmi makui,
mohoti katehe a yai ihamɨ kɨpɨ kõpɨoma kure ha. 

--- Kamiyë iha kɨ kõpɨoma, ipa suwë rë kɨ pëpɨtario ë, ya kɨ no xi
hirapɨaɨ! 

A puhi xomi ha kunɨ, a rĩ hĩtari makui, a kãi krĩhipɨ, a kãi warapisi
makui, a kuaama, a huxutou puhio yaro, suwë kɨpɨ hãkɨkopɨma, ɨ̃hɨ naha
kɨpɨ kupɨaɨ ha, e kea nokarayoma. 

--- Hëyëmɨ, hei kë ya! --- e xomi kupɨma, kama ihamɨ mi makui, a nohi rë
yaipɨpore pëkɨ no aihiwë yawëtëtao horayoma, yĩɨ e kɨ kasi kɨ ĩyë wakë
xurixurimoma, kuopë hamɨ: 

--- Kihi yai, kihi pëkɨ yai! --- e kɨ kupɨɨ kãi upraoma. 

Kuɨ tëhë, e xomi wãyamoma, wãyamou taomi makui, e xomi wãyamou kãi
kepɨrayoma: 

--- Ae, ae, ae, napemi, hei nii ipa ya, ya e kɨ, kɨ rë ma, makepo, pore
a wa yãhi ha yurɨnɨ, hei ipa suwë kɨ, kɨ rë kõpɨpohe ha yãhi ta hipëaɨ
ayaonɨ a --- Kama ihamɨ mi makui a xomi kuma. 

Kama posi moko kero, kero. Ɨnaha të kuɨ ha, kɨpɨ mamo ha axëpraikunɨ pei
të rĩ ha kɨ rupɨrayo hërɨma, kɨpɨ mamo ha axëpɨpraikunɨ:

--- Kuxu! --- kahi u këxëɨprao hërɨma. --- Kuxu! --- kɨ kupɨ hërɨm., 

Kɨ yaxupɨrema, kihi Yamonamowãro pëkɨ hamɨ kɨpɨ katitipɨke herayoma.
Kama e nahi yawëtëa yaro, kɨ ha ukupii kuhurunɨ kɨ yakapɨpari hirayoma.
Kama a rë kui a kõpema. 

--- Hɨ̃ɨɨ! --- a puhi kutarioma --- Exi wahë të taɨ? Ɨnaha kuwë të kɨ
kupɨonomi! Weti ihamɨ wahë kɨ huimaɨ kuhe?

--- Kahë ihamɨ! Kahë pë napë kõo. Pë tararema yaro, a ta hapo! 

 

A hurihipɨa mɨ parema. Pei kɨpɨ parɨkɨ ha a makepɨpoma. Pei a rĩ ha a mɨ
hetutupɨrema. Kuaaɨ ha, wãriti të rë wayamore, të no preaama, a ɨ̃kɨma. 

A mohe poarioma, moo pë pata makui kama a rĩ hĩtari hamɨ, a kãi
motepɨoma, mamo kasi pë motepɨmou haikioma. Kuwë yaro a huxutarioma,
rope a hurayo hërɨma, a no yai rë hërɨpɨ hiraɨwei a ahetea yaro, a rërëa
xoakema, a pehi yuaɨ ha, e rope ha nakarënɨ a ha rërëikunɨ:

--- Pei, hërɨ a ta hio! Ya ta wapa, kamiyënɨ! Suwë kɨ hore kõpɨpohe ya
kɨ a no ta prepɨmapo! Ware ã no prepɨma kɨhë! Peheti re ya wã! Hei kuikë
rë ya wapaɨ xoao! Ãtahu kɨ ta hio! --- e kuma. 

Si ihehewë a wã kãi hupɨrema. 

--- Pei, kutaenɨ ipa a ta yurë! Hei anɨ, a ta wapa! A ta wapa! Wapëpraa!
Weti kë? e kɨ kupɨma.

--- Weti mai! Wahë a nohi taɨ waikire! Yamonamowãro! Wahë pë xami taɨ
waikire! Yamonamowãro pë wãha xami rë kure! 

Kama a kuhupɨmi makui të kuma, a maxixipɨ yaro mohekɨ wãritiwë makui, a
kuma. Hei pei he rë kui ha, pëma kɨ henakɨ rë kurenaha henakɨ kuwëmi
makui, he tapramou wãritiwë makui, a kuma, hũxipë kãi wãritiwë makui, a
kuma, mamo xatio kãi wãritiwë he parohowë makui, a kuma. 

--- Peheti rë kë?

--- Awei, peheti hei kuikë ya wapaɨ! 

--- Wã haɨ yaio, pei yurë hërɨ! 

E ha kõponɨ a xurukou xoa peheroma. 

Naroriwë nɨɨpɨ e ã hama, pë nɨɨnɨ a hore wasima: 

--- Xei! Exi të ha wa xurukou kure? Kuo pëtao! Kuo pëtao! --- e hore
kuma --- Mihamɨ, xei, exi wa të taɨ? Exi wa të taapraraɨ ayaa kure?
Mihamɨ! A hurɨhë, a huɨ ta yanɨkɨ taru! --- e hore kuma --- A hua
kõrɨhë. Kama të pë hesiopɨ! Kama të pë! Kamiyë ihamɨ të pë tao puhi kuɨ
ayao tihë! --- e kuma, pë nɨɨ --- E kɨpɨ hũkakɨ kahuaɨ ha! --- e kuma. 

Praha makui a rĩ warou yaro. 

--- Ɨhɨ exi të ha mai! Ma! Ya huɨ! 

Kihi a rë kayapamo rahari ha, kiha a haɨpraɨ puhiopë yaro. A haɨprapë! E
rërëkema. E hua xoarayo hërɨma, kayapa hi kɨ hõra pë ramapou tëhë. 

Kou! Kou! Kou! Kou! Ãaaaaaõoooo! Aë, aë, aë, aëëëëëëëëëë! Katehe e a
karëtoma. Të kɨ wayu wapaɨ ha, reha a wapama. Ãroko hi pë hamɨ reha pë
marë sutiowei, reha pë marë tuo xĩroaɨwei, ɨ̃hɨ a ha ukuikunɨ, a hõra
përao tëhë. 

Ãroko hi pata rë upraawei iha:

--- Tarararararara! --- reha e kupe hërɨma. 

--- Kuxu! Kuxu! Kuxu! Ya wapaɨ ta yaio kë. Ya hamiwë ta wapaxo! --- e
kuma. Ɨhɨ reha rë a horaɨ parɨoma. A horama. Ɨhɨ rehanɨ: 

--- Paha! Hɨtɨtɨ anɨ! Paha! Paiii! 

Õramisi yoamarema. Ɨhɨ re e uno. Ɨhɨ rë õramisi rë tɨhɨyëprarenowei
hamɨ, ɨ̃hɨ rë pë õramisi kuprawë, a yuo mɨma. Ɨhɨnɨ e kɨ hamianomi. E kɨ
okearema. Ãroko hi rë kui, e ko hi henakɨ porepɨ kea tahiarayoma. Ɨhɨ
kama reha a rë kui a rë horaɨwei e torere pe hërɨma. E nomaɨ makui, e
kenomi. Ɨha kama rë a harorayoma. 

 

--- A no huxuaɨ tikooma ta yaitanɨ --- e matarioma, reha. 

Yai ha e kua yaro. Ɨhɨ a ha taakɨnɨ, õramisi ha yoamarënɨ, katehe a yai
hikarimopë ha e ukua piyëkema. E kɨ no horawëapraɨ rë mai! Mixiã tiremaɨ
rë mai. 

A rerekeaprarou tëhë, Paha! Paha! Ɨha rë e rë yutuprore, hei reha a rë
horare, e ã hama. Ɨha rë hei rë e rë kui a rë horaɨwei a hamirayou yaro,
e he tiherikema. 

--- Pei, a ta moyawëiku! Ware a hore horahe! Naro a rĩ wãriti rë
hĩtariweinɨ, ware a hore wapaɨ parɨohe, miha kahë a ta moyawëiku! Kahë
wa napë hore huimahe --- e kuma. 

--- Hɨ̃ɨɨ! Ya yuo mai kë të! Wãriti anɨ ware a waɨ mai kë të! ---
Yamonamowãro e kuma. 

A wãti kua ma waikirayoi ha. 

--- A wã no hore huxuoma ta yai tanɨ ɨɨɨ!

A kuu tëhë, a rerekeaprarou tëhë, hãyokoma a tuyë taɨ tëhë: Paha! Paha!
Paha! 

A ɨnaharë, ɨnaha e kɨ horaa takema, wa no no tetepɨ ma rë mai, ɨhɨ
naxomi të kɨnɨ. 

Të pë no tetemaɨ tao marë mahei, ɨha e xi wãrihiprou xoarayoma,
wãrihiprarunɨ, e porepɨ kerayoma, a ha kerɨnɨ: Sarai! Hesiopɨ kɨpɨ
ihamɨ, ɨha rë a keheropë, a no tetenomi, ɨ̃hɨ e ha kõtaponɨ:

--- Pei! Hĩa ya rë saihia tikorɨhe ë, ya rë saihirɨhe ë! --- e kuuroma,
porepɨ --- Kaɨyë wakë ta yëpɨpa xë ɨ̃ɨ --- e kuma. 

A maɨ waikia hërɨa yaro. A hore huripɨɨ mɨ paa, huripɨɨ tëhë e
kepɨrayoma. 

Ɨhɨ Yanomamɨ të pë ɨ̃kɨɨ hiraɨ ha, e kɨɨ̃kɨpɨma, kɨɨ̃kɨpɨ xoaoma, puhi
wayuyopɨma yaro, puhi mɨrapɨɨ yaro, kɨ hãkɨkɨpɨaprou, kɨ mɨa kãi
ɨpɨhɨpɨaprarou, hẽaropɨ a nomawë përɨopë hamɨ, kɨ tikukupɨaprarou, kɨ
tirurupɨaprarou, ɨ̃ha wãriti e të xomi nikeropɨma, a ma horanowei e xomi
ɨ̃kɨma, e xomi kiriri ɨ̃kɨma, e ã piyëkoma:

--- Õasiwë të pata yaio! Xëtëwë të pata yaio! Xëtëwë a pata xamia përaru
haikë! Asi yai! Ipa e pata pëpɨtou ta kõro! --- a xomi uhuti kuma. 

E kuɨ tëhë kihamɨ e kɨpɨ rurayoma, të ã hiripɨɨ puhiomi yaro. Ɨhamɨ e
xomi tiporepɨa ha kõrɨnɨ, kihamɨ e xomi kuaaɨ kõoi maɨ tëhë, kɨ rupɨa
kõrayoma kihamɨ kɨ rii ɨ̃kɨpɨma, ɨnaha kɨ kupɨama, kɨ huxupɨtarioma yaro,
a ɨ̃kɨɨ makui, mapuu kɨ kãi hanomi, xomi kuu pëoma, a xomi ɨ̃kɨɨ pëoma. 

Kuo tëhë, e ha kuokuopo hërɨnɨ, hesiopɨ kɨpɨnɨ a yapɨpema, yaaɨ tëhë,
wãriti të rë kui a hãtoprario hërɨma, a no tokupɨma marë! A xomi
hãtoprario hërɨma. 

A ha tokurɨnɨ, Yanomamɨ të pë xëpraɨ hiraɨhe ha, të pë mayo nosi pou
hiraɨhe ha, të pë nosi yaupraa xoakema. 

--- Mayo taeiwei, mayo taeiwei, taeiwei! --- kuaaɨ xoatarioma. Mayo nosi
hikepou xoatariomahe --- Ɨhɨ rë a wãritinɨ a xamimararei kuhe. A wãriti
kua marë mai, a wãriti hore ma rë tokurɨhe. --- Pei! Kama rë ã! --- a
noã taɨ xoaomahe --- Pë unokãi përɨkei ha yamuku! Ɨhɨ pëma a no payeri
tapraɨ xoao. Kamiyë pëma kɨ përɨamɨpɨ hore xëprarɨhe! --- të pë kuma. 

A nosi yure hërɨmahe, a nosi ha yurë hërɨhenɨ, ɨ̃ha a napë kõkaprouhe,
tokonari pënɨ, xĩrianari pënɨ, mamisipreimari pënɨ, koyeri pë kãi, pë
kuaama. Kiritari pë kãi pë nosi yauama. 

Ɨhɨ pei mayo nosi yai rë hikepore, hiima kurenaha hũka kɨ yakëɨ rë mai!,
Resimaretawë a nakaremahe. Resimaritawë a hayuremahe, mayo nosi
katitipraimihe yaro. 

Maa ma pë hamɨ, maa ma pë kurenaha, maharixitoma pë rë kure hamɨ, mayo
tokua waikirayou yaro. Hei a rë xëprare, kihamɨ mayo rë tɨhɨyëmatiiiiii,
wawëri pë pata ha mayo tokurayoma. A rë tokurɨhe, kihi ai të pë rë
heniyomo pohori. Maupuuwei u rë papohori naha, kihi Naroriwë a hitëo
kupohori, pei makɨ pata ha, a heakaprou piyërayoma. 

A tokuu puhiopë yaro. Hii hi ha Naroriwë a tua piyëkema. Ɨnaha
komorekomore ãroko hi kurenaha ɨ̃hɨ pei makɨ pata komore xatiopë ha, a
rukërayo tayoma. A he ruketayoma. Pei makɨ tuyëmapë. Ãroko hi yai hamɨ,
Naroriwë a rukënomi. Ɨhɨ kɨ ha, a he rĩya ha ruonɨ, a rukërayoma. 

Unokãi hixehixemou xoaoma. A mikema. Unokãi. A mɨprou yaro, a xëprarema
yaro. Të pë unokãimou hiraɨ ha, a unokãimoma. Ɨha a mitaoma. 

Ɨha Resimaritawënɨ, kama mananaepɨ hamɨ, mayo nosi rë hikepouwei yɨmɨka
takema. 

 

--- Hëyëha! --- e puhi kuma, e yɨmɨka takema, a mitakema. 

Kana ã hãkipoma, ai ã huuporanɨ. Õoooo, õooo, õooo! E mixiã kɨ kãi
horeheoma, e yɨmɨka takema. 

--- Õooo, õooo, õooo! Yamonamowãro! Yamonamowãro! Wa të oarei kuhe,
ɨ̃naha pë rĩya tapraɨ kuhe! --- e hũhũrua kuma, pei e hũhũrua makui, a
noka hekaɨ ha. --- Ma rë kui! --- Ɨnaha pë rĩya tapraɨ ɨɨɨ, Yamonamowãro
ɨ̃naha pë rĩya tapraɨ ɨɨɨ! 

Hũhũrua kuma, kuɨ ha hĩri tarënɨ, e rarɨa xoarayoma ɨ̃ha, e kirirarioma
pei a wãha, e rarɨa xoarayoma. 

--- Hɨ̃aaaaaaaaaëë, aaaaaaëë, hëyëha Naro a rĩ wãti rë hĩtariwei hëyëha
wãriti a wãti makui ware a no ka hore rë heka yahi aaaaaë! --- e kuma
--- Hei yahi a rë kupe, ɨ̃hɨ kihi rë të ma pata koro kumopë ha, wa të no
tupraaɨ haɨo taopɨ rë mai! 

Ɨhɨ pei rë ma pë nanoka rë kuinapë pënɨ ma pë ha karomaɨhe ha, të pë
ãtahu nanoka pata rë taamaɨwehei naha, makɨ nanoka pata taamamahe, pë
siparapɨ kua ma mai rë, a xëpraɨ he yatirayomahe. Ɨhɨ të rë kui, a ha
xëprarɨhenɨ, a rë xëprarenowehei, napë pë mi makui, pë siparapɨ rë
kure. 

Parawari pë ha nakarehenɨ, ãrimari pë no preaaɨ yaro, arari pë makui pë
siparapɨ yahekioprou no preaoma yaro, pë poopɨ yahatomou no preoma,
pokorari pë makui hapa të pë waɨ tipuramaɨ no preomahe, rukëimi ha, pë
nakaremahe. Parawari pënɨ, pë rë hiraawei, pë hiraa rë pepiawei, pë
nakaremahe. 

Ɨhɨ hei wãriti tënɨ, katehe Yamonamowãro iha a sioha rë kuaawei, ɨ̃hɨ e ã
harayomahe: 

--- Pei! Hei wama kɨ rë kui, hayë pë yai ta nakatarɨhe, xoayë pë yai rë
hirare, miha hayë pë yai hiraa kupe, hëyëha, ahete ha pë hiraa kupe, ɨ̃hɨ
rë pë siparapɨ yai kua kure --- e kumahe. 

--- Pei! Pei! Pei! Oxei, rërëiku, ɨ̃hɨ kahë rë wa! Pë ta kõa! Pë ta
kõaxë! --- e pë kuma, a noã tamahe: sioha ai hei, tëëpɨ yupoma: 

Pë kõremahe. Kõrehenɨ, pë ha waroikunɨ, ɨ̃hɨ pënɨ e ma kɨ napë kea
xoakemahe. Kama pë puhinɨ, pë siparapɨ wãriao puhiomi yaro, pë no preaaɨ
maopë, ɨ̃hɨ ai pë siparapɨ rë wãriaonowei, të pë hemata no wai preprawë
ha, kama pë puhinɨ, ma kɨ huxomi pata ëpëhëprare hërɨmahe, hii hi
kurenaha, kama pë puhinɨ, makɨ huxomi pata ha ëpëhëpramarɨhenɨ, të kɨ no
pata tuyëwëapranomihe. 

--- Kraxi, kraxi, kraxi --- të kɨ atahu pata yutuhamaɨ piyëkopehe, e
makɨ napë pata kea mɨ hetutua xoakemahe. Ɨhɨ wãrimahi pë pata rë
kurenaha, të makɨ pata rii kuwë, kihi kë të kɨ posika pata yawakaɨ
haɨtao huimati, të kɨ posika pata yawëtou kuimati, hawë makayo kɨ ka
pata taamamahe. 

--- Kɨ mamo no ta yërehe --- e pë kumahe Parawari. 

Ɨhɨ të pë husi ma rë tirehetai, ɨ̃hɨ kama xoati pë siparapɨ hĩiprawë,
sipara rë pënɨ, ɨ̃hɨrë pënɨ pehi rë tuyënowehei, pë siparapɨ hĩihiwë. 

Hei kamiyë pëma kɨ yaro, Yanomamɨ pëma kɨ mamo yëo tëhë, mayẽpɨ a pao
ha, pei pëma pë husi wãha rë hiraɨ, sipara e pë, sipara rë a kepou kure.
Wa si wërema, wa si nohi no tapɨ rë mai! Të pë nohi ma rë kure. Kraxi,
kraxi, kraxi! --- ɨ̃hɨ hei rë të rë kutaruhe naha, të pehi wëkëama
hërɨmahe. No tetenomihe. Kru tu tu tu tu tu! Të kɨ ãtahu pata xɨrɨkamaɨ
piyëkomahe, të kɨ pata rë tihatomouwei heinaha kuwë, të kɨ ãtahu pata
ĩtatarayo hërɨma. Hũaaaaaaa! Hei siparapɨ wãriaopë teri të pë husi
onohowë no rë preprai, kihi të kɨ mɨ pata puruwë. 

--- Ɨnaha kë yai ë, kɨ wã kãi rë haɨma piyehei ë --- e pë kuma. --- Kru
tu tu tu tu! të pë pata kutaama. Ĩsitoripɨ heinaha kuwë të amo hõro wai
tikëa hërɨma, pë siparapɨ hawëɨ he yatioma, hei të wainɨ të kɨ pata
huwëpoma. Kuwë makui, këprou haɨtaonomi, kurenaha rë të pata:

--- Pei kë! Aaaaaaaaooooo! 

Pei të kɨ koro pata reiwë totihiwë rë a makui! Të kɨ pata hãweteo rë mai
kë! 

--- Aaaaaëëëë, e pë kuma. --- Pei kë kahë wanɨ, kihamɨ të pata ora këɨ
kuaa hërɨpë! 

Ihamariwë a noã tamahe. Ei të kɨ ora pata ɨpɨamaɨ puhimahe, të pata
hãweteproimi ha, të pata këɨ haɨami ha, ɨ̃hɨ hei rë të kɨ ha yõriyõri rë
a rë taare, a rë ututuare anɨ, Ihamariwënɨ të kɨ ora pata ɨpɨretayoma.
Pë siparapɨ hawëa he yatirayoma. Kihi të kɨ pata ma rë tuyëɨ
hërakɨrɨhei, kiha e të kɨ ora pata ututuretayoma. Ihamariwënɨ: 

--- Ware a xëprarɨ! --- të kuɨ ha maonɨ, të pë kirihou marë mai!

Ɨhɨnɨ kɨ yai ɨpɨrema. Kama yõriyõri enɨ, ihiya masita rë kurenaha, e të
no owawëmi makui kɨ ora pata ɨpɨre kirioma, xĩroxĩro të ha yami
yaiikutunɨ, xi kãi ha kirirɨnɨ mai! Ku tu tu tu tu tuuuuuuuuuuuu! Hawë
kihi të parɨkɨ pata ketayou të kɨ pata tẽreremotayoma të kɨ pata kë
hërɨɨ yaro. 

Naroriwë a mɨa no preoma, a ɨ̃kɨma, mapuu kɨ kãi kei mɨ tëapraroma, a
puhi õkii yaro. 

--- Tuuuuuuuuu, tẽẽẽrererere! 

Kihi Ihamariwë pehi rë yokëɨ nokarakiri, hëyëmɨ të keo maopë. 

--- Ãaaaa! --- A pehi rë yokëɨwei, kihamɨ të no preaaɨ a pehi keo
kurakiriiiiiiiiiiiiii, a yaupraye kirioma, pei imisinɨ, të wãti yaupraye
kirioma, kuprao tëhë, pei e makɨ pata rë këre, ai te hi kɨ pata rë
hayuyare, e te hi kɨ hawërayoma, ahehe hamɨ moi pei a rë yaukenowei
kirionowei, a xëkei. 

Ɨhɨ tëhë a rë wãriarɨhahei. 

Ɨhɨ hei yarori pënɨ, ara pënɨ, mayẽpɨ pënɨ, pokorari pënɨ, hororomɨ
pënɨ, paruri pënɨ, kuremɨ pënɨ, katara pënɨ, katauri pënɨ, yãpi pënɨ,
Yanomamɨ kurenaha të pë kuoma yaro, hei kurenaha a nohi warokemahe yaro.
Ɨhɨ rë ĩyë pë rë hɨprɨarahei të pë warokema. Ɨhɨ rë ĩyë pë wai praaɨ rë
marahei, ĩyë pë napë warokemahe. Ãrima pë wai makui, të pë wai rë hãrei,
axiaxi ĩyë pë wai rë hɨprɨanowehei ɨ̃hɨ rë të pë. Kiritamɨ pë makui të pë
yaitawë, ai të pë he wakë, ai të pë he axi, pë mamokasi kɨ kãi axi,
ɨnaha pë rë kure, ĩyë pë hɨprɨamahe, të pë rë kuonowei hamɨ, ɨ̃ha rë pë
xi wãrihopë ĩyë pë hɨprɨaɨ kurahei, ai kihi pë rë itorɨhe hamɨ, pë
hãtopɨ nahi mɨonomi. Ɨharë. 

Ai pë kõa mɨ ha yaparɨnɨ, pë ni kuopë ha, pë haropë mai! Ɨhɨ hei rë ĩëpë
rë yãarahei, hɨprɨarahei, haikorahei, ɨ̃ha rë pë ha waikiprarunɨ, pë yëo
xoaokema. Pë xi wãrihiprou xoarayoma. Urihi a haikia xoaremahe. Urihi a
hëprou rë mai! 

Ai të pë parɨkɨ ĩxi yahetiataroma. Të pë õramisi hãhɨataroma, mamo kasi
kɨ hɨprɨataroma, të pë henakɨ kãi hɨprɨataroma, ai të pë yaro he marë
rohorei, pei hẽoxipë wai hɨprɨamahe, pei xĩhipë makui, xĩhipë kãi hɨprɨa
haikiaremahe, xĩhi pë praonomi, ɨ̃naha të tamahe. 

Pei makɨ ɨpɨarewënɨ xĩhipë hɨprɨama, ɨ̃hɨ ũ kɨ marë rohorei, ɨ̃hɨ pei rë
xĩhipë, pei hẽoxipë hɨprɨamahei, pei kõhesi ha, ĩyëĩyë e të wai
hɨrɨkɨkema. 

Kuwë yaro ɨ̃ha ĩyë pë rë waikarahei, pë yëokema. Pë marayoma. Ɨhɨ të rë
kui, të maprarioma, ɨ̃ha kama pë rii huokema yaro. Ɨha kama pë xi rii
wãrihiprarioma yaro, pë rë huoi kuhe, ɨ̃hɨ tëhë të wãisipɨ maxi rii hamɨ,
ai të rii kua notia, ai. Ai të ã. Ei të rë kui, të rë tanowehei, të ã
kuprarioma. Ihamariwë. Naroriwë pehi rë tuyënowehei të ã. 

Ɨhɨ rë te he rii tikëa kure. Ɨhɨ rë a urihi ha, urihi a wãha pei pë rë
përɨonowei a wãha. A wãha kãi kua. A wãha urihi rë taponowehei. Pëma kɨ
urihipɨ wãha marë kuprai kurenaha a wãha kuopë ha, të pë përɨhɨoma. Ɨhɨ
kama katehe a rë xënowei, ɨ̃hɨ Moyenapɨwei a parɨkɨ ha, a përɨoma. Katehe
a yai rë takenowei, Moyenapɨwei a parɨkɨ ha a përɨoma. Ɨhɨ Moyena pei e
makɨ wãha kuoma. Ɨhɨ kɨ të pëpoma, Yamonamowãronɨ. Ɨhɨ kama pënɨ, kamanɨ
Moyenapɨwei të wãha yupomahe. Moyena a urihi poma. 

Hehu makui ɨ̃hɨ Moyenapɨwei e kɨ kuoma. Katehe a urihi yaro a tapoma. Ɨhɨ
naxomi a himo rarou tëhë, urihi a rĩ pata hirakaatima. A rĩ pata hirakaa
xoaa, kama a përɨo no kuopë ha. Ɨha kama a no hekama përɨa xoaa,
Moyenapɨweiteri pë wãha rii kua. Ɨha Moyenapɨwei teri pë përɨa hëa, kama
pë wãha kuoma yaro, pata pë wãha kuoma yaro. 

Ɨhɨ të mɨ amo ha, te he tikë ha, hëyëha Iwariwë a rii përɨoma. Ɨhɨ
Yamonamowãro a xamiano rë kure ha a urihi mɨ amo ha he tikë ha a rii
përɨoma. 

 

 

%\chapterspecial{A {transformação} {dos} {quatis}}{}{}
 

\letra{O}{s animais} moravam em xapono; os quatis, as cutias e as antas, as
queixadas, os cuatás, os beija"-flores, os passarinhos moravam em grupos
como nós.

Àquela época, a mesma metamorfose ocorreu com todos esses animais,
exceto o caititu. Ele não andava como nós; não morava como nós, ele
sempre andou como ele ainda anda hoje. É um animal, e sempre foi. Sempre
andou como animal, assim como os cuxiús, os inambus, as cutias
vermelhas, os veados roxos.

Havia cinco espécies de animais. 

Já existiam todos os animais que há hoje na floresta? Não, somente
esses. Os jabutis não existiam, não andavam, não existiam como animais e
nem como Yanomami.Nem os tatus. 

Vocês comem o quati, apesar de ele ser Yanomami. Os Yanomami se
transformaram em quatis. Tornaram"-se animais à época de Horonamɨ. 

Como eram Waika, eles se metamorfosearam. Eram os ancestrais dos Waika.
Quando foram de \emph{wayumɨ}, todos se transformaram. 

--- Querido! Meu nariz se rasgou! 

--- Meu nariz se rasgou assim também!

--- Õãaa! Xiri! Meu nariz arrebitou! 

Foi assim com todas as crianças. 

--- Õãaa! Xiri! Avô! Meu nariz também arrebitou! 

--- Õãa! Xiri! --- disseram todas as crianças ao avô delas. 

Transformaram"-se enquanto estavam de \emph{wayumɨ}. Os ancestrais Waika
se metamorfosearam. São os primeiros moradores; eles se transformaram.
Ocuparam toda a floresta. Não sobrou nenhum xapono em torno do qual não
vivam quatis. 

Esse rio grande, rio abaixo, do qual vocês comem muitos peixes, na sua
parte média, apesar de ser rio abaixo, dá para avistar a pelagem muito
vermelha dos quatis. Eles andam por lá. 

--- \emph{Fĩfĩfĩ}! --- dizem. 

As mulheresWaika moravam aqui nessa época. O xapono vizinho
dos Koteahiteri era o xapono dos ancestrais dos Oramisitarematari,
Pescoços Podres. Moravam aqui. Eles comiam carne demais, não paravam de
comer carne. Eram verdadeiros carnívoros, por isso o nome deles era
Oramisitarematari. Eles gostavam dos Koteahiteri, eram amigos deles. Os
vizinhos com os quais se visitavam mutuamente chamavam"-se
Oramisitarematari. 

Eles os convidaram para ver e pegar a árvore dos cantos. Eram xapono
vizinhos. É o nome dos ancestrais. Os vizinhos dos Oramisitarematari
eram os Waika. Os que se transformaram em animais era o grupo dos
Waika. 

Os quatis são Waika. Foram os antepassados dos Waika, que se
transformaram indo de \emph{wayumɨ}. O nariz quebrou e se arrebitou.
Transformaram"-se no meio da floresta. Nunca mais voltaram a morar em
xapono. Eles se transformaram. A imagem deles se alastrou por toda a
floresta, como também a imagem dos jabutis.

 

 

 

%\chapterspecial{{Të} {pë} {rë} {yaruxeprarionowei}}{}{}
 

 

 


Yarori pë rë kui, hei yaro pë hirapramoma, yaruxe, tomɨ, xama, Yanomamɨ
kurenaha të pë rë hiraonowei, warë, paxo, të pë rë hiraonowei, tẽxo,
kiritamɨ pë kãi hiraoma, hei kurenaha. 

Ɨhɨ tëhë kama të rë wãrihore, poxe pë rë kui, Yanomamɨ të pë rë kurenaha
a huɨ taonomi, a përɨaɨ taonomi, kama xoati a huɨ rë xoaonowei. Poxe
yaro kë a rë hunowei ɨ̃hɨ a wai mahu a xĩro huma. Poxe, wɨxa, hõrama,
xĩhɨna, prẽari. 

Ai, ai, ai, ai, ai, ɨ̃naha të pë huɨ kutaoma. Hei kurenaha urihi a
hɨtɨtɨoma, yaro pënɨ? Ma, hei të pë xĩro kuoma. Totori pë kãi kuonomi,
pë kãi huɨ taonomi, totori Yanomamɨ pë kuonomi, opo pë kãi hunomi. 

Hei yaruxe, Yanomamɨ a makui wama a waɨ, ei të pë rë kui Horonamɨ a kuo
tëhë, pë xi wãrihoma. 

Yanomamɨ të pë yaruxeprarioma, Waika kë pë. Kama Waika pë wãha kuoma,
Yanomamɨ. Yaruxe pë rë kui Waika pata pë përɨoma. Waika pata të pë
hiraoma. Ɨhɨ Waika pë rë hiraonowei, ɨ̃hɨ kama Waika rë pë rii yaro pë xi
wãrihoma. Ɨhɨ rë pëma pë Waika yaruxe waɨ, pë xi rii rë wãrihonowei.
Wayumɨ pë xi wãrihou haikoma. 

--- Õasi ya hũka kɨ nohi hëtɨa yairëhe! Ɨnaha ya hũka kɨ rii hëtɨa
kurayou kuhe. Ihiru hɨtɨtɨwë të pë kãi rë kui. 

--- Õaaa! Xiri! Xoape ya hũka kɨ rii hëtia kurayou kuhe. Õaaa! Xiri! ---
e pë kuma. 

Ɨhɨ wayumɨ rë pë xi rii wãrihiprarioma. Waika pata kë pë xi rii
wãrihoma. Xomaomɨ të pë, pë xi wãrihoma. Urihi kë a haikia xoaremahe. Ai
të kɨ yahipɨ xee hëama rë mare! 

Hei të u pata koro hamɨ, hei wama të u pëɨxokɨ no yuripɨ waɨ rë xoape,
të u koro pata makui hamɨ, yaruxe të kɨ wakë pata xopoi kunomai të kɨ
pata ma rë huɨ korowë piyëkei. 

--- Fĩfĩfĩ! --- pë kutoma

Ɨhɨ Waika kama nohi patama pë kuoma kutaenɨ pë wayumɨ huɨ xi ha wãrionɨ
pë xi wãrioma. Pë hũkakɨ ha kërarunɨ, pë hũkakɨ mɨ yaprekewë. Pë hũkakɨ
hëtɨtoma.

 

Hëyëha, Wãikayoma pë kãi hiraomahe. Wãikayoma pë kãi rë hiraonowehei.
Koteahiteri pë yahipɨ he tikëre hamɨ, hei kë pë yahipɨ.
Õramisitarematari pata kë pë yahipɨ. Hëyëha pë rii hiraoma. Ɨhɨ pë rë
hirare pë yai naiki he ropao totihioma. Yaro a waɨ no teteonomihe. Pë
yai naikioma. Ɨhɨ kutaenɨ Õramisitarematari pë wãha kuoma. Ɨhɨ
Õramisitarematari pënɨ Koteahiteri pë nohipomahe. Kama nohi e pë kuoma.
Pe he rë waroyouhe rë tikëkonowei Õramisitarematari pë wãha kuoma. 

Ɨhɨ rë pë iha amoa hi nohi kõapraɨ piyëkomahe. Kama pë yahipɨ he
tipëtëmoma. Pata të pë wãha yai. Oramisitarematari pë yahipɨ he tikë ha,
Waika pë xi rë wãrihonouwei pë rii hiraoma. 

. Urihi ha, pë xi wãrihia xoararioma. Pë përɨaɨ kõonomi. Ɨhɨ pë xi rë
wãrihiraruhe, pë no uhutipɨ huomopotayoma, urihi a haikiremahe. Urihi a
he tatohowë rë kurenaha, a urihi yaitawë makui, a urihi haikiremahe,
yaruxe pënɨ, totori pënɨ. 

 

 

%\chapterspecial{A {proliferação} {do} {fogo}}{}{}
 

\letra{E}{u vou contar} a história de Jacaré. Seus conterrâneos sofriam por causa
da escuridão à noite, porque não conheciam o fogo e, então, comiam cru.
Depois de apanharem frutas \emph{kaxa}, eles as comiam cruas, pois não
havia fogo. A região chamada Kaxana era a região de Jacaré. Ele ocupava
essa região. Ele bebia a água do rio Kaxana. Era o nome dessa região.
Ele comia as frutas \emph{kaxa} cozidas às escondidas, aquele que
detinha o fogo na sua boca. Essa região Kaxana fica no centro da
floresta. No meio dessa região, há o rio Kaxana. 

Os jovens que moravam com Jacaré estavam sofrendo e delirando por causa
da comida crua. Ele comia cozido sozinho. Ele não oferecia a eles. Até o
paladar das mulheres sofria com a comida crua. A história dos
companheiros de Jacaré, que, depois de refletirem, encontraram o fogo,
guardado por Jacaré, ocorre no meio da nossa história. 

Alguém acordado à noite ouviu o som baixinho daquele que mastigava
escondido comida frita. \emph{Kãrɨ, kãrɨ, kãrɨ}! Por causa de sua boca
emitindo esse som baixinho, um deles percebeu o que estava acontecendo. 

--- Será que ele está comendo algo frito?

Ele pensou assim, apesar de não ver lenha queimada no chão, como
acontece quando a gente acorda e, ainda deitado, olha para o chão. 

Depois de torrar os alimentos na sua boca, ele os comia com sua esposa à
noite. \emph{Kãrɨ, kãrɨ, kãrɨ}! Os dentes dos dois faziam esse som
baixinho. 

Depois de seus companheiros acordarem, falavam de Jacaré, baixinho. 

--- \emph{Hoaaaaaaaa}! Vamos, vamos procurar! --- disseram logo. 

Não foram nossos antepassados que descobriram o fogo. Nós não conhecemos
o fogo por ele ter aparecido de repente. Nossos antepassados sofriam,
eles endoideciam por causa da comida e pareciam doentes. Sofriam e
endoideciam por comer carne crua. 

Salvaram"-se com o fogo de Jacaré, que pegaram e espalharam. 

Eles viviam sempre tristes; o fogo quase não saía, quase não existia.
Viveriam sempre tristes se o fogo não tivesse existido e sido
repartido. 

Jacaré ensinou os \emph{napë} a fazer fogo; eles o usaram porque Jacaré
lhes ensinou, pois ele guardava o fogo na sua boca.

Depois de acender as brasas em um tipo de forninho, ele colocava a
comida nas brasas. A comida cozinhava escondida em cima das brasas; ele
fritava as caças e as frutas às escondidas, assim ele fazia. Ele
guardava o fogo com ciúme, quase não revelava o fogo. 

Todos se juntaram para descobri"-lo. Agruparam"-se, onde Jacaré comia
queimado e guardava as folhas dos embrulhos. Ele as enterrava, cobria"-as
com terra; cavava o chão levemente e colocava as folhas queimadas; assim
fazia. Sozinho, saboreava a comida cozida. É um jacaré, como a gente
diz.

Mas ele cedeu o fogo? Não, ele sovinava o fogo. Vocês só veem um pedaço
da língua dele por causa do fogo. É um jacaré. A língua não queimou à
toa, o fogo acabou com a língua, ficou somente um pedaço no fundo, pois
ele guardava o fogo na boca. 

Avisaram todos os xapono, convidando"-se a se reunirem. Chegaram e se
reuniram. Jacaré não ficou sabendo e saiu à procura da
fruta \emph{kaxa}, enquanto eles se reuniam. 

Onde se revelará o fogo, queriam se alegrar rapidamente com a comida
cozida; conseguiram encontrar lascas queimadas de comida. Revelou"-se o
fogo com Jacaré; os \emph{napë} e os Yanomami terão o usufruto do fogo,
por isso comemos cozido.

No início, quase não conhecíamos a comida cozida. Onde nossos ancestrais
comeram cozido? Onde comiam cru com sangue? Quase todos nós comíamos
cru! Você escuta as minhocas e os minhocões, eles estão
aliviados, \emph{Tɨ̃ɨɨ}! \emph{Tɨ̃ɨɨ}! Quase você comia esse tipo de
minhocões, quase você sofreria, assim como os antepassados sofriam. 

Reuniram"-se, como nós agora. Na reunião, proibiram Jacaré de sair. 

--- Vamos, avô! --- seus netos disseram. 

Um de seus netos, aquele tipo de calango que fala
assim, \emph{Serororo}!, aqueles calangos pequenos que caem na água e
que fazem assim, \emph{Serororo}! Nós chamamos de \emph{temoa}. Era um
dos netos dele, quando era Yanomami. 

--- Vamos, meu tio! Fique parado! Você vai escutar. Escute os que se
reúnem. Deixe para ir amanhã de novo à procura de comida! --- disse seu
neto. 

--- \emph{Haaa}! --- disse ele com uma voz rouca. 

--- \emph{Hai! Hai!} Por quê? --- ele não perguntou. 

--- \emph{Ho}! Está bem! --- disse. Assim disse logo. 

--- \emph{Hëëë}! --- disse ele, da mesma forma que diz hoje. 

Enquanto se dirigia a ele, ele ficou parado. 

--- \emph{Ho}! Sim! Vamos esperar logo! --- disse à sua esposa. 

Ela tem uma voz harmoniosa, como a que se escuta nas cabeceiras dos
igarapés. \emph{Pẽi! Pẽi! Pẽi!} Os que falam assim, na beira d'água, os
que sempre ficam dentro da água, os \emph{pëipëimɨ}. Era a esposa de
Jacaré. Apesar de ser esposa de Jacaré, era parecida com a perereca que
tem desenhos na coxa. 

Quando o neto disse aquilo, eles pararam, ele e sua esposa. \emph{Hɨ̃ɨɨ}!
Ele sabia de quê se tratava. 

--- Eu não revelarei o fogo que eu possuo! --- pensou logo. 

Ele se zangou logo, sem razão. Ele foi dormir, para não dizer nada
àqueles que estavam reunidos. \emph{Hɨ̃ɨɨ}! Ele afastou sua rede, virou
seu rosto para o outro lado e adormeceu. Não se mexia. 

Os outros queriam fazer rir Jacaré e se revezavam. Os \emph{kiritari} se
revezaram. \emph{Ão, ão, ão, ão, ão, õoooo!} Eles tentaram, mas ele não
se mexeu. Aquele que possuía o fogo, cujas brasas brilhavam, guardou a
boca bem fechada e parecia não ter boca. 

Cada pássaro ia no meio do xapono dançar, se revezando.
Os \emph{huimiri}, os \emph{taɨsarakari} e os galos da serra iam deitar
com ele, se revezavam. Os bicos"-de"-brasa se sentavam no chão, sentavam,
sentavam, mexendo as asas; queriam fazê"-lo rir e se revezavam. A esposa,
de costas, disse logo: 

--- Não se mexa! Fique aí quieto! --- disse ela.

Em seguida:

--- Não se mexa! Durma! --- Falava ela, do seu canto.

--- Vocês vão fazer assim! --- Jacaré não disse, apesar do barulho. 

Bem depois, eles cansaram. Ele não se mexia. Ficaram cansados, os
coitados se mexeram tanto, tiveram tantas dores, que começaram a
sofrer. 

Beija"-Flor não se levantou logo. Ele estava no centro, afastado com
Tohomamoriwë, irmão dele. Eram os dois: o mais velho, Beija"-Flor, e o
mais novo, Tohomamoriwë, daqueles beija"-flores pequeninhos. Beija"-Flor
agiu primeiro. 

--- Irmão maior, sua vez! Você tenta logo, eu vou olhar primeiro! ---
falou o irmão mais novo. 

Apesar de o irmão mais velho se mexer, Jacaré não reagiu.

--- \emph{Sĩo! Sĩo! Sĩo! Tõu, tõu, tõu!} --- disse.

Mexendo"-se à frente de Jacaré, ficando parado, ele colocou peninhas
brancas no seu ânus, nas suas patas, as peninhas estavam repartidas
igualmente nos dois lados do ânus. A língua de Beija"-Flor saía e os
olhos de Jacaré se abriram. Deu uma olhada e adormeceu de novo. 

--- Não com você! --- pensou. 

Beija"-Flor mexeu duas vezes e quando todos terminaram, ficaram cansados.
O irmão mais velho disse ao mais novo: 

--- Vai, irmão menor, tua vez, depressa! Com você é bem capaz de ele
rir.

Aí o irmão mais novo disse:

--- \emph{Waooo!} Olhem isso! Sou eu mesmo! Olhem para mim! --- disse
ele. 

Todo o mundo olhou. Bico"-de"-Brasa estava triste, sentado, como um
doente, aguardando o fogo para pegá"-lo logo. Aquele que estava esperando
se levantou e aproximou"-se. 

--- Bem, olhem só como eu faço: \emph{Sĩo! Sĩo!} --- dizia a certa
distância, de onde se levantou. 

Parecia o som do carapanã \emph{tëërëkë.}

--- \emph{Sĩo! Sĩo! Sĩo! Tëɨ! Tëɨ! Tëɨ!}

Um som bonito começava a ser ouvido e o dono do fogo arregalou os olhos,
olhou para o que acontecia.

Beija"-Flor logo se levantou, para dançar à frente de Jacaré.

--- \emph{Wão!} Você vai ver!

Ele ficou à frente de Jacaré, continuando a dizer:

--- \emph{Sĩo! Sĩo!}

Parecia voar para sempre. 

--- \emph{Tëɨ, Tëɨ, Tëɨ!} --- ele falava. 

Ele parecia pendurado com as pernas abertas. Ele fazia assim: a cauda
dele ficou virada para cima, e os dois irmãos ficaram um ao lado do
outro. Eles faziam como se fosse um ritual. 

Nessa altura, Jacaré se levantou. Ele estava sério, mas se endireitou e
sentou na rede, e deu uma gargalhada. Com as gracinhas de Tohomamoriwë,
ele entregou o fogo. Tohomamoriwë estava na altura dos olhos de Jacaré: 

--- Ho, ho, ho, ho! --- Jacaré riu. 

\emph{Prohu}! A brasa pulou. \emph{Tou}! Bico"-de Brasa esperava o fogo e
o apanhou. \emph{ɨɨɨ, krihi, krihi}! Ele apanhou o fogo, que queimou seu
bico, por isso o bico dele é vermelho, pela queimadura do
fogo. \emph{Hɨ̃}! Como o fogo era pesado, ele não conseguiu voar alto,
quase caiu de volta com o fogo. 

--- \emph{Xĩapo}! Disse ele.

Ele estava esperando, empoleirado mais em cima. Ele foi levar o fogo.
Wẽooo! Pois ele é maior. Tu tu tu tu tu tu tu! Levou as brasas em cima
da árvore abiorana murcha, ele as colocou na ponta do tronco da árvore. 

A mulher de Jacaré se levantou. Quando Bico"-de Brasa quase apanhava o
fogo, ela jorrou urina. Apesar de a mulher de Jacaré jorrar sua urina,
Japu levou o fogo mais alto e a urina não conseguiu alcançá"-lo. Enquanto
o fogo estava em cima e não apagava, ela não acertou o fogo, a urina não
alcançou. Ela disse:

--- Vocês pegaram o fogo, então vocês chorarão quando cremarem os seus
mortos, vocês sofrerão e chorarão pelos seus mortos cremados! 

Ela falou verdade, pois quando morremos, nos cremamos; ela disse a
verdade: nós praticamos a cremação e nos cremamos. Ela disse que era
para ser nossa tradição. 

--- Eu vou ao igarapé e ficarei feliz lá com meu marido para sempre ---
disse --- Vocês sofrerão com o fogo. Ele se tornará eterno. O fogo
derreterá seus olhos! 

É verdade o que ela disse, a esposa disse a verdade. Nós nos cremamos,
nossa carne queima, ela falou certo. Se isso não houvesse acontecido,
não nos cremaríamos. Dito isso, ela e seu marido, Kruxu! Kopou! Os dois
foram às águas para nada de ruim acontecer, para eles não ficarem
doentes, não pegarem diarréia, nem dor de cabeça, nem conjuntivite, nem
terem dor nas pernas, nem conhecerem a malária. E assim serão. 

Eles tem problemas de dentes como nós? Não, tem não! Não conhecem dor de
dente. 

Esses eventos aconteceram para que seja assim. Eles ainda estão
felizes. 

--- Pĩri! Pĩri! Pĩri! Pĩri! Pĩri! Pĩri! Pĩri! Pĩri! Pĩri! Pĩri! Pĩri!
--- cantou a esposa até chegar às águas. 

 Eles pegaram o fogo e os dois moram ainda nas águas.

%\chapterspecial{{Iwariwë}}{}{}
 

 

Kama Iwariwë hei pë rë kui, ya pë wãha rë tare, ruwëri pë no prepramoma,
kaɨ wakë tanomihe, riyëriyë yaro pë wamahe. Riyëriyë atayu pë wamahe.
Riyëriyë kaxa a pata yaxaamahe. Kaɨ wakë kuami yaro. Ɨhɨ Kaxana kama e
urihi ha, Iwariwë a rii përɨoma, Kaxana a urihi ha. A urihi rii poma.
Kaxana u rii koama. Urihi ɨ̃hɨ Kaxana a wãha kua. Rɨpɨrɨpɨ ɨ̃ha a waɨ
hãtooma, kaɨ wakë rë pore, kaɨ wakë rë titipore pei kahikɨ hamɨ. Mɨ amo
yai hamɨ, Kaxana a urihi kua, Kaxana u mɨ amoa. 

Hëyëha Iwariwënɨ kama e të pë hei kurenaha kama huya e pë kãi rë
përɨawei makui riyëriyë të pënɨ të pë xi wãrii no preo ayao hei. Kamanɨ
rɨpɨrɨpɨ yaminɨ a wama. Të pë toonomi. Suwë makui riyëriyë të pë ha, të
pë no kãi preaama. Hei të pë rë kuinɨ, yakumɨ të pë puhi ha tao hërɨnɨ,
Iwariwë iha kaɨ wakë he rë haaɨ rë piyërahei, hei mɨ amo ha të kua. 

Mi titi hamɨ harɨkano kë kɨ waɨ rë hãtoawei kahikɨ ĩsitoripɨ wai hõra ha
hĩria he harurehenɨ, ai të rë rãa he haruawei. Kãrɨ, kãrɨ, kãrɨ! Kahikɨ
wahato kuɨ makure. 

Ei harɨkano rë kɨ hõra nohi waɨ kui, a nohi karɨkarɨmouwei! E puhi kuɨ
ha kuikunɨ, wãri rë të wakëxi wai praama mare makui, të pë ha rãmanɨ, të
pë mamo kãi ma rë përɨmapouwei. 

Ɨhɨ pei kahikɨ ha, kɨ rë harɨkaɨwei, kɨ wapɨɨ he harukei, hesiopɨ xo.
Kãrɨ, kãrɨ, kãrɨ! E kɨ nakɨ kupɨma wahato. 

Ɨhɨ e pë ha rarɨnɨ, pë noã rë wahato taɨwei, taɨwei, taɨwei: 

--- Hoaaaaaaaa! Pei këëë! Pë ta taeo! --- e pë kuɨ xoatarioma. 

Ai patama kamiyë pëma kɨ patamanɨ kaɨ wakë ha tararënɨ, wakë rë
pëtarionowei, wakë kupronomi. Kamiyë pëma kɨ no patama preaama. Riyë të
pë ha kamiyë pëma kɨ no patama xi wãriprou no preoma, hawë rãakãi të pë
kuaama. Yaro riyëriyë pënɨ të pë xi wãriprou no preoma. 

Ɨhɨ iha wakë he ha ha piyërehenɨ, Iwariwë a prukaremahe, kaɨ wakë ha,
wakë yua piyëkëapotayomahe. 

Të pë puhi rë owahataohe të pë xĩro õha kunoha, kaɨ wakë mori hanomi,
wakë mori kupronomi. Ɨhɨ a mao ha kunoha, hei riyëriyë pëma të pë waɨ
xoawë. 

Ɨhɨ Iwariwënɨ të pë ha hirakɨnɨ, wakë tapraɨ ha hirakɨnɨ, napë pënɨ wakë
tapraremahe. Të wakë pë kãi hĩihaɨhe, ɨ̃hɨnɨ të hirakema yaro. Pei kahikɨ
ha wake titipoma yaro.

Hapa anamahu wakë rë kui, wakë ha homoprarutunɨ, heinaha hawë maramahe
kure të ha, wakë anamahu pata makeama, ɨ̃hɨ wakë anamahu ha kɨ yaropɨ
rɨpɨpɨo hãtooma, kaxapɨ hãrɨkɨpɨo hãtooma, ɨ̃naha të tama. 

Wakë no xi ɨmapou he parohoma, wakë kãi mori wawënomi, kurenaha ɨ̃hɨ a
napë kõkaprou rë piyërahei, ɨ̃hɨ re e wakë wawëmapehe, a napë kõkaprou
piyërayomahe. 

Ĩxiĩxi a rë iaɨwei hamɨ, hena pë si pomama, pita pë rëmama, maxita pënɨ
të pë paterimama, pita pë ka ha titanɨ, hena pë ĩxi ɨnaha të tama, kama
yami a totihotii yaro, rɨpɨrɨpɨ të pë ha. Hei, iwa kë a, wa rë kuɨwei. 

Kaɨ rë wakë no xi ɨ̃hɨtapoma? Wakë noãpou he parohooma. Ɨhɨ wakënɨ wama
aka hemata tapraɨ. Hei iwa kë a, wa rë kuɨwei, pei aka no watëno pëwëmi,
ɨ̃hɨ kaɨ wakënɨ aka rë haikiarenowei, korokoro të wai hemata hëa. Wakë
hopoma yaro. 

A napë kõkaprou piyërayomahe, e të pë waroo piyëkema, kihamɨ e të pë noã
no waxuoma, të pë rë nakayouwei, të pë kõkamoma. Kama a mohotio tëhë, ai
kaxa a yuaɨ mɨ napë kuo tëhë, a napë kõkaprou heamahe. 

Ɨhɨ wakë rë wawëmare hamɨ, rɨpɨrɨpɨ të pë ha, të pë puhi yai toprarou
haɨtao puhiopë yaro, ĩsitoripɨ të nakaxi wai taa he ha yatirarɨhenɨ. Ɨhɨ
iha wakë rë wawërɨhe, napë, Yanomamɨ pënɨ wakë ha piyëarɨhenɨ,
piyërehenɨ, pëma kɨ iaɨ, rɨpɨ të pë ha. 

Hapa pëma kɨ mori ianomi. Ɨhɨ weti ha pata pënɨ rɨpɨ të pë wamahe? Ĩyë
makui, të pë kãi wamahe? Ɨnaha pëma të mori pruka tama. Horema
korimorewë: Tɨ̃ɨ̃ɨ̃! Tɨ̃ɨ̃ɨ̃! Wama pë kuɨ hirii? Ɨhɨ naxomi wama xi kɨ mori
wama, ɨ̃naha pë no preaaɨ kuoma. 

Ei a napë rë kõkapraruhahei, hei kurenaha e të pë kõkaproma. Rë
kõkapraruhe, a wasitaremahe. 

--- Pei! Xoape! --- pë xɨɨ hekamapɨ e kuoma --- Serororo! --- të pë wai
rë kuɨwei, ɨ̃hɨ hekamapɨ ree. 

Hawë xãraima ãto pë wai rë kure mau u pë ha, të pë wai keo ha:
Serorororo! Të pë wai rë kuɨwei. Temoa, hawë reha pë rë kure naha, të pë
rë kure, hekamapɨ Yanomamɨ e kuoma, hei kurenaha e kuoma. 

--- Pei! Xoape! A ta yanɨkɨtaru! Ai wa të pë ã hirii. Hei të pë rë
kõkamorɨhe, të pë ã wayou ta hiri! Henaha yai, henaha wa kɨ napë yai huɨ
kõopë --- hekamapɨ e kuma. 

--- Haaa! --- e kutarioma, pë a marë wahëi 

--- Hai! Hai! Exi të ha? --- e kunomi --- Ho! --- e kutarioma. 

Ɨnaha kama a kuɨ xoawë:

--- Hëëë…! --- a rë kuɨwei, pei a wã haɨ. 

Hei a noã taɨ ha, a yanɨkɨtarioma. 

--- Ho! Awei! Pei pëhë kɨ no tatou! --- hesiopɨ̃ mau u pë hawaro hamɨ të
pë ã rë karëhouwei. 

Pẽi, pẽi, pẽi! Mau u pe he tatoopë ha të pë rë kuɨwei, pẽipẽimɨ, ɨ̃hɨ
hesiopɨ e kuoma, Iwariwë hesiopɨ. Ɨhɨ hesiopɨ e makui, hawë moka e wai
kuwë. Waku kɨ wai oni, ɨ̃hɨ e kuɨ ha, e yai yɨkɨtarioma. 

Ai të pë hapa rë kuprore makui: Hɨ̃ɨɨ! Ɨha a puhi kua yaro 

--- Ma, hei kuikë ipa ya wakë rë tapore, ya wakë wawëmaɨ kuami --- a
kutou xoarayoma. 

A xomi huxutou nokarayoma. A ha miikunɨ, të pë rë mɨre, a kuɨ maopë.
Kama pë praopë ha: Hɨ̃ɨɨ! Të pë mohekɨ mɨ marë kutaowei. Pëkɨ yawëtëa
xoaparioma. A kãi karihipronomi. Ɨhɨ a rë miore, hei të pë noka rĩya rë
ĩkamouwei, hei e të pë wai yaiataroma. Kĩritari e pë wai yaiataroma. Ão,
ão, ão, ão, õoo! Të pë xomi kuma makui, a kãi kupronomi. 

Ɨhɨ ëyëha rë wakë rë titipore, hëyëhë të wakë anamahu pata ma
watawatamope, ɨ̃naha husi kua xoakema ĩkari, hawë kõmikõmi. Ɨnaha kahikɨ
kua xoa parɨoma. 

A rë kure tëhë, të pë pëɨxokɨa hërɨɨ, yaiatarou, huimiri e yaiatarou,
tãɨsarakari e pë yaiatarou,ẽhoamɨri e pë yakaaprarou, yõreketerari e pë
roaroaaprarou, roroaprarotiii, roroaprarotiii, e pë kãi
yahuyahuapraroma, a ka ĩkamaɨ xi totihitaoapehe, e pë yaiataroma.
Hesiopɨ e rë kui kihi e yaipë rëa kure: 

--- A kupro tihë! Miha a kuaaɨ kuparuhë! 

E kuɨ nokamoma: 

--- Kuaa tihë! Miparu hërɨ! --- e kuɨ nokamoma. A mɨ hururanɨ! Të pë ã
ma teteo tëhë:

--- Ɨnaha rë wama kɨ kuaaɨ kupe --- e kunomihe. 

Yakumɨ, të pë ha motarɨnɨ, të karihiproimi ha, të pë ha waximirɨnɨ, pei
të pë wai kuaaɨnɨ, të pë wai ha ninirɨnɨ, të pë no preaatii yaro. 

Ah! Kihi Tẽxori a makui, e kãi hokëprou haɨonomi. Kɨpɨ yai, hei kɨpɨ mɨ
amoa. Hei kɨ yawëtëpɨa kupiyei. Tohomamoriwë a yai, kihinɨ ɨ̃hɨ kihi pata
e rë kui, kihi Tẽxoriwë, ɨ̃hɨ a rë kui pata kee. Tohomamowë oxe kee, të
pë wai rë sinapii, ĩsitoripii, ɨ̃hɨnɨ kihi e kuaama, pata e rë kui e
kuaama: 

--- Pei! Apa, kahë! Wa wapaɨ xomao, kamiyë ya mamo yëo parɨo --- e
kuma. 

Pata e kuaaɨ makure:

--- Sĩo! Sĩo! Sĩo! Tõu, tõu, tõu! --- e kuaama. 

Kuaaɨ makure, pei mɨ tarɨ ha, hëyëha e katioma, Tẽxo e husi ha
horoiprarunɨ, e mamikɨ kãi ha horoiprarunɨ, e posi wai horoi
wauhuapraroma. E aka kãi nianiamoma makui, e mamokasi homotarioma, a mɨɨ
ha, miaa kõkema. 

--- Kahë iha mai! --- e puhi kutarioma. E kuaaɨ ha horohoprarunɨ, hei të
pë rë hɨtɨtɨpraruhe, të pë waximi yaro. 

--- Pei! --- pata e ã hama, oxe iha --- Pei! Õasi! Kahë! Haɨpraru! --- e
kuma. 

--- Kahë iha pei të yai xi kirihiprario --- kuɨ ha: 

--- Waooo! Pë mamo ta yëparu! Ɨhɨ rë! Wamare a ta mɨ! Kamiyë yai! --- e
kuma. 

E të pë mɨ puruparioma. Yõreketerariwë hei të no pretaa kure. Hawë e
rãakãi e kutaoma. Ɨhɨ kaɨ rë wakë no tapou kure. Wakë rĩya ha nokaanɨ.
Hei e no rë tare, e hokëtarioma. 

--- Pei, ɨhɨ rë, pë mamo ta yëparu! Ɨnaha ipa të kua. Sĩo! Sĩo! --- e
kuma, ɨ̃hɨ kiha ree kuɨ kure. 

--- Sĩo! Sĩo! Sĩo! Tëɨ! Tëɨ! Tëɨ! 

Hawë të krë e ã taa xoamakema. 

Ɨnaha a kuɨ tëhë, ɨ̃naha e të kuɨ totihiatarou hapa hërɨa ha, hei tënɨ
kaɨ wakë rë tapore, mamokasi homoprou nokarayoma. E mamo homotarioma.
Ɨhɨ të kupë hamɨ, e mamo xatiprakema, e hokëtou nokarayoma hokëprarioma:

--- Wão! --- hawë e yë hërɨpë --- Sĩo! Sĩo! 

E ku hërɨma pei mɨ tarɨ ha e wai kutou xoarayoma. 

--- Tëɨ, Tëɨ, Tëɨ! --- e kuaprarou ha. 

Hawë e yaua e wai kupario ha, e wai ha wayakarɨnɨ, e kuma. E të texinakɨ
mɨ wai ha yaprekerɨnɨ, e kɨpɨ mɨ wai hetutupɨkema. Hawë hekuramou e kɨ
wai tikutikupɨapraroma. 

Kuaaɨ ha Iwariwë a hokëprarioma. A puhi yopraa makupe, a ha hokëprarunɨ,
a ha tipëtarunɨ, ĩkawã wã no rarɨɨ xoaopë. Tohomamoriwënɨ kaɨ wakë no xi
ɨhɨtarema yai, heinaha e të mɨ wai tarɨaprarou kupe, hei, hei, mɨ wai
taɨaprarou, mɨ wai tarɨaprarou tëhë: 

--- Ho, ho, ho, ho! --- e kurayoma. 

Prohu! Kaɨ wakë anamahu pata yutuparioma --- Tou! --- Hei të no rë tare,
Yõreketerariwënɨ, e wakë nokarema. Ɨɨɨ, krihi, krihi --- pei husinɨ,
husi no rë watënowei, ɨ̃hɨ husi marë wakëɨi, ɨ̃hɨ kaɨ rë wakë unosi, e wakë
nokarema. --- Hɨ̃! --- wakë pata hute yaro heinaha e wakë kãi pata
kuapraroma. Wakë mori kãi pata kei mɨ yapaoma. Hei kihi a no rë tare. 

 

 

 

--- Xĩapo! --- të pë rë kuɨwei, kihi e paoma, kihinɨ: Wooo! Ɨhɨnɨ wakë
pata nokare hërɨma, pë ma rë prei, ɨ̃hɨ. Tu tu tu tu tu tu tu! Apia hi
pata hëwë ha, pei hi nanoka pata ha, wakë pata anamahu makeketayoma,
hesiopɨ e hokëprarioma. 

Hei Yõreketerariwënɨ e wakë mori mapramaɨ tëhë, naasi hirekoma. Hesiopɨ
e naasi hirekei makure, Xĩaporitawënɨ wakë kãi tirerayoma, naasi
hawërayoma. Ɨnaha e kuma, wakë kãi tirerayou ha, wakë misi rupramanomi
yaro, wakë tanomi yaro, naasi notikema yaro. 

--- Ɨnaha kë wama wakë ma tarenowei, ɨ̃naha kë wama wakë ma tarenowei,
wama wakë imi preaɨ ha kë, wama kɨ ã no rĩya preo hëo, wama wakë imi
rĩya preati hëo! --- hesiopɨ e kuma. 

Peheti a kuma, të pë nomaɨ, kamiyë pëma kɨ yaayopë, ɨ̃hɨ pëma kɨ pehi
wãha hiraprarema, kamiyë pëma kɨ ĩximayou hëopë, pëma kɨ pehi wãha
hiraprarema. 

--- Hawaro ha kë, xoati hawaro ha kë, ipa kë ya të ã kãi rii rë
topraowei kë, ya rii makui ha! --- e kuma --- Wakë imi preaɨ ta përahe,
wakë imi parimi preaɨ ta përahe, wama kɨ mamo rĩya protomotou kë! --- e
kuma. 

Peheti e të takemahe. Hesiopɨ peheti a kuma. Pëma kɨ yaayou, pëma kɨ
yãhi kɨ ĩxipë, të pë pehi wãha hirapraɨ katitirayoma. Kutaenɨ pëma kɨ
yaayou. Ɨnaha të kuprou mao ha kunoha, pëma kɨ yaayoimi. A kuma. A ha
kunɨ, kama, hẽaropɨ xo: Kruxu! Kopou! Mau u ha kɨ rii kupɨ kiriopë, kɨ
kupɨprou rë mai! Ai kɨ kãi pëpɨɨ, kɨ kãi kriipɨprou, kɨ he kãi
hayupɨprou, mamorinɨ kɨ kãi yupɨaimi, kɨpɨ matakɨ kãi nini no preaaimi,
hurapɨɨ taimi, kɨ kupɨopë. 

Hei kamiyë pëma kɨ nakɨ rë kurenaha ɨ̃hɨ pë nakɨ kuwë? Kuwëmi. Nakɨ kãi
niniaɨ taomi yaro. 

Ɨnaha të kuopë, të kuprarioma. Wã totihitawë kɨ marë kupɨa xoare. Kupɨa
xoaa. 

--- Pĩri! Pĩri! Pĩri! Pĩri! Pĩri! Pĩri! Pĩri! Pĩri! Pĩri! Pĩri! Pĩri!
--- e kuɨ morokaɨ hëparioma. 

Wakë yureihe ha! Ɨharë kɨpɨ, përɨpɨa xoaa. 

%\chapterspecial{O {surgimento} {do} {cupim}}{}{}


\letra{Aí}{há cupim}!--- nós dizemos. 
Por causa do rio no qualos Yanomami se afogaram, outros subiram nas
árvores por medo, ensinando"-nos. Ensinaram"-nos a subir. Subiram,
subiram. Conforme o rio subia, eles também subiram, sempre mais, até a
copa das árvores e, em seguida, se transformaram em cupinzeiros, mesmo
não existindo cupinzeiros nesse tempo. 

Os cupins grudaram aos troncos. Os Yanomami se tornaram cupim. A casa
dos cupins fica sentada em cima dos galhos. As casas sentadas são a
imagem daquela que existiu, oriunda dos Yanomami. 

Eles se sentam na forquilha das árvores. Os Yanomami ficaram sentados
nas forquilhas das árvores por medo. Apesar de quererem fugir, eles não
conseguiram. Depois da transformação dos cupins por causa do dilúvio,
aqueles que se afogaram boiaram à deriva na água e se tornaram
jacaré"-açu. Alguns se transformaram em jacaré, outros em peixe, outros
em capivara. Caíram na água, os {Yanomami}. Transformaram"-se
assim pela água. Não foi obra de ninguém! 

Quem os teria feito? Os cupinzeiros eram {Yanomami}. Desde a
transformação, feito isso o cupim está sentado e grudado às árvores nas
beiras de rio.

Uma casa gruda, outra está pendurada, outra está enfiada lá em cima. As
casas de cupim ficaram na posição na qual os {Yanomami} estavam.
É assim. 

São do tamanho de uma criança; uns quase ficaram na terra seca, uns
caíram na água por medo. Os que eram um pouco maiores caíram na água por
medo e se transformaram em cabas\footnote{ Vespas ou marimbondos.}  \emph{xaxa}. Vão se
chamar assim de imediato. Assim se transformaram. Após a transformação,
sua imagem se alastrou. Ocuparam todas as regiões onde moravam os
Yanomami. 

 

 

 

%\chapterspecial{{Yanomam}ɨ {pë} {rë} {ãrepopronowei}}{}{}
 

 

--- Kihi ãrepo kë ko --- pëma kɨ kuɨ. 

Ɨhɨ ko pë rë kui, Motu unɨ pë mixi ha tuo kuikunɨ, ɨ̃hɨ hei suwë ya wãha
yuaɨ rë taprarɨhe hamɨ, ɨ̃hɨ u nɨ të pë kiriri toreroma, pëma kɨ hiraɨhe
ha. Pëma kɨ tuopë, tuo hiraɨ ha. Të pë tuoma, të u pata tirei ha, të pë
kiriri ĩhetoma, ãrepo ko pë maoma makui, ɨ̃hɨ u pata ha rërërɨnɨ, ãrepo
ko pë kua hërarioma. 

Ãrepo ko pë yërëkoma, yërërarioma. Ɨhɨ Yanomamɨ pë rë kuonowei, ɨ̃hɨ ko
pë. Ko pë tikëprawë, hii hi pë poko hamɨ, ɨ̃hɨ të pë rë kuonowei, të pë
no uhutipɨ, tikëkëwë, ãrepo ko pë rë kui. 

Ko pë rë yakaroprai, hii hi pë hakarakɨ hamɨ, ɨ̃hɨ Yanomamɨ të pë kiriri
yakaropramoma. Të pë tokuu puhio makui, pë tokunomi. Ɨhɨ unɨ të pë xi
wãrihou ha kuikunɨ, ai pë mixi rë tuoprou he rë yatianowei, pë rĩya ha
pokëpronɨ, pë xomi rë niaaprarou huxomionowei, pë kãi ãopanaprarioma. Ai
a kãi iwaprarioma, ai a kãi yuriprarioma, ai a kãi kayuriprarioma, të pë
keparioma, Yanomamɨ të pë. Mau unɨ pë xi rë wãrihonowei, ɨnaha pë
kuprarioma. Taprano mai! 

Wetinɨ pë ha tapraɨ kuikunɨ të pë? Yanomamɨ kuoma makui, pë xi re
wãrihonowei, pë yai. Ɨnaha të pë kuaama. Të pë kuaaɨ ha kuikunɨ, hii hi
pë hamɨ, pata u kɨ hamɨ, ãrepo ko pë tikëkëwë, ko pë sutiprawë. 

Ai ko pë kãi yauprawë, ai ko pë yakaroa, kiha ai të ko wai hĩiatayoa,
Yanomamɨ të pë kuaanowei naha, ko pë kuwë. 

Ai të pë ihiru rë kurenaha, të pë mori haxɨrɨoma, të pë kãi keparioma,
kiriri. Hei kurenaha të pë kiriri keoproma, Xãxa na pë kuprarioma. Ɨnaha
të pë kuprarioma. Të pë kuprou ha kuikunɨ, të pë no uhutipɨ
praukurayoma. Yanomamɨ kutarenaha të pë urihipɨ hamɨ të pë kurarioma.
%\chapterspecial{{Os} {levados} {pelo} {rio}}{}{}
 

\letra{H}{avia um pajé} chamado Xiritowë, queveio a se chamar Keopëteri. Os
descendentes moram lá ainda. Os Xiritowëteri, depois, se chamaram
Keopëteri. 

Qual era o nome do rio onde eles afundaram, aquele que deu o nome de
Keopëteri? Esse rio se chama Xitipapɨwei. Eles bebiam dessa água, que
cobriu tudo. Caíram nas águas do Xitipapɨwei. Eles não morreram. Ainda
existem ali. 

Um dia, Xiritowë mandou seu genro buscar os visitantes
numxapono\footnote{  É o nome da casa coletiva circular onde vivem os Yanomami. Cada casa
dessas corresponde a uma comunidade ou assentamento. [\versal{OBS}. Não estou
grafando em itálico porque considero que pode ser usada tranquilamente
como palavra portuguesa].}  amigo para dançar durante a festa, e assim
nos ensinou a nos chamarmos e a nos convidarmos mutuamente. 

O genro de Xiritowë foi convidar os parentes e amigos deles, os
Anahupɨweiteri para dançar durante a festa, assim como fazemos até hoje.
Ele correu até o xapono dos Anahupɨweiteri. 

O pajé Xiritowë chamou seus conterrâneos para a festa sem imaginar que o
rio inundaria o xapono; ele pensava que todos iriam morar lá para
sempre. 

Como se chamava esse rio {antigo}? Esse rio se chamava
Xitipapɨwei.Aqueles que as águas cobriram, antes bebiam dessa água. 

 Xiritowë tornou a se chamar Keopëteri (os afundados). Depois de eles se
afundarem nas águas desse rio, tornaram a se chamar Keopëteri, porque
caíram nas águas desse rio que os levou. Assim, ainda moram lá. Esse rio
se chama Xitipa; os Keopëteri caíram nas águas do Xitipapɨwei. Eles não
morreram. Existem ainda. 

A filha de Keopëteri e seus parentes não se afogaram, apesar de estarem
no fundo do rio, o rio os levou. Eles vivem sempre lá. Tornaram"-se
eternos. Tornaram"-se esses monstros que nunca morrem; eles afundaram.

 


%\chapterspecial{{Keopëteri}}{}{}
 

 

Xiritowëteri a yai hekura përɨoma, kamiyë kurenaha, ɨ̃hɨ a wãha
Xiritowëteri. Kama Keopëteri pë wãha kukema kutaenɨ, kama Keopëteri ɨ̃ha
pë hiraa xoaa. Xiritowëteri Keopëteri pë wãha kuprarioma.

Motu u pata wãha rë kuonowei, weti naha u wãha kuoma? Motu unɨ pë rë
kepenowei, Keopëteri a wãha kuprarioma. Pë kepema yaro, pë yurema yaro.
Motu Xitipapɨwei u pata wãha kua. Ɨhɨ kama e u pata rë makepenowehei,
kama pënɨ u rë koanowehei, u makui, Xitipa kama e u wãha pata yua
xoaopemahe, Keopëteri a kepema. A nomanomi. Ɨha a kua xoaa. Ai pë wãha
xoaa. 

Siohapɨ e matoto rë ayonowei, ɨ̃hɨ weti naha kuwë pë xoama? Ɨhɨ hei kama
nohi e pë rë kuonowei, ɨ̃hɨ pë yahipɨ hetikëkëoma, pë rĩya ha praɨmahenɨ,
hei pëma kɨ nakayou hiraɨ ha, pëma kɨ xoayou hiraɨhe ha, pë nakayoma.
Xiritowëteri a hekuranɨ pë reahumou tëhë pë rë nakanowei. Pë kãi hakë
përɨotii tarei. A keo hëomai tao, të pë puhi ha kunɨ, matoto a ayoma,
kamiyë ipa reahu të kurenaha të kua yaro, a rërërayo hërɨma, kama pë rë
përɨonowei, kama maxi norimɨpɨ Anahupɨweiteri pë praɨmamahe.
Anahupɨweiteri kama norimɨ e pë yai kuomahe. Ɨhɨ pë xĩro hëtarioma. 

Ɨhɨ ei Keopëteri pë tëëpɨ rë kui, hei Motu unɨ pë mixi tuo taonomi, pëma
kohomoo waikio tëhë, kama unɨ pë rii yurema, pë rii kepema. Xoati ɨ̃ha
kama pë kua. Pë ma rë kure, parimi pë kupropë. Yai të pë nomaɨ rë mai pë
kupropë, pë kepema. 

 

%\chapterspecial{A {queda} {do} {céu}}{}{}
 

\letra{A}{ história do céu:} há a história da derrubada do céu. No início, ele
estava lá em cima.Este solo é o céu caído, portanto, os primeiros
habitantes foram esmagados. Esta terra os esmagou. Tornaram"-se, então,
os Amahiri, que moram em gruposcomo nós, mas moram lá em baixo. 

Aquele céu caiu. Os nossos antepassados saíram bem no meio; se tivessem
sido lerdos como os que foram amassados, não estaríamos aqui nesta
floresta. 

Depois daquele primeiro céu eterno, este segundo sobreviveu. Este novo
céu sobreviveu. Aquele primeiro, que estava em cima, envelheceu e caiu. 

Depois de sua queda, nós surgimos nesta terra mesmo, pois nossos
antepassados se multiplicaram nesta floresta. Apenas os habitantes de um
xapono sobreviveram para que possamos existir, mas quase que eles não
existiram. Em cima desta terra, nossos antepassados se reproduziram e
nasceram, e depois deles, os seus descendentes. 

O céu não exterminou todos nossos antepassados, ele esmagou alguns. Os
Amahiri se agruparam lá em baixo da terra. Os que foram esmagados se
chamaram Amahiri. 

Apesar de estarmos nas alturas, nós existimos. Assim, se essas montanhas
não houvessem existido e se não se erguessem, nossos antepassados não
haveriam saído, não haveriam existido. 

Eles ficaram espertos por causa dessas cavernas.

Quando começou a estrondar lá em cima, quando o evento se
aproximava, ele avisou seu povo. Os outros não sabiam:

--- Vamos, meus jovens, nós da região central vamos escapar por essas
montanhas, pois eu não sou tolo! Vocês não serão esmagados. Nós apenas
sobreviveremos nesta região, e se for preciso, passaremos ao outro lado
do céu. Limpem essa montanha! --- disse--- Abram o caminho! Limpem"-no! 

Dito isso, eles limparam ao redor da montanha para se abrigar. Quando o
céu ia cair, quando esse momento se aproximava, a montanha já estava
limpa. Na hora de cair, o céu se arrebentou, porque ele era antigo.
Enquanto se arrebentava, os que escaparam entraram na caverna. 

--- Vamos, enquanto o céu ainda está alto! Venham! Depois de amanhã,
depois de amanhã, o céu vai descer até o chão! 

Depois de ele dizer isso, o céu caiu. Esmagou os Esmagados. Moravam
nesse lugar. Quando o céu caiu, esmagou os que ficaram e os que estavam
na caverna não foram esmagados. O céu ficou por cima da caverna. Assim,
chamaram"-se os Derrubadores de Céu; era o nome deles. Queriam
derrubá"-lo, por isso se chamaram assim, com o mesmo nome. Apesar de o
céu quase os amassar, eles escaparam. 

Enquanto caía, o filho mais novo e o cunhado pularam e, assim, se
prepararam. O pai mandou"-os enfrentarem, mandou arrebentarem o céu. Fez
que o arrebentassem. Apesar de o céu parecer indestrutível, mesmo assim
os dois o arrebentaram. 

Aquele que arrebentou o céu, aquele que tinha esse nome, rebentou"-o
mesmo:chamava"-se Hutukarariwë. Atacou logo, sofrendo por causa do
sangue, sendo cortado pelos pedaços do céu, cortado perto dos
olhos. \emph{Kreti! Kreti! Kreti!} Fazia assim. 

Ele sofreu. Chamava"-se, assim, Hutukarariwë. Escaparam por onde ele
arrebentou o céu; por essa abertura, só o grupo dele escapou. Os que
sobraram, escaparam e saíram. Ninguém mais saiu de outro lugar. Nossos
ancestrais se reproduzirama partir daqueles que conseguiram escapar.
Continuaram de imediato a se reproduzir. Assim aconteceu. 

--- Enquanto o céu está vivo ainda--- disse o pai aos dois filhos---
vocês vão juntos! Vão embravecer! Não cortem em silêncio! 

A partir do momento em que nossos antepassados se reproduziram, surgiram
também os Waika. Dividiram"-se.

Antigamente, existiam também outros grupos. Outro nome importante dessa
época é o espírito Hemarewë. Todos esses nomes são nomes de espíritos.
Hemarewë vivia também nessa época como pajé, ele foi um dos primeiros
habitantes da região.

Mas não é o nome dos nossos antepassados. A história do nosso grupo
Parahiteri se encaixa no meio da história dos Yanomami. Alguns grupos
foram se extinguindo e de gerações posteriores foi que apareceram nossos
antepassados. 

 Nossos antepassados surgiram na região central. Nós ficamos nessa
região central, onde surgiu a primeira mulher, pois nossos antepassados
moravam lá. Ela nasceu nessa região central, ficava lá. Esses moradores
foram chamados de habitantes da região central. Foi assim que nós
surgimos.

 

 

%\chapterspecial{{Hutukarariwë}}{}{}
 

Hetu misi të ã: hetu misi rë tuyënowei, të ã kãi kua. Hapa a kuoma hei
heaka hamɨ, hei keno kë a. Kutaenɨ ei yëtu hamɨ të pë rë kuonowei,
hëyëmɨ pata të pë hɨ̃kɨwë. Hei a patanɨ pë xëprarema. Pë ha xëya hërɨnɨ,
Amahiri pë kuprarioma. Ɨhɨ misinɨ Amahiri pë xëye hërɨma. Hei pëma kɨ rë
kurenaha kuwë pë hiraa. Temɨtemɨ të pë kohomowë.

Kuwë yaro hei a kerayoma. Ɨhamɨ kamiyë pëma kɨ no patama rë hare, mɨ amo
hamɨ, ai të pë rë kui ɨ̃hɨ xëamorewë xoaonowei pë rë kui, pë puhi rë
mohotionowei rë kuonowei të pë kuo ha kunoha, pëma kɨ kuami, hei a urihi
ha.

Kihi a rë kui, kihi parimi a rë kui, a hëtarioma makui, kihi a
hëtaritayoma. Tukutuku a pata hëtariotayoma. Hei a rë kui a ha
horepɨonɨ, a kerayoma.

Ɨhamɨ pëma kɨ rë pëtore, të pë pararayoma. Hei a urihi ha, hei a rë kui
ha, mahu të pë rë hëtarionowei. Mahu a yahi pata hëtarioma. Ɨhɨ pëma kɨ
kupropë, ɨ̃naha pë mori kuonomi. Hei a rë kui hamɨ, pëma kɨ no patama
pararayoma, wawërayoma. Ha parahërɨnɨ, kamiyë pëma kɨ no hekama
kuprarioma.

Hei a rë kuinɨ, kamiyë pëma kɨ no patama xëaɨ haikionomi. Ai të pë rë
kui, të pë xëyë hërɨma. Kama pë rii Amahiri hiraopë, hëyëha. Amahiri pë
wãha kua, pë rë xëparenowei!

Tirewë makui, pëma kɨ kua. Hei kurenaha, makayo kɨ kuo mao ha kunoha, kɨ
pata ĩtao mao ha kunoha, pëma kɨ no patama kuami, mori hanomi. Ɨhɨ kɨ
kua yaro, të pë moyaweoma.

Moyawë yaro, kihamɨ të pata kimotayou tëhë, të si wëtɨkɨpraaɨ ha, pë noã
waxurema. Kihamɨ ai të pë mohoti no prepramoma.

--- Pei, huya pë, ipa wama kɨ rë kui, hei makayo kɨ rë kui, pëma kɨ
tokupë, mɨ amo yai hamɨ pëma kɨ rë kui, ya puhi mohotimi. Wama kɨ xëaɨ
maopë, ɨ̃hɨ pëma kɨ mahu rë hëtore, a urihi hamɨ, pëma kɨ no mɨhɨpropë.
Hei kɨ rë kui, kɨ ta wawëahe! --- e kumahe --- Kɨ no mayo ta takiohe! Yo
ta hõkɨhe!

E ha kuikuhenɨ, makayo kɨ të wawëaremahe, pë tokuu puhiopë yaro. E kei
kuketayomahei, hei a aheteprou tëhë kɨ kope kuo waikioma. A kei
aheteprou tëhë, të pata ha hëtɨrutunɨ, kama a rohote yaro, a hëtɨrayoma.
A ha hëtɨrɨnɨ, a hëtɨɨ tëhë, pë rë tokure, pë rukërayoma.

--- Pei! Haɨmohe! Ai të henaha, ai të henaha! A pata pitaɨ kurakiri,
ɨnaha a pata tireo kuo tëhë --- e kuma.

A ha kunɨ, a pata kerayoma. Xëapëteri pë rë kui, pë xëparema. Ɨhɨ pë
përɨoma. Xëapëteri pë xëparema. Kama pë wãha rë kuonowei. A ha kerɨnɨ,
pë rë hëpraruhe, pë xëyë hërɨma, makayo kɨ ha pë kua yaro, pë hɨ̃kɨanomi.
E pata hietiye kiriomahe, kuaaɨ tëhë, kama a wãha rë yehiponowehei, ɨ̃hɨ
hei a rë tuyënowehei, kama pë wãha kuoma. Ɨhɨ a wãha yehipomahe. A tuyëɨ
puhiohe yaro, pë wãha kuoma, kama a wãha rë kurenaha, pë wãha kure, pë
xëparioma.

A kei tëhë, a napë praɨpraɨmopɨmahe, ihirupɨ oxe pe heri ɨ̃naha a napë
kuaamahe. A napë ximɨkema, pë hɨɨnɨ. A no hëtɨpɨma mai ha, kurenaha a
pata hëtɨpɨmarema. Kama a yai rë hëtɨmarenowei, ɨ̃hɨ kama a wãha rë
yehiporenɨ, a yai hëtɨmarema. Ɨhɨ a wãha Hũtukarariwë a wãha kuoma. A napë kea xoakema, a ĩyë no preaaɨ makui, a rë hanɨɨwei, pei mamo kasi kɨ
hamɨ. Kreti! Kreti! Kreti! --- a tama.

A no preaama. Ɨhɨ a wãha Hutukarariwë kuoma. Ɨhɨ a rë hëtɨmare hamɨ, pë
harayoma, hei a ha. Ɨhɨ mahu a rë hëprore hamɨ, pë harayoma. Pë ha
harɨnɨ, ai të hare ha kunomai! Pë mahu ha harɨnɨ, kamiyë pëma kɨ no
patama raroma. Rarou xoao hërɨma. Ɨnaha të kuprarioma. 

--- Ɨnaha të pata yoprao kuo tëhë --- pë hɨɨnɨ pë noã tama. Kɨ noã
tapɨma, pë hɨɨnɨ. --- A napë ta patotoa xoaikuhe xë! Pei a napë ta
ɨramorɨhe xë!Mamikãi kãi tuyëatihehë xë!

Ɨhɨ tëhë, pëma kɨ no patama ha pararɨnɨ, Waika pë kãi kuprarioma. Pë no
patama xereroma. 

Kamiyë pëma kɨ mɨ amo hamɨ pëma kɨ no patama harayoma. Mɨ amo hamɨ, a
suwë rë kepraruhe a hamɨ, kamiyë pëma kɨ mɨ amoprarioma yaro, pëma kɨ no
patama mɨ amoprarioma. Mɨ amo hamɨ, a keprarioma, mɨ amoprou xoarayoma,
mɨ amo hamɨ pëma kɨ wãha hirapehe yaro. Ɨnaha pëma kɨ kuprarioma. 

Ɨhɨ pë wãha rë kui, kama pë xĩro rë përɨonowei, ai pë xoaa. Kamiyë pëma
kɨ no patama, napë pë no patama kãi, waiha opi pë pëtou, waiha të
warokei. Ɨhɨ te he tikë ha, hei yahë pë wãha rë waikare, ɨ̃hɨ të maxi ha,
mɨ amo ha, Hemarewë a hekura përɨoma. Pata të wãha, hapa Hemarewë.
Hekura kɨpɨ wãha. A hekura rë përɨonowei, pei kë a wãha Hemarewë kuoma,
hapa yai. Ɨhɨ a hekura përɨo mɨ hetuoma kutaenɨ pei a wãha Hemarewë pata
totihiwë. Urihi a kãi hapa përɨaɨ mɨ hetuoma kutaenɨ, a wãha yuamou,
pata. 

Hei kama të pë no patama rii rë pëtouwei të pë wãha. Hei kamiyë pëma kɨ
mai! Të ã yaia, kamiyë pëma kɨ no patama rii rë kuprarionowei. Kama pë
no hekama rii rë përɨhɨi, të pë wãha rii. 

Hei kamiyë mɨ amo hamɨ të ã kua. Ɨhɨ te he tikë hamɨ, hei ai xapono. Të
pë përɨhɨmou mao tëhë, kama të pë rë pëtarionowei. Të pë rë përɨonowei,
ai a xoaa. Ɨhɨ pë wãha. Ɨhɨ të rë kure hamɨ ai xapono, pë rë përɨonowei,
ɨ̃naha të pë kuoma. Hei ai pë waikou waikioma, hei pë rë hëaaimati, pë ha
maprarunɨ, hëyëha kamiyë pëma kɨ no patama rë kuprarionowei, ɨ̃hɨ tëhë mɨ
amo ha të yai kuprarioma. 

 

%\chapterspecial{O {sangue} {da} {lua}}{}{}
 

\letra{N}{o início,} os dois que flecharam a Lua também existiam, antes de nossos
antepassados Yanomami se misturarem. Eram espíritos. O irmão mais novo,
Uhutimari, morava com seus irmãos. Eram somente eles, junto ao seu irmão
mais velho, Escorpião. Eram três. O mais velho tinha o nome daquele
inseto que faz doer muito, o escorpião. Por isso, chamava"-se Escorpião.
O do meio tinha o nome da árvore paricá.\footnote{ A árvore paricá fornece as sementes e a casca com as quais os Yanomami
produzem um pó alucinógeno utilizado em diversos rituais.} 

Quem realmente flechou a Lua, o verdadeiro flecheiro da Lua, foi
Escorpião. 

Por que a flechou?Naquela época, a Lua era baixae ficava sentada na
terra. Sendo muito faminta de carne, ela devorava sempre as crianças.
Devorou o filho de Paricá. 

Era alta, assim como hoje? Não! Perto da casa de Escorpião, erguia"-se um
jatobá reto, onde a Lua se empoleirou desajeitadamente. Ela se sentou em
uma forquilha baixa. Paricá e seu grupo moravam junto com os dois
outros, Escorpião e Uhutimari. A Lua devorava as filhas quase formadas,
os filhos quase crescidos, as filhas quase moças. Ela comia as crianças
dessa faixa etária. 

Como ela as queimava, chamando assim o ódio de todos, os dois irmãos
mais novos quase a flecharam quando ela desceu para
atacar{.} Paricá se deslocou de \emph{wayumɨ} e ensinou aos
demais a se deslocarem de \emph{wayumɨ.} É por isso que hoje os Yanomami
vão de \emph{wayumɨ} onde há o bacabal para se alimentar. Os irmãos iam
de \emph{wayumɨ} e, assim, nos ensinaram. Foram por lá.

O xapono deles era como o nosso. Ele chorou como nós choramos, ficou
muito abalado. Depois da cremação do corpo do segundo filho, que a Lua
comeu em seguida, Paricá cobriu as cinzas no meio do xapono. Saíram
de \emph{wayumɨ}. 

Em determinado momento, um dos integrantes teve de voltar correndo,
tendo esquecido os dentes de cutia, outros ficaram sentados a certa
distância. Quando chegou à entrada do xapono, ele viu a Lua comendo as
cinzas no meio do xapono, que a gente sempre mantém limpo. A massa da
Lua se mexia.

--- \emph{Ũũũũũ}--- fazia um ronco assim. 

Ela comia até o carvão, devorava gulosamente as cinzas. 

\emph{Hɨ}! Ele ficou com medo e retornou rapidamente: 

--- Será que o monstro grande está fazendo isso? Ele está comendo? O
monstro está comendo as cinzas! \emph{Sãrai}!--- disse retornando de
medo. Ele foi buscar e avisar o pai da criança morta, Escorpião, que
fará a Lua sofrer as conseqüências. Ele buscou o pai. 

--- O monstro grande está comendo lá! Ele está comendo o que fez você de
luto. Ele está comendo gulosamente os restos, ele está comendo as cinzas
do seu filho--- disse. 

--- Ele está comendo as cinzas do meu filho!?--- disse o pai, desolado e
chorando. 

--- Vamos! Vamos! Vamos, meu irmão!--- disseram os dois irmãos mais
novos, apesar de eles não serem bons flecheiros. 

--- Já aprontamos as pontas de nossas flechas--- disseram. 

Escorpião observava os dois flechando em vão, pensando que eles não
conseguiriam, apesar de a Lua não estar muito alta, pois os dois eram
péssimos flecheiros. 

A Lua se empoleirava e olhava para si mesma, porque tinha comido o
filho. A Lua estava mole, digerindo mal. 

\emph{Taɨ xiri, taɨ xiri, taɨ xiri}! Os dois estavam flechando, mas suas
flechas, infelizmente, não acertavam o alvo. Fizeram a Lua subir,
espantaram"-na. Fizeram"-na subir, de tantas flechas que atiraram. 

Ela ficou altíssima, rodando e subindo, e os dois insistindo. \emph{Tai,
tai, tai, tai}! A Lua fez as flechas se tornarem espíritos. Por fim,
Escorpião, o pai da criança morta, conseguiu vingá"-la. 

Nossos antepassados não sabiam fazer guerra, foi ele quem nos ensinou. 

A Lua subia em direção à sua casa, sua rede estava lá, lá em cima. A sua
casa e a sua região estavam escondidas. 

Quando a Lua, que era diferente desta, passou pela porta, ele a flechou.
Quando entrou, ela estava cansada e deitou"-se lentamente na sua rede. 

Escorpião se moveu, levantando"-se e ao mesmo tempo apontando a flecha
para cima. 

Quando ela estava pronta para entrar, quando ela ia se sentar na sua
rede: \emph{prãoo, kroxooo}! Ele não falhou: apesar da altura, ele
acertou plenamente. Apesar de o vento sempre soprar muito nessa região,
a flecha não desviou, a flecha voou direto através do vento. Quando a
Lua se inclinava para deitar, a flecha se fincou entre as duas
escápulas.Escorpião a fez balançar. O sangue jorrou. \emph{Ho, ho, ho,
ho, ho, ho, ooooooo}! A Lua! \emph{Taka, taka, taka, taka, taka,
taka}! \emph{Ha}! Ele a fez estremecer. 

O sangue caído, as gotas de sangue caindo de lá para cá não se
estragaram. O sangue caía se transformando logo em Yanomami, mas em
Yanomami ferozes. O sangue se transformou em Yanomami, que imediatamente
flechavam. As gotinhas de sangue voaram sem se espalhar bem. O sangue
desceu flechando e não se esgotou. Os Yanomami, transformados a partir
do sangue da lua, mataram os habitantes do xapono do flecheiro da Lua. 

Ninguém sobreviveu, nem Escorpião, que se tornou espírito. Os dois que
conviviam com Paricá tampouco sobreviveram. O Sangue da Lua queria se
tornar Yanomami; queria se tornar ferozes Yanomami. Queria se tornar
matadores de Yanomami. Aconteceu, assim é. 

Foi então que surgiram nossos antepassados, a partir do Sangue da Lua. 

Assim aconteceu: o Sangue da Lua exterminou todos os que moravam em
baixo. Somente os espíritos sobreviveram, apesar de eles viverem como
nós. Os Yanomami Sangue da Lua não guerrearam com os espíritos.

Pouparam Kasimi e seu neto. 

 

%\chapterspecial{Pẽ{ripo} ĩ{yë}}{}{}
 

 

Pẽripo a rë nianowei, ɨ̃hɨ hapa kɨ kãi përɨpɨoma. Kamiyë Yanomamɨ pëma kɨ
yai rë kui no patama pë koyokoo mao tëhë, Pẽripo a rë niapɨnowei, ɨ̃hɨ kɨ
kãi kupɨoma, përɨpɨoma. Hekura kuoma. Kama oxe Uhutimari e pë kãi
hiraoma. Kama pë xĩro hiraoma. Suhiriwë pata, oxe ɨ̃naha e pë kuoma. Të
pë no nini he rë parohoi, ɨ̃hɨ patare e kuoma, Suhiriwë. Ɨhɨ rë a wãha
Suhirina kuoma. Ɨhɨ rë të pë wãha rë yehiponowei, ɨ̃hɨ rë a wãha Yakuana
kuoma. Suhirina, Uhutimariwë oxe, ai a, ai a, ai a, ɨnaha e pë kuoma,
kama pata. Kɨpɨ rë përɨpɨonowei, ɨ̃naha të kuoma. 

Pẽripo a rë niaɨwei hei a yai, Pẽripo niaprarewë, Suhirina. 

Exi të ha a niama? Kama Pẽriporiwë a yahatotoo tëhë, kama e pë waɨ yaro,
Yakuana a wãha rë kuonowei ihirupɨ wama, Pẽriponɨ. A naiki yaro.
Yahatoto ha kihi naha pita ha e të pata roo parɨoma yaro. 

Kihi a rë kui ɨnaha a kuoma? Hɨ̃ɨɨ! Kuonomi. Kama yahipɨ ahete ha,
pukature hamɨ të rë kurenaha, motua e ãhi wõroropë ha, e të pata
wahehetaoma. Yahatoto ha e të pata hãkioma. Ɨhɨ kɨpɨ kãi përɨpɨoma. Ɨhɨ
Yakuananɨ pë kãi rii përɨoma. Ɨhɨ pë tëë, suhe mo kɨ, ihirupɨ e mori
pataɨ, pë tëë mori mokou, ɨ̃naha e pë kuprou tëhë, e pë waɨ nokamoma. 

E pë iximama kutaenɨ, kama Pẽriporiwë a napë ha rurorɨnɨ, oxe kɨpɨnɨ a
mori niapraɨ makure, Yakuana a wayumɨ hokërayoma, të pë wayumɨ huɨ hiraɨ
ha, të pë wayumɨ huɨ, kihamɨ të pë iaɨ, hoko ma pë kuo pë hamɨ, të pë
wayumɨ hu hërɨɨ, kurenaha të hiraɨ ha, pë kuaama. Ɨhamɨ pë ukukema. 

Ɨhɨ hei ipa xapono kurenaha, pë xaponopɨ kuoma. Kamiyë pëma kɨ ɨ̃kɨɨ rë
kurenaha, a ha ɨ̃kɨrɨnɨ, puhi kãi no preaama. Wakë ha yëarɨnɨ, ihirupɨ a
rë nokaaɨ kõowei ha, ĩxino heha yohopa hërɨnɨ, a hokërayoma. 

Kihi naha ai tëka hikari rë prare naha kama e tëka praa, kihi naha të pë
kuke herayoma, roa kupe herayoma. Ai të rë rërëɨwei hei kurenaha, tomɨ
nakɨ nohi rë mohotuaɨwehei, të rërëimama, rërëimama, rërëimama. Pei yo
ka ha a mori kutou tëhë, hei kure naha, të rë kure, të rërëimama. Xapono
mɨ amo a yai ha, ĩxino a wai ha, ɨ̃hɨ, einɨ Pẽriporiwë a iaɨ tararema,
ĩxino a ha. Xapono ha të pë ma rë wawëtoi. Kihi yo ka rë kurenaha, yoka
kua yaro. Pẽriporiwë a iama, ĩxinoma a ha, xapono a mɨ amo ha të pata
hamorimoranɨ: 

--- Ũũũũũ… --- të pata kuma. 

Ẽxehẽroma kɨ makui, kɨ wama. Të pata wëhërɨwëhërɨmoma. Kuwë ha:

Hɨ! A kirirarioma, a kiriri he tatoprarioma. 

--- Yai rë të pata iaɨ ta yaipiyei, ei rë wa të no pata kiriaɨ nosi ë!
Yai të pata, ĩxinoma a ha, yai të pata rë iapiyei! Sãrai! --- e kua mɨ
yaparayoma. Ɨhɨ kama pë hɨɨ a kõaɨ ha, a yɨmɨkamapë. Suhiriwë a
yɨmɨkamapë. Ɨhɨ Suhirina iha, ɨ̃hɨ ihirupɨ, wama yaro kutaenɨ, a noa
prearema, ihirupɨ. A kõrema.

--- Pei! Yai të pata ia harayou. Hei wa wã no rë preohe, të pata iaɨ. Të
pata wëhërɨwëhërɨmou. Ɨhɨ ĩxinohë a waɨ --- e kuma. 

--- Ipa ĩxino a wa harayou --- e kuma, mɨa kãi. 

--- Pei, pei, pei, oxei, haɨmo! Pei, pei oxe, yai kë ũũũ! --- oxe kɨpɨ
kupɨma, e kɨpɨ nɨhɨtepɨmi makui. 

--- Pei! Ipa ya kɨ hĩihaɨ waikire! --- e pë xomi kuma. 

A no tapɨma mai! Oxe e kɨpɨ ma kupɨnowei, e kãi mori no mɨhɨpɨanomi,
kɨpɨ ninipɨo he parohooma. 

Hei kurenaha Pẽriporiwë a wahehetaoma, a mɨprou ha, ihirupɨ e warema
yaro. A ëpëhëoma kutaenɨ, a xi wãrima. E kɨpɨ nini prepɨamoma. Hei naha
të pata kuu makui:

--- Tai xiri, tai xiri, tai xiri! --- heinaha të pë xerekapɨ pata
hãrokaɨ kuaaɨ tikooma. Ɨhɨ të rë kui a tiremarema a ha yaxuprapërɨnɨ. 

A pata tirepɨmarema yakumɨ, xereka pënɨ hei a pata rë yaxupɨre, të pata
rë hamorimoimati, kuaa hërɨɨ tëhë. Tai, tai, tai, tai! Xereka rë pë ma
kuonowei, hekura rë e pë pehi pata hãrokoa xoarayoma. Hei a rë kuinɨ, pë
hɨɨ a yainɨ, a no yurema. 

Kamiyë Yanomamɨ pëma kɨ no patama niayonomi makui, ɨ̃hɨ të pë hirama. 

Pẽripo a tirewë makui, hei kama pëkɨ kihamɨ pë kɨ kuotayoma, ei yoka ma
rë kurati, yoka kua, kama yahipɨ, urihipɨ pesia yaro. 

Hei yoka kurenaha, inaha e yoka kua yaro. Ɨhɨ yoka kua xoaa. Yoka tapoma
makui, hei a rë kui a kuo mao tëhë, a niama. A hatayou tëhë, a
waximirayoma, kama pëkɨ hamɨ, a përɨo ka kuaa hërɨma. 

Hei e rë hokëtore, e rë hokëtou nokare, ɨ̃hɨ kama pëkɨ ha a përɨaɨ tëhë a
nokaaɨ puhiopë yaro, a pehi mɨ nonomarema. 

Kama e yoka ha të pata ha horeikutunɨ kama pëkɨ hamɨ, të pata tipëatayou
tëhë: Prãoo, kroxooo! Yai hamɨ e kãi morokõnomi, tirewë makui, xereka e
kãi hawënomi. Watori hiakawë a kua makurati, xerekapɨ kãi yarënomi. Ɨhɨ
watori a pata huxomi hamɨ xereka e morokoa katitirayo hërɨma. Pei pëkɨ
ha të pata wëkëatayou tëhë, pei a përɨatayou tëhë, mapuupëka pata
nokaretayoma. A pehi kãi pata wahehiamatayoma. Ĩyë kɨ pata rë
hirekerati. Ho, ho, ho, ho, ho! Ooooooo! Pẽripo, Taka, taka, taka, taka,
taka, taka! Ha! A pata porepɨ tamama. 

 

 

Ɨhɨ rë ĩyë kɨ rë kerati, mokure kɨ rë kui, ĩyë kɨ wãrimou kateheo maopë,
ĩyë kɨ kei kurati. Mokure të ĩyë pata kerayoma, hei a pita hamɨ. Ɨha ĩyë
ketayou yaro, ɨha ĩyë pata niayou rë xoarati, niayou kuketayoma. A ĩyë
rë keimati hamɨ, Yanomamɨ kurenaha, kama ĩyë yanomamɨprou xoarayoma.
Pẽripo ĩyë rë kui ĩyë nakaxi, remaxi yëɨ kateheonomi. Ɨhɨ ĩyë kãi niayou
itorayo tayoma. Ĩyë kãi mapronomi. Ɨhɨ ĩyë kɨ kepë ha, hei a rë niare,
kama yahipɨ haikiarema. 

Ai a hëpranomi. Kama Suhiriwë a kãi temɨ kutonomi, pei a no uhutipɨ
hekura hurayoma. Pei a no uhutipɨ xĩro rë hõriprariowei, a hekura
kuprarioma. Kama kɨpɨ rë kui Yakuananɨ kɨ kãi rë përɨpɨawei, e kãi
hëtonomi. A temɨ kãi përɨo hëonomi. Ĩyëpë yanomamɨprou puhio yaro.
Pẽripo ĩyë Yanomamɨ xëtimi pë kuprou puhio yaro. Të kuprarioma. Ɨnaha të
yai kua.

Kamiyë yama kɨ ĩyëpë ha ta kuio! Të pë puhi kuu maoma. Kamiyë pëma kɨ no
patama, ɨ̃hamɨ pë rë pëtouwei të kuprarioma, Pẽripo ĩyë hamɨ. 

Ɨnaha të kua: Pẽripo ĩyënɨ pepi hamɨ pë rë përɨonowei, pë haikiaremahe.
Hekura pë xĩro hëprarioma. Hei kurenaha pë xĩro përɨoma, pë kua makure,
hekura pë niaonomi. 

Ɨhɨ pë kãi hëpraremahe. Suhiriwë yesipɨ, hekamapɨ, ɨ̃hɨ e kɨ wai
hëpɨpraremahe, Yakuana nɨɨpɨ, hekamapɨ.

 

%\chapterspecial{{Kasimi} e o {seu} {neto}}{}{}
 

\letra{K}{asimi e seu neto} alcançaram os espíritos para não viverem como viviam
os {Yanomami}. Os espíritos viviam escondidos no mato. Todos foram
exterminados pelo Sangue da Lua.

Ou seja, a avó e o neto alcançaram os espíritos. Ela está lá ainda. É assim, não morreu. Aonde vão os espíritos, ela também vai, ela os
alcançou. Ela alcançou o lugar dos espíritos para se tornar eterna. 

Kasimi alcançou o xapono dos primeiros espíritos, os Ihiruwëteri, nome
de espíritos. O primeiro xaponodos espíritos era o dos
Ihiruwëteri (crianças), Kasimi chegou a eles, ela chegou aos
Oxemawëteri, que é outro nome deles. 

Kasimi alcançou a moradia dos Parawari Yokënamari, todos solteiros e que
estavam dançando. Ela chegou durante uma festa, e, assim, nos ensinou a
fazer festa. Kasimi era o nome da mãe daquele que flechou a Lua.

Os antepassados, que moravam espalhados, não resistiram aos Yanomami
oriundos do Sangue da Lua, de onde nós viemos. Fizeram isso para nós
brigarmos, para nós guerrearmos. 

Quando começaram esses eventos, os antepassados logoficaram espertos,
que antes nãoeram. Quando brigavam, era como uma dança, e
simplesmentenão paravam de rir. Quando havia guerra, não sabiam reagir
e {só} faziam {xamanismo}. Antes, eles não se vingavam.
Durante a guerra, eles festejavam; era isso que eles faziam! 

No início, existiam os que ensinaram os Yanomami a morar. Os espíritos
existiam e eram parecidos com os Yanomami. Os espíritos não foram obra
de ninguém. Eram assim, como os Yanomami. O xapono deles era tão limpo
como o meu, não era fechado. Moravam juntos. 

Dizem que moravam assim sem ninguém ensiná"-los. Eles viviam em um xapono
igual ao meu. Eles andavam sempre no limpo. Eles faziam amizade,
conversavam e se visitavam. Se não agissem assim, teriam nos exterminado
há muito tempo, pois eles nos comem. Se eles morassem ainda no limpo,
todos vocês, rapazes, diriam: \emph{ea, ea, ea}! Todos vocês seriam
pajés. 

Hoje, os espíritos não são mais visíveis, pois não moram mais no limpo.
Eles dançavam, faziam festas no limpo, como os Yanomami. Eles dançavam
como dançam os Yanomami, no limpo. 

Faziam os rituais de \emph{himou} e de \emph{wayamou}, cantavam como os
pajés cantam. Eles cantavam assim. Eles também brincavam de roubar
esposas como fazemos durante as festas. 

Moravam na planície, em terra plana, não moravam naquele tipo de
montanha. Depois, eles foram morar lá nas montanhas, foram logo assim. 

Como eram todos Yanomami morando no limpo, a imagem da minha avó, apesar
de ser espírito, ainda alcança os pajés, porque ela era Yanomami. A mãe
daquele que flechou o monstro chegou aonde moravam os matadores de
monstros, que moravam no limpo. Ela deve ter chegado enquanto eles ainda
eram visíveis. 

Aquela que fechou a casa dos espíritos se chamava Kasimi, mesmo. O segundo nome dela era Maxikomi. 

Ela carregava um grande cesto. Não existia porta como essa. Ela
carregava esse cesto andando no caminho dos espíritos. 

Enquanto eles olhavam para Kasimi, o cesto apareceu. Não havia carregado
o cesto até esse momento. Depois de o cesto ficar visível, apesar de
haver uma grande entrada, o cesto não passava pela entrada. 

Ela tentava entrar com o cesto, que bloqueava a entrada, por isso os
espíritos deram uma gargalhada. Enquanto estavam rindo, ela se mexia
para conseguir entrar. A casa dos espíritos estava se fechando devagar,
fechando devagar e o xapono acabou fechando totalmente. 

O xapono onde moravam os espíritos e cuja entrada ficou fechada tinha
nome: Yoararopɨwei. Esse xapono se chamava Yoararopɨwei. Era muito
bonito. Gostavam muito dele. Apesar de ser um xapono, ele era muito
brilhante, como um espelho, possuía uma luz própria. Colocaram o nome de
Yoararopɨwei.

--- É o xapono de Yoararopɨwei--- diziam. 

Kasimi chegou lá. 

--- Meus queridos! Meus queridos! Estou chegando com meu netinho deum
grande sofrimento. Esperem"-me! Esperem"-me aí! Abram a entrada!--- dizia
vindo.

--- \emph{Õ}! É a voz de quem? Quem é, será?

--- Meus queridos, abram a entrada! Estou chegando com meu netinho!---
disse, vindo--- Eu estou chegando e sofrendo de fome! Agradem a meu
netinho! O único que restou, agradem"-no! Ele é meu único!

--- Quem é você? De quem é essa voz?--- disseram. 

Eles queriam que ela pronunciasse seu nome. 

--- Quem pode ser? Quem é você mesmo? Essa voz de mulher, de quem pode
ser?\emph{ Ãaaaaaõooooo}!--- disseram. 

--- Sou Kasimi! Sou Kasimi! --- disse.--- Queridinho, não pergunte quem
sou! Sou Kasimi!

Escutava"-se o som dos pulos. Infelizmente a entrada fechou. Ela fez
assim, como quando alguns \emph{napë}\footnote{  É o termo yanomami para designar seres humanos não Yanomami,
particularmente os não índios.}  ficam presos na
cadeia. Foi assim. 

A história dos espíritos foi obra de alguém? Não pensem assim! Não foi
obra de ninguém! Ela aconteceu através de Kasimi. Essa é a verdade! 

Eles também comiam, comiam banana"-pacovã, faziam festas à época da
pupunha, faziam também beijus, sabiam caçar, comiam anta, quando faziam
festas; era assim que viviam os espíritos, no início. Tiravam lenha,
assim faziam. Plantavam bananeiras, enquanto moravam no limpo,
ensinando"-nos, assim, a plantar. Nós continuaremos a plantar os
alimentos como eles os plantavam. 

Atualmente, quase que nós não comeríamos pupunhas. Foram os Japiins que
espalharam as sementes de pupunheira. Não foram os antepassados
dos \emph{napë} que criaram essas pupunheiras. Eles não inventaram as
sementes. 

Omawë plantou pupunheiras, depois de inventar as sementes? Não, não foi
ele quem fez isso!\footnote{  Aqui o narrador se refere a personagens desse ciclo.} 

Quando nós desconhecíamos a pupunha, os japiins se agrupavam no chão, as
pupunheiras se erguiam perto daquele xapono cuja entrada ficou fechada. 

Somente eles faziam festas à época da pupunha. As pupunheiras não
existiam lá onde moraram Yoawë, Omawë, Ruwëri e Pore. Foram os Japiins
que criaram as pupunheiras. Os Japiins moravam com seus irmãos mais
novos, os Jaloacas. Foram eles mesmos quem deram essas palmeiras. 

Não foram os Yanomami que conseguiram as sementes para podermos comê"-las
hoje. Não comemos pupunhas hoje devido a um antepassado
dos \emph{napë} que as houvesse feito aparecer. Elas nos foram
transmitidas. Foi Japiim quem as conseguiu. A pupunha se espalhou. Foi
assim. Assim foi,na verdade. 

Apesar de serem espíritos, foram eles quem ensinaram os Yanomami a fazer
festas. Os Yanomami seguem o ensino da festa daqueles dois que faziam
festas, mesmo sozinhos. Eles moravam na região central. A partir daí,
nós faremos festas. 

Ele nos ensinará a fazer o ritual de \emph{yaɨmou}. Qual era o nome
desses dois? Será que alguém ensinou a vocês o nome desses dois? Se
um \emph{napë} perguntasse o nome dos dois que moravam juntos, alguém
pronunciaria o nome dos dois? Como eles se chamavam? Esses dois
ensinaram a festa e o ritual do \emph{yaɨmou}. 

 

 

%\chapterspecial{{Kasimi}}{}{}
 

 

Hekura pë iha, hei kurenaha pë kuopë ha, kɨ waropɨopë, kɨpɨ
hëpɨprarioma. Waropɨkema, awei, kɨpɨ si rë poaɨwei. Ɨhɨ hekura pë iha, a
rë waroore, Yanomamɨ pë përɨaɨ rë kuonowei naha, a kuprou maopë. 

Hekamapɨ kãi warokema. Ɨhɨ Pẽripo ĩyënɨ pë no rë watëɨwehei, pë yesinɨ
hekamapɨ a kãi warokema, hekura pë iha.Ɨha a kua xoaa. Ɨnaha të kua. A nomanomi. Hekura pë huɨ tëhë, a hupë, kama a warokema. Parimi a rë
kuowei, a warokema. 

Hekura, ɨ̃hɨ hapa pë rë kui, Ihiruwëteri pë xaponopɨ kuopë ha, Kasimi a
warokema, hekura pë wãha. Ihiru të pë yosika rë rararenowehei pë yaro,
pë wãha Ihiruwëteri kua, hekura. Ɨhɨ pë iha Kasimi a warokema. Hekura
Ihiruwëteri, Oxemawëteri pë iha a warokema, Kasimi a warokema. 

Parawari Yokënamari pë yai hiraopë ha, xĩro të pë pata praɨpraɨpraropë
ha, Kasimi a yai warokema. Kamiyë pëma kɨ reahumou hiraɨhe tëhë, pë
reahumou tëhë, e warokemahe. Kasimi a wãha kuoma. Ɨhɨ Pẽripo a rë
nianowei nɨɨpɨ wãha. Pë nɨɨ e wãha Kasimi kuoma. Ɨhɨ a hëprarioma, yami
xĩro, hekamapɨ xo. 

Ɨhɨ ĩyë ha yanomamɨprarunɨ, hei kurenaha, ĩyë ha pëtarunɨ, ai të pë
hëpranomihe, hei pata pë rë përɨhonowei. Pëma kɨ xëyopë, pëma kɨ
niayopë, të tapraremahe. 

Hapa ai të rë mohotimouwei, të kuonomi. Hapa të pë puhi mohotioma. Të pë
xëyou tëhë, praɨɨ kurenaha, të pë ka xĩro ĩkaprou pëotima. Të pë ma
niaɨhe tëhë, të pë imɨkɨ kãi rërëkëapraroma. Hapa të pë no yuayonomi. Të
pë ma niaɨhe tëhë, të pë kãi praɨma. Ɨnaha pë kuaaɨ pëoma. 

Ɨhɨ hei Pẽripo ĩyënɨ të pë rë niare a patamorayoma, Yanomamɨ hei të pë
kupropë. 

 

Kamiyë Yanomamɨ pëma kɨ përɨaɨ rë hiranowehei, hapa të pë përɨkema. Hapa
hekura pë rë kui Yanomamɨ kurenaha pë kuoma. 

Taprano hekura pë kuami. Kama xoati. Yanomamɨ kurenaha. Hei xapono ipa
kurenaha, wawëtowë yahi kurenaha pë kuoma kãi, kahuhuwë mai, pë
hiraoma. 

Kama pë rë përɨohe hõra, ai tënɨ pë rë hiranowei, hekura pë kuami. Hei
kurenaha ipa xapono hamɨ Yanomamɨ kurenaha pë kuoma. Wawëtowë pë kãi
hutima. Pë nohimoma, pë ã kãi wayoma, pë kãi hama huma, kuaaɨ tëhë, pë
kuoma makui, ɨ̃naha pë kuo ha kunoha, kamiyë yëtu hëmɨ, pëma kɨ
haikiareihe. Pëma kɨ rë waɨwehei. Ɨnaha pë kuo xoao ha kunoha, hei huya
wama kɨ rë kui. Ea, ea, ea! Hɨtɨtɨwë wama kɨ kuɨ. 

Pë wawëtowë kua yaro. Hei kuikë pë wawëtoami. Yanomamɨ kurenaha pë kãi
praɨma, praɨama, wawëtowë, pë kãi reahumoma. Yanomamɨ praɨpraɨmou
wawëtowë kurenaha pë kuaama. 

Pë kãi himoma, pë wãyamoma, ɨ̃hɨ pë rë hekuramore, amoa kurenaha pë kuma.
Amoa a rë taɨwehei ɨ̃naha pë kuma. Pë hesiopɨ kãi hãkɨoma, pë kuaama. 

Wawëtowë pë hirao tëhë, yarɨta ha. Kama pë yarɨtaoma yaro, kihi ma pë rë
kui, ɨnaha të pë kuonomi. Ɨhɨ kihi kama pë yahipɨ hamɨ pë përɨhɨrarioma.
Ɨnaha pë kurarioma. 

Hekura taprano pë mai! Wetinɨ pë taprarema? Kama xoati pë Yanomamɨ rë
përɨkenowei pë yaro, kuoma makui, pë wawëtowë kua ha, ɨ̃hɨ yayë a wãha no
uhutipɨ rë hirore, hekura a makui, a waroo xoa. Yanomamɨ a kuoma yaro.
Ɨhɨ yai të rë nianowei nɨɨpɨ ɨ̃hɨ pë xĩro kuopë ha a ha waroa taronɨ,
wawëtowë pë kuopë ha, e kõo kateheo ha maohenɨ. 

Pë pëka rë kahuprarenowei, Yanomamɨ a wãha, ɨ̃hɨ a wãha Kasimi kuoma.
Maxikomi ai a wãha. Pei pëka yai rë kahuprarenowei, Kasimi a wãha yai
kua. Maxikomi oraora a wãha. Korokoro a wãha yai Kasimi kuoma. 

Ɨhɨnɨ yotema a pata ha yehitarɨnɨ, hei yo ka rë kurenaha, yo ka kuami
makure, pë peipɨ yo pata rë haawei hamɨ të pata yehitarema. E pë mɨ
puruama, pë mɨ puruaɨ tëhë, e të pata pëtarioma. A yehiponoma mai! Të
pata ha pëtamarɨnɨ, të yosi ka pata hore kuwë totihiwë makui, e pata
hõkikeyoruma. 

E të pata ha hõkiikionɨ, të kãi pata rë hare, ɨ̃naha a kuaaɨ tëhë, e të
pata hõkikei ha, hekura e pë ka ĩkapraroma, no ka ĩkaprarɨ he, ĩkaprarɨ
he kurenaha e të pata yãɨka hërɨɨ tëhë, hekura pë yahipɨ ka
komaaimatayoma, komaa hërɨma, komaa hërɨma, komaa hërɨma, pë pë ka
kahuprarema. Kihi pei ma pë ma rë kõmii. Pë ka kahuprarema. Tei, tei
huxomi hamɨ pë tapramarema. Hekura pë huxomipramaɨ xoarayoma. Të yai
kuprawë. 

Kama hekura pë pëka rë kahumanowei, pë xaponopɨ rë kuonowei, pei pë
xaponopɨ wãha. Ɨha pë rë përɨonowei, ai të pënɨ pë wawëtowë rë
përɨonowei, pë rë reahumonowei, e wãha xapono yupraɨ taohe? Ɨhɨ xapono
Yoararopɨwei a xapono wãha, Yoararopɨwei. A riëyëhëo he parohoma. A nohi
toaɨ totihiomahe. Xapono a makui të xĩi pata hãtohãtopraramamahe. Mɨre
pë xĩi rë kurenaha e xĩi kuaaɨhe yaro, a wãha Yoararopɨwei tapomahe. 

 

--- Yoararopɨwei kë a --- pë kuma. Ɨha e warokemahe. Kasimi a warokema. 

--- Oxeiwë pëë! Oxeiwë pëë! Ipa xëtëwë! Ya no kãi wai preaaimi! Wamare
no ta tapa! Miha wamare no ta tapa! Yo ka ta karokɨhe! --- a kuimama. 

--- Õ! Weti kë a wã? Weti wa wã ta tawëëë?

--- Oxei, oxei, yo ka ta karopahe! Ya hekamapɨ wai pararuaɨ kë a kure!
--- e kuimamahe --- Ipa ya të ohiri no kãi wai preaaimi! Ipa të nohi wai
ta toahe! Ipa të wai mahu rë hëprouwei, ipa të nohi wai ta toahe! Ɨhɨ
ipa të wai mahu yaro. 

--- Weti kë wa? Weti kë wã ta tawë? --- e pë kuma. A wãha yupramapehe
--- Weti naha kuwë pei wa të ã yai ta tawë? Weti kë wa wãaaa ëëë? Weti
naha kuwë suwë wa të ã yai ta tawë, ãaaaaaõoooooo? --- e të pë pata
kuma. 

--- Kasimi kë ya wã, Kasimi kë ya wã --- e kurayomahe --- Oxei weti ma.
Kasimi kë ya wã! 

E të pë mamikɨ pata haruharumopë. Pë pëka kahupraɨ tikorayoma. Napë pë
pëka kahumaɨ hiraɨ ha, pë pëka ma rë kahuowei, ai pë pëka mare kahuowei
kurenaha të taprarema. Ɨnaha të kuprarioma. 

Hekura taprano pë wãha! Pë puhi kuɨ mai! Kama taprano pë mai! Kama
Kasiminɨ ɨnaha pë yai taprarema. Ei të ã yai. 

Pë kãi iama, kurata pë kãi wamahe, raxa a kãi reahuamahe, naxi hĩ pë kãi
ramamahe, pë naiki kãi taoma, xama pë kãi wamahe, pë reahuamamahe,
hekura pë kuaaɨ parɨoma. Kãɨ ãxo pë kãi tamahe, pë tarɨkɨ taoma, ɨnaha
pë kuaama, kurata si pë kãi keamahe, hei kurenaha pë wawëtowë kuo tëhë,
pëma kɨ̃ hiraɨhe ha. Ɨhɨ pënɨ pëma kɨ ni rë keore, pëma të taɨ hëopë. A nii keamahe. 

Kihi kamiyë pëma kɨnɨ kihi raxa pëma kɨ mori waɨ hëonomi. Ɨhɨ pënɨ ɨ̃ha
ãyakorari pë iha mo rë piyëarahei hamɨ kihi të si kɨ, piyërenowei hamɨ.
Napë pë no patapɨnɨ raxa mo tapraɨ taonomihe. Mo kãi tapranomihe. 

Omawënɨ mo ha takɨ hërɨnɨ, mo keke hërɨma? Keanomi. 

Pëma kɨ tao mao tëhë, ãyakorari pita ha pë përɨoma, ɨ̃hɨ pë pëka rë
kahuprarenowei hamɨ, kihi të si pë. 

Ɨhɨ pë mahunɨ raxa a reahuamahe. Kihamɨ, Yoawë, Omawë, Ruwëri, Pore kihi
si kɨ kuonomi. Ãyakorari pë yaia kure. Ãyakoari, Kuyarori oxe pë xo, pë
përɨoma. Ɨhamɨ hei si kɨ rii piyëwa notiwa. 

Yanomamɨ të pënɨ kihi, mo ha piyëprarɨhenɨ, pëma të pë waɨ hëami. Napë
iha napë a patanɨ mo kɨ ha pëtamarɨnɨ, pëma kɨ waɨ hëami. Ɨhɨ piyëwa. Pë
aka praukurayoma. Ɨnaha a kuprarioma. Ɨnaha të kua, hei të yai. 

Hekura pë makuinɨ pëma kɨ reahumou hirakemahe, Yanomamɨ. Reahu a rë
hiranowei, pëma kɨ reahumopë, yami kɨpɨ makui, ɨ̃hɨ të mɨ amo ha, kɨ
përɨpɨoma. 

Ɨhɨ të he tikëa ahete ha, reahumorewë pëma të tapë, yãɨmorewë pëma të
tapë, a rë hiranowei. Ɨhɨ weti naha kɨpɨ wãha kupɨoma, ɨ̃hɨ kɨpɨ wãha
hirapɨkemahe, kama wama kɨ iha? Ai të pënɨ napë pë iha, të wãrihihe ha,
kɨ pata rë përɨpɨo mɨ hetuonowei, kɨ wãha yupɨamahe? Ɨhɨ weti naha kɨ
wãha kupɨa? Reahu të rë hirapɨnowei, yãɨmou të rë hirapɨnowei. 

%\chapterspecial{O {pássaro} {siikekeata}}{}{}
 

\letra{K}{eora,} kɨrɨɨɨ, keora kɨrɨ!--- disse assim. 
Os dois fugiram ensinando o ritual do \emph{yaɨmou}. Apesar de serem
somente dois, faziam festa, ensinaram a encher os cestos de fruta
conori. Não sabiam matar anta, porém tinham muita anta moqueada. 

Como se chamavam os dois avós de Siikekeatawë? O neto deles se chamava
Siikekeatawë. 

\emph{Siikekeã, keã, keã}! Vocês escutam esse canto de passarinho? Era o
nome dele. Não era Yoawë. Era outro irmão mais velho. Yoasiwë era o nome
do irmão mais novo. 

As famílias yanomami em geral são numerosas. 

Esses dois apareceram na seqüência desta história. Sim, os dois
apareceram. Aquele que eles chamam Yoahiwë, aquele se
tornou \emph{napë}, os dois foram ao rio Tanape. Omawë e Yoahiwë se
tornarão \emph{napë}. Não foram esses dois que ficaram. Não foi Omawë, o
irmão mais velho e bonito. Nem Yoasiwë. Não foi aquele
passarinho \emph{uxuweimɨ} bonito que ficou.\footnote{ N\,E\,Ver as histórias de Omawë.} 

Trata"-se aqui de outros dois. Ficaram somente os de aparência triste e
velha. Esses dois ficaram.

Onde fica a foz do rio, cuja parte inferior olhamos? Onde se encontram
os dois rios? Foi nesses dois rios que os dois se dividiram. É assim. É essa a história dos dois que se dividiram. 

Esses dois irmãos mais velhos ficaram, ficaram fazendo festa, ensinaram
os descendentes a fazer festa. Terminou a festa, a festa acabou quando
juntaram a comida, colocaram a carne de anta em cima dos conoris.
Enquanto isso, eles cheiraram paricá. Ensinaram o ritual
de \emph{yaɨmou}, apesar de estarem sozinhos. Os dois conversaram,
fazendo o ritual de \emph{yaɨmou} no meio do xapono. 

--- Sou teu irmão e pergunto, como vamos falar? Nós vamos discutir e,
depois, nos fazer comer mesmo!--- disseram os dois, ensinando. 

Disseram o que dizem os Yanomami quando fazem o ritual de \emph{yaɨmou}.
Os dois deram exemplo aos Yanomami. 

--- Vamos encher a barriga até cair!--- falavam brincando, como se
houvesse muita gente ao redor. 

--- \emph{Aë, aë, aë, aë, aë, aë, heeeee}--- um deles diz. 

Enquanto dizia isso, o neto deles saiu levando o arco \emph{haowa}.
Ouviu a voz forte do pai. O menininho escutou a voz forte do seu pai, a
voz vinha subindo da curva do rio. Apesar de ele não ser de grande
tamanho, apesar de ele ser pequeno, a voz forte espantou os dois até os
expulsar. 

Enquanto flechava passarinhos, o neto correu até certa distância e,
enquanto flechava, a voz forte surgiu. 

Os dois avós faziam o ritual de \emph{yaɨmou}, se batendo no
corpo. \emph{Hëë, hëë, hëë, haëëë, haëëë, haëëë}! (Eu mesmo sou surrado
pelos meus parceiros ao fazer o ritual de \emph{himou}!) 

Enquanto faziam assim, o neto rodeava para matar passarinhos. Nesse
momento, a voz grande surgiu. A voz não era fina. 

--- \emph{Siiiikekeã, kea, kea}, rasgar a pele, rasgar, rasgar!--- dizia
a voz, chegando. 

A floresta estava tremendo com essa voz. 

--- \emph{Siiiiikekea, kea, kea}! dizia a voz, chegando. 

Apesar de ser pequeno, ele estava andando ali, ouvindo essa voz. 

Enquanto os dois avôs discutiam, o neto retornou. Na entrada, soltou o
arco. Soltou"-o de medo. 

--- Avô!Avô! Parem com esse barulho! A voz terrível do monstro se
aproxima! Ele vem rasgar a pele de vocês, a voz já está perto, avô,
parem!--- disse para os dois. 

Quando os dois pararam, assustados, pararam de repente, a voz surgiu
naquele instante. Quando ficaram silenciosos, ouviu"-se logo a voz
forte: 

--- \emph{Siiiiikekea, kea, kea}!--- disse aos dois. 

--- \emph{Hɨ̃ɨɨ}!--- gemeram os dois, e se levantaram assustados. 

O que lhes aconteceu?

--- Vamos! Vamos, querido! O monstro vai rasgar nossa pele, parente
querido, vamos, vamos, depressa!--- disseo mais velho, cujo nome já
pronunciei. 

Enquanto a voz dizia isso, o irmão mais velho também queria se
transformar. 

--- Vamos, venha, meu irmãozinho, venha, meu netinho! Vamos depressa! O
neto rodava na frente dos dois. 

--- \emph{Keora kɨrɨ! Keora kɨrɨ! Keora kɨrɨ}, vamos cair na água, lá
embaixo!--- disseram logo os dois. 

Os dois disseram isso, embarcaram e voaram acima da água. Lá, os dois
caíram rio abaixo. Enquanto os dois prosseguiam, a cauda vermelha do
neto saiu. \emph{Prohu}! 

O neto seguiu os avós, aqueles que voam. Por que ele tem esse nome de
Yoasiwë? Quando os passarinhos pousam em galhos fincados na água, a
cauda deles não é vermelha? O neto se transformou nesse passarinho. Onde
ele se transformou, onde foram os dois avós, a imagem do neto ficou, se
multiplicou e ficou voando acima das águas. Os dois avós caíram, levando
à frente deles seu verdadeiro neto, de quem restou somente a imagem.
Assim foi. 

Onde os dois avós caíram e onde a sua imagem ficou, se soterra o fim do
rio. É lá que os dois moram, eles não morreram e ainda moram lá. Eles
não morrem de doença. 

Os dois ensinaram o ritual de \emph{yaɨmou}, da mesma forma que faziam;
é por isso que nós, Yanomami, fazemos festas. Nós perpetuamos os
rituais. Não foram outros que nos ensinaram a fazer festa. 

 

%\chapterspecial{{Siikekeatawë}}{}{}
 

 

--- Keora kɨrɨɨɨ, keora kɨrɨ! --- a ha kurɨnɨ, kɨ rë tokupɨrayonowei të
tamahe, ɨ̃hɨ kɨpɨnɨ yãɨmou të yai hirapɨkema, yami kɨpɨ makui kɨpɨ
reahumoupɨɨ ha kuponɨ, momo, a yami makui të pë yorehi pata taɨ hirapɨɨ
ha, kɨ no xamapɨ mi makui, xama të pë pata rë tapɨamanowei. 

Ɨhɨ Siikekeatawë xɨɨ e kɨ rë kupɨonowei, weti naha kɨ wãha kupɨa?
Hekamapɨ e wãha. Siikekeariwë e wãha kuoma. 

Siikekeã, keã, keã! Wama të pë kiritamɨ kuɨ hiriaɨ tao? Ɨhɨ e wãha
kuoma. Ɨhɨ hei Yoawë, ɨ̃hɨ mai, ai rii pata yai, ɨ̃hɨ ai a rë kui oxe a
wãha, Yoasiwë oxe ai a wãha, ai pë marë Yanomamɨ tetei, kɨ rë përɨpɨtou
hërayonowei, ɨ̃naha kɨpɨ hëpɨtou kurayoma. 

Yoahiwë, të pë rë kuɨwei, ɨ̃hɨ kihamɨ hei të pë napëpropë, hëyëmɨ pata u
hamɨ, kɨpɨ hupɨnomi. Ai kɨ rii hupɨrayoma. Napë të pë kupropë. Hei kɨpɨ
rë hëpɨtore, Yoahiwë, Omawë pata a yai hei katehe a yai rë kui ɨ̃hɨ mai!
Ɨhɨ Uxuweimɨ si pë rë riëhëi ɨ̃hɨ mai! 

Të si pë puhi rë ahii, të si pë wai tapramou rë rohotei, ɨ̃hɨ hei kɨpɨ
përɨpɨa hëkema. Ɨhɨ pata e hëkema. Ei ai kɨpɨ rë kui, kihamɨ ai a
kurayoma, kihamɨ Totewë a kurayoma. 

Weti ha pei u kɨ he tikëpɨa kure yai? Pëma u pë koro rë mɨɨwei. Kama u
kɨpɨ weti ha u kɨpɨ he tikëa?Ɨhɨ u kɨpɨ hamɨ kɨpɨ rë xereronowei. Ɨnaha
të kua. 

Kɨ yai rë xererepɨprarionowei ɨ̃hɨ të rii. Ɨhɨ të rë kui hei pata e kɨ rë
hëpɨre, kɨpɨ reahumopɨɨ hëoma, notiwa të pë reahumou hiraɨ ha. Kɨpɨ
waikipɨprou, reahumou tiraprou, nii pë pata kõkapɨrarei, momo pë pata
heaka hamɨ xama pë pata makekepɨrei, kupɨaɨ tëhë, ẽpena a kopɨama. Të pë
yãɨmou hiraɨ ha, yami kɨpɨ makui. Kɨpɨ noã tayou, mɨ amo ha kɨpɨ
yãɨmapɨyoma, pata a makui xo. 

--- Hei aihë ya rë kui weti naha yahë kɨ kupë? Yahë kɨ noã tayou ta
kuparunɨ, yahë kɨ iamayou ta yaio këëë --- kɨ kupɨma, të hirapɨɨ ha. 

--- Weti naha Yanomamɨ të pë kuɨ, të pë yãɨmou tëhë? --- Të pë kuma? Ɨhɨ
e kɨpɨ kumahe. 

--- Pëhë e kɨ pëtɨrɨ praparei! --- kɨpɨ xomi kuma, ai të pë kuama mai
ha. Pë ĩtapɨpraapë pë kopemapɨma. 

--- Aë, aë, aë, aë, aë, aë, heeeee --- e hore kuma. 

Kuɨ tëhë, hei hekamapɨ e horepɨkema. Haowa e ha hayurëma. Ɨhɨ ya e wãha
rë yuprarɨhe, e horekema. E wãha rë kuonowei. Ɨhɨnɨ kama pë hɨɨ a wã
pata hĩriapë. Ɨhɨ pë hɨɨ e rë kuonowei, ihirupɨ wainɨ e ã pata hirirema.
Pei e të ã mɨ pata yamou kuimi, e të ã pata ãyoriaɨ kuimi. Ɨnaha të kasi
pata kuwëmi makui, të wai ma ihirupɨwei, të huxomi ha preonɨ, kɨ rë
yaxupɨre të kupropë. 

Hekamapɨ e rërëpɨpe kirioma, kiritamɨ pë wai niaɨ tëhë, pë wai tɨhɨyëmaɨ
tëhë, e të ã pata pëtarioma, yaɨmopɨɨ tëhë, pë xɨɨ kɨ ã karëhopɨɨ tëhë,
kɨ si rë paipɨyouwei: 

--- Hëë, hëë, hëë, haëëë, haëëë, haëëë, kɨ kupɨɨ tëhë, kiritamɨ e pë
napë tiporema. (Ware si marë yoaɨwehei, ware a himomaɨhe tëhë.) 

Kuɨ tëhë, të ã pata pëtarioma. Wã okaroonomi. 

--- Siiiiiikekea, kea, kea, keããããã! --- e të pata kutoimama. 

Urihi kë a pata haruharumotamama. 

--- Siiiiikekea, kea, kea! --- e të pata kutoimama. 

Të wai ma ihirupɨwei rë, kiharë të wai ma hure. 

Kuɨ tëhë, pë xɨɨ kɨpɨ noã tayou tëhë, a he tatoprarioma. Pei yo ka ha e
haowapɨ huheprakema. E kiriri huheprakema. 

 

 

--- Xoape! Xoape! Kɨ ã ta mapɨiku! Yai të ã no pata kiriaimi, wahë kɨ si
rĩya ikekaimi, të ã pata ahetea waikire, xoape, kɨ ã ta mapɨpo! --- e
kupɨma. E kɨ ã mapɨprapema kiriri, mapɨprao tëhë, e ã pata pëtou
nokarayoma. Kɨ mata waipɨprao tëhë, të ã pata nokakeyoruma. 

--- Siiiikekea, kea, kea! --- e të pata kupɨtarioma. 

--- Hɨ̃ɨ! --- e kɨ kupɨma, e kɨ kiriri hokëkëpɨprou ha xoaronɨ, ɨ̃hɨ weti
naha e kɨ kupɨprarioma?

--- Pei kë! Pei kë! Yai tënɨ pëhë kɨ si rĩya rë ikekaimi ë! Pusi pei
këë! Pei këëë! Pei a ta haɨpraru xë! Pata e kuma, ya e wãha rë
yuprarɨhe.

--- Haɨmo pusi, xei, hapoxë pei kë, haɨmoxë, hekamapɨ e mɨ wai
tarɨapraroma. 

--- Keora kɨrɨ! Keora kɨrɨ! Keora kɨrɨ! --- e kɨpɨ kuɨ xoaoma. 

Ɨhɨ rë e kure, porakapɨ kɨ rërëpɨa xoarayo hërɨma, pei u heaka hamɨ.
Kihamɨ kɨ kepɨ kiriopë. Kɨ kupɨa hërɨɨ tëhë, hekamapɨ ya e wãha rë
takɨhe, ɨ̃hɨ kɨ kupɨa hërɨpë hamɨ, e texina ĩyë. Prohu! 

Ɨhamɨ e pë nosi yaurayo hërɨma. Ɨhɨ rë pë yëɨ. Ɨhɨ exi të pë Yoasiwë?
Mau u pë hetu hamɨ, të pë texina rë ĩyëi? Ɨhɨ e kuprario hërɨma. Ɨhɨ xɨɨ
e ku hërɨpë hamɨ, a no uhutipɨ rë hëaanowei, mau u pë hetu hamɨ pë yëɨ.
Ɨnaha pë kuprarioma. Pei a yai rë kui, a kãi kepɨye kirioma. No uhutipɨ
a hëpɨmake hërɨma. Ɨhɨ rë a pararupɨre hërɨma. Ɨnaha të kuprarioma. 

Ɨhɨ kɨpɨ rë kepɨora kiri, kɨpɨ no uhutipɨ rë përɨpëye kirionowei hamɨ,
pei u koro titia. Ɨhamɨ kɨ përɨpɨa. Kɨ rë nomapɨno rë mai, ɨ̃hɨ kɨpɨ
xoaa. Ai kɨ ha wãritipɨprarunɨ, kɨ nomapɨnomi. Kɨ ha xawarapɨprarunɨ, kɨ
kãi nomapɨ taimi. 

Ɨhɨ kɨpɨnɨ yãɨmou, të rë hirapɨnowei, ɨ̃naha të tapɨma yaro, Yanomamɨ
pëma kɨ reahumou. Pëma të pou hëa. Ai të pënɨ reahumou tararenowehei, të
rë hiranowehei, kuami. Ɨnaha të kua.

\chapterspecial{A {árvore} {dos} {cantos}}{}{}


\letra{N}{ós vamos cantar.} No início, não havia canto, não havia, ninguém
cantava. Onde se erguia a árvore dos cantos, os dois foram caçar. Dois
moços Wakusitari -- dois não, um só moço, que a descobriu em sua região.

Os Katarowëteri eram os amigos dos Yãrusi, cujo líder se chamava Yãrusi.
Do outro lado da planície, eles, os Wakusitari encontraram a árvore dos
cantos. 

Outros dizem que foram os Koteahiteri que descobriram a árvore cantando,
e que chamaram os Katarowëteri para pegar os cantos. 

Graças à árvore, os Koteahiteri se enfeitaram com penas de cauda de
papagaio, pintaram"-se com elegância, colocando crista de mutum, e
dançaram. Era uma região bonita e plana onde crescia somente a planta
ária. Eles ocupavam uma bela região. 

Por isso, dois moços Koteahiteri foram caçar. 

--- Vamos entrar na mata, lá adiante! 

O irmão mais velho e o irmão mais novo foram caçar. A floresta parecia
mais baixa por causa da luz forte, como a luz do dia na roça. Foram
embora naquela direção, andando. Andavam no meio do brejo, andavam no
meio, ouviram os ecos dos cantos. 

Não havia sujeira no chão onde encontraram a árvore dos cantos dançando,
para frente e para trás. Havia somente areia bonita e muito brilhante. A
árvore dançava. 

--- Ãë, ãë, ãë, e, e, e, e, e, ãë, ãë, ãë, ãë!--- encontraram a árvore
cantando assim. 

--- Ë, aëë, ëaëë, ëaëë, ëaëë, ëaëë, ëaëë! -- cantava a árvore. 

Enquanto isso, o irmão Katarowëteri, o filho mais velho, disse: 

--- Õo, irmão menor! Dá pra ouvir um canto, lá onde há uma luz grande
acima do pântano, o som do canto vibra lá, escute isso! Provavelmente é
o som de um grande monstro! Esse som, naquela direção, mais adiante!
Vamos nos aproximar por ali, abrir um caminho no areal! Venha aqui!
Vamos, irmão menor! Vamos logo olhar de perto! 

--- Será voz de gente?--- disseram os dois. 

Onde a árvore dançava, a luz forte batia na areia bonita. 

--- Õoooãaaa! Vamos, irmão menor, vamos! A árvore dos cantos está
dançando, vamos, vamos, vamos até nosso pai, para avisá"-lo!--- disse. 

O irmão menor subiu em uma árvore bonita \emph{matomɨ} inclinada, para
ver se havia gente por perto, se via algum movimento, subiu e ficou
no alto. 

Ali, na areia, a luz brilhava de todas as cores, repousava bem no
centro, e a árvore dançava devagar para frente e para trás, cantando. A boca da árvore era bem bonita, e a árvore dançava para frente e para
trás. 

O irmão menor desceu e disse:

--- Õooãaaa! Irmão mais velho! Irmão mais velho! Nossa! Está lá cantando
e dançando, de uma maneira tão bonita, é a árvore dos cantos! Querido,
parece que essa árvore canta, essa árvore tem cantos bonitos! 

--- Vamos! Vamos até nosso pai! 

Os dois disseram e correram imediatamente. Chegaram correndo.

--- Prohu! Chegamos! 

Eles encontraram esse som e se enfeitaram por causa da árvore dos cantos. 

--- Meus queridos! Enfeitem"-se para pegarem cantos bonitos!--- disse o
líder dos Koteahiteri. 

O irmão mais velho fez o \emph{himou} com o pai, contando"-lhe sobre a
árvore dos cantos.\footnote{  O \emph{himou} é uma modalidade de diálogo cerimonial usada para trazer notícias, ou fazer um convite para uma festa.}

--- Tãrai! Ha! Meu pai! Pai! Olhe! Sou teu filho, olhe! Você não sabe
por que voltei logo correndo! Você nem sabe! Pai! Pai! Pai! Você nem
imagina o canto bonito que meus ouvidos ouviram! De arregalar os olhos!
Meu pai! Meu pai! Meu pai! Você que mora aqui, eu sou seu filho, eu não
lhe diria para proibir as mulheres se enfeitarem!--- disse. 

--- É claro! É claro! Queria ouvir isso mesmo, meu filho mais velho,
querido!--- respondeu seu pai. 

Fez o \emph{himou}: 

--- Vamos! Õoooãaaaaõoãaõoãa! Ele viu uma bonita árvore dos cantos!
Õõoo! --- gritaram. 

Ficaram animados.

 


\chapterspecial{{Amoa} {hi} ã {he} {rë} {haanowehei}}{}{}
 

Amoa pëma a tapë. Hapa amoa a kuonomi. Kuonomi, ai të pë kãi amoamonomi,
tëhë amoa kama hi rë upraatayowei hamɨ, kɨ ramɨ hupɨrayoma. 

Kutaenɨ hi ã eë hapɨrema Wakusitari a huyanɨ, kɨpɨnɨ mai, yami a huyanɨ.
Katarowëteri pë rë kui, Yãrusi kama nohi e pë wãha kuoma, Yarusi përɨamɨ
a wãha kuoma yaro. Ɨhɨ ai maxi yarɨ hamɨ Wakusitari pënɨ amoa hi ã he
haremahe.

Ɨnaha ai të pë kuɨ: Hei Koteahiteri pë yainɨ amoa kë hĩ ã he haamahe.
Katarowëteri kë pë ha nakarëhenɨ, amoa kë hĩ ã toamahe. 

Amoa hi nohi pauxiamahe. Werehi xina pata huuhamahe. Pë onimoma, pë no
aiama, ɨ̃kɨmo a huuhamahe, pë ha kuaanɨ, amoa hi nohi praɨamahe. Urihi
katehe kë, kuma kë masi he pata yarɨmoma, urihi katehe a pomahe. 

Pouhe yaro, ɨ̃hɨ Koteahiteri huyahuya kɨ ramɨ apɨyo hërɨma. Kɨ ramɨ ha
apɨro hërɨnɨ. 

--- Kiha pëhë kɨ ha paikutunɨ! 

Hei pë pata, hei pë oxe, ɨnaha ramɨ kë kɨ hupɨma, ramɨ kë kɨ apɨyo
hërɨma. Kutaenɨ, hɨ̃ɨɨ! e të xĩi pata yahatotoa hërɨma. Hikari kurenaha e
të xĩi pata kuaa hërɨma. Kuaa hërɨpë hamɨ, kɨpɨ katito hërɨma. Matotapɨ
hërɨma. 

Yãmaro kë xĩi pata hamɨ kɨ mɨ amopɨa hërɨma, mɨ amopɨa hërɨɨweiiiiiii,
mɨ amo yai ha amoa kë ã wa karëhoma, mɨ amo yai ha hẽka a praopë ha
kunomai, amoa kë hi tirurou he hapɨrema, makamaka katehe kë a pata
yaixĩi no aihimou totihiopë hamɨ, amoa kë hi tiruroma. 

--- Ãë, ãë, ãë, e, e, e, e, e, ãë, ãë, ãë, ãë! --- amoa e hi kupɨɨ he
harema. --- Ë, aëë, ëaëë, ëaëë, ëaëë, ëaëë, ëaëë! --- amoa e hi kupɨma.
Kuɨ ha, Katarowëteri pata a rë kui, ihirupɨ pata e rë kui: 

--- Õo! Õasi! Amoa a nohi karëhorati kihi të xĩi pata rë makerati ha,
kihi amoa, kiha amoa kë a morokaɨ kurati, yɨmɨka ta taprao, ɨ̃hɨ rë --- e
kuma --- Yai të ã pata pë wëë! --- pata e kuapraroma. Ɨhɨ Koteahiteri
pënɨ. --- Kihi të ã rë morokarati hamɨ, mihi të pata makamakapɨ rë
matoto piyëhërɨ hamɨ, wa yo ha reikimapaharunɨ, a ta ahehetetaru! ---
pata e rë kuyaronowei --- Pei! Oxei! Pëhë të ta mɨɨ ahetou xoao! 

--- Yanomamɨ rë të pë ã mata tawë! 

Kɨ noã tapɨyoma, amoa kë hi tirurupë hamɨ. Hɨ̃ɨɨɨɨɨ! Makamaka katehe kë e
xĩi pata makeoma. 

--- Õooãaaa, pei kë, oxei, pei kë oxei, amoa hi rë tiruropiyei ë, pei
kë, pei kë, hayë kë ihamɨ ëë, hayë pëhë a yɨmɨkamapëë! --- e kuma. 

Matomɨ katehe hi pata kaɨopë hamɨ, oxe e tukema, Yanomamɨ të pë mɨɨ ha,
të pë xurirou mɨɨ ha, e ha tuikunɨ, e ha tirehetarunɨ. 

Kihi makamaka kë a xĩi pata no aiwë makeaɨ kupiyei, mɨ amo yai hamɨ amoa
kë hi wa kãi opi tirutirumoma, të hi kahikɨ no aiwë no kirii, e të hi
tiruroma. 

 

--- Õooãaa! --- e ha nihoroto hërɨnɨ --- Apa! Apa! Kurahë katehe kë të
wã kãi tirurou kuopiyei. Apa, amoa kë hi ë! --- a kuma --- Pusi amoa kë
hĩ ã no taië, pusi amoa katehe kë hĩ ã rë taië! --- e kuma. 

--- Pei kë! Hayë kë ihamɨ ë! Kɨ ha kupɨnɨ, kɨ rërëpɨa xoape hërɨma. Kɨ
ha rërëpɨpo hërɨnɨ: 

--- Prohu! --- e kɨ kupɨma. 

Amoa hi nohi pauxiaɨhe ha: 

--- Pusi pei kë pë ta pauximo xë! Amoa katehe wama a toapë! --- përɨamɨ
Koteahiteri e kumahe. 

Hi nohi himopɨama pë patanɨ e hi nohi himoama, amoa hi wãha nohi wëaɨ
ha, pë hɨɨ iha. 

--- Tãrai! --- e kuma --- Ha! Napemi! Napemi! Ha! Hei yarohë ya rë kuii,
ha! Weti wa të taɨ ha, wa të rërëɨ mɨ yapa a ta kuponɨ? Wa puhi kuoranɨ
ha kunomai! Napemi! Napemi! Napemi! Hei ya yɨmɨka ha amoa katehe ya rë
hiritaɨwei! Ya mamo rë ikeketouwei, napemi! Napemi! Napemi! Hei kɨ suwë
rë kui, kɨ pauximomaɨ mai! E roa yai a ta përa! Yarohë ya rĩya kuoranɨ
kunomai! --- e kuoma. E kuɨ ha: 

--- Hão! Hãooooo! Noa taɨ yai a ta përaxëa. Pusiwë! Pusiwë! Ha! Ɨnaha rë
kë, ɨnaha rë kë --- pë hɨɨ e kuma. 

E nohi himoama. 

--- Pei kë! Õoãaõoãaõoãa! Amoa katehe hi he hõra rë harenowë! Õoooo! ---
e pë kuma. 

E pë xi wã kãi toaama.

 

 

 

 

 

\chapterspecial{{Monstro} {këyakëya}}{}{}
 

\letra{H}{avia também} os que viviam na região centro"-sul, os
Yãimoropɨwei, que ficaram presos, pois moravam dentro da terra com o
monstro Këyakëya--- que, portanto, não era gente. 

Os que asfixiaram Këyakëya existiam bem antes de nossos antepassados.
Këyakëya morava dentro da terra, na vizinhança do xapono dos
ancestrais.\footnote{ Os xaponos são as casas coletivas circulares onde moram os Yanomami. Cada casa corresponde a uma comunidade; em geral não se fazem duas casas numa mesma localidade.} 

Apesar de ser um monstro, Këyakëya era líder dos Yãimoropɨwei. Os
companheiros de Këyakëya moravam dentro da terra e a casa deles tinha um
respiradouro, como o da casa do tatu. A casa de Këyakëya também tinha um
respiradouro. Moravam ali também os Motuxi, que se dividiram e se
espalharam. 

Os Prãkiawëteri asfixiaram Këyakëya, tentando matá"-lo. Asfixiaram"-no,
foi assim que nos ensinaram a matar. Eles não o mataram com flecha. 

No início, não havia matança, não havia inimizade, não havia briga
mortal. Os \emph{napë} também não existiam.\footnote{  O termo \emph{napë} designa os estrangeiros, em geral os brancos, ou quem adotou seus costumes.} Os nossos antepassados não sabiam manifestar ira nem raiva. 

Ele conseguiu escapar sob a forma de espírito. Ele não se transformou à
toa. Os companheiros dele, como nós, sempre padeciam de fome; todos
morreram pela fumaça que entrou no buraco. 

Këyakëya nos legou o sentido de vingança por causa da filha de quem?
Qual é o nome do pai cuja filha foi vítima da crueldade de Këyakëya, que
chegou e entrou no xapono? A vítima que, brutalmente, Këyakëya fez
descer da rede e sair era a filha do líder dos Naɨyawëteri. Era uma moça
bonita, realmente muito bonita. Ela estava na primeira menstruação, e
mesmo assim, ele a arrancou da reclusão.

Apesar de ser monstro, Këyakëya existia e vivia como gente. Como
morava dentro de um buraco, depois de trucidar a menina menstruada, ele
e os demais membros do grupo foram asfixiados pelos Prãkiawëteri. Mas
apenas Këyakëya conseguiu fugir, se tornando eterno na forma de
espírito. Ele ainda existe como espírito. 

Naɨyawë desgalhava um pé de fruta \emph{naɨ}\footnote{ Segundo Lizot, uma balateira, \emph{Manilkara bidentata}.} em uma roça
distante. Aooo, aoooo, aoooo, aooo! Fazia assim para sua gente. 

Enquanto eles comiam a fruta \emph{naɨ}, Këyakëya arrancou a menina do
seu recluso, matou"-a e a devorou. Ele a comeu sozinho. 

Fez lascas pequenas da carne das demais crianças, que também havia
trucidado, para oferecer a todos seus companheiros. Amontoou as lascas de
carne que ele colocou no seu grande cesto, chamado \emph{yotema}.
Carregou todos os restos das crianças massacradas e levou junto o irmão
da menina menstruada, que estava vivo e bonito. Ele o fez sentar em cima
dos cadáveres dentro do cesto. 

O menino vivo, que ele levou, transformou"-se em papagaio durante o
percurso. Këyakëya saiu do xapono dos Naɨyawëteri e andava a passos
largos, foi então que o menino, já de longe, disse: 

--- Kuao! Kuao! Kuao!

Esse som se tornou o som dos papagaios. Esses pássaros voam; ele pousou
em um galho e assim ficou. Këyakëya olhou para a beira do cesto,
querendo ver se o menino ainda estava sentado. Fez o filho de Naɨyawëse
tornar papagaio. Como o menino não estava, ele retornou àquela direção.
O menino se tornou a imagem do papagaio que grita: Kuao! Kuao!
Kuao!

--- Ouça! Meu xerimbabo! Onde você pousou? Kuato, kuato, kuato! ---
disse Këyakëya voltando e correndo.--- Em qual paragem você
ficou? Kuato, kuato, kuato!

--- Õiyaoooo!--- disse o papagaio. 

Assim disse aquele que, apesar de ser filho de gente, tornou"-se
papagaio. 

É a história dos antepassados. Tambem existiam monstros com outros
xaponos, sendo essa a história de Këyakëya e dos Yãimoropɨwei, que
moravam em xaponos pouco distantes um do outro. 

Depois, aparecerá o nome do rio que tirará e levará muitos ancestrais
Yanomami. É somente depois da história dos Yanomami levados pelo rio que
vem nossa história. Os Waika a contam de uma maneira diferente, eles a
contam conforme seus antepassados lhes contaram.\footnote{   O par \emph{waika}/\emph{xamatari} parece ter sido usado originalmente para designar outros grupos yanomami vivendo em região geográfica diversa de quem fala, os primeiros ao norte e oeste, e os segundos ao sul, reconhecendo-se neles conjuntos de características que os particularizam. Os termos foram atribuídos em diferentes momentos pelos brancos para designar grupos específicos de forma estável e, no caso de \emph{xamatari}, para designar a própria língua do tronco yanomami usada pelos Parahiteri que fizeram este livro.} 

Os companheiros de Këyakëya não sobreviveram, morreram todos pela
fumaça. Eles os asfixiaram a todos, somente Këyakëya sobreviveu, se
transformando em espírito eterno. Esse sobrevivente alcançou o xapono dos
espíritos, pois se tornou um deles, quando ainda eram Yanomami e moravam
como nós. Ele os alcançou e ficou lá. 

Não mora mais onde o asfixiaram. Somente restou o marco dele. Não pensem
que os companheiros de Këyakëya sobreviveram e se agruparam enquanto ele
alcançava os espíritos! 

Não houve sobreviventes do grupo dos Naɨyawëteri. Acontecerá depois. Os
sobreviventes eram os que afundaram, não os outros antepassados. As
águas sobem devagar e os que afundam são os únicos sobreviventes. 

Depois, os que tinham o mesmo nome que as montanhas também
sobreviveram.
\chapterspecial{{Këyakëya}}{}{}
 

Kama pë rë kuonowei koro ha mɨ amo ha, pë xi rë wãrionowei, pë rii rë
titionowei, yai tënɨ pë kãi titioma. Këyakëya, Yanomamɨmi makui, a
përɨoma. 

A rë yarënowehei. Kamiyë pëma kɨ no patapɨ përɨo mao tëhë, të pë rë
përɨo xomaonowei të pë wãha xomaa. Ɨhɨ pata pë yahipɨ he tikë ha,
pëɨxokɨ ha, yai të titioma. 

Këyakëya përɨamɨ a wãha, yai të makui. Ɨhɨnɨ Yãimoropɨweiteri pë kãi
përɨoma. Këyakëyanɨ pë kãi rë titionowei, mahu hẽremopɨ kuoma, opo pë
hẽremopɨ rë kurenaha Këyakëya yai të hẽremopɨ kuoma kutaenɨ, Motuxi pë
pata xereremou piyëkëmoma kutaenɨ, kama e pë kãi rë përɨonowei,
Motuxiwëteri pë kãi titioma. Këyakëya ei pë wãha. 

Wetinɨ Këyakëya pë kãi rë titiaɨwei, weti naha pë wãha kuoma? Ɨhɨ
Prãkiawëteri pënɨ Këyakëya a yarëmahe, a xëpraremahe. Ɨhɨ pënɨ pëma kɨ
ixou hiraɨhe ha, Yãimoropɨweiteri pënɨ Këyakëya a unokai yarëmahe. A xëprapehe, a yarëmahe, kamiyë pëma kɨ xëprayopë. A nianomihe. 

Hapa niayou të kuonomi. Pëma kɨ napëmayou, xëpraɨ të kuonomi. Napë pë
makui, pë kãi kuonomi. Pëma kɨ nohi patama waiterimou taonomi, huxuo kãi
taonomi, ɨ̃hɨ pë xëremahe. 

Këyakëya a xëpraɨ puhioma makuhei, a xëpranomihe. A hekura tokua he
yatirayoma, kama a kuprou pëonomi. Hei kamiyë kureneha kuwë të pë no
xĩro preaama, të pë hɨtɨtɨwë nomarayoma. 

Këyakëyanɨ weti tëëpɨ noã prearema? Këyakëyanɨ pë tëë e napë rë
itorayonowei, e napë rë harayonowei, weti naha pë hɨ̃ɨ e wãha kuoma?
Naɨyawëteri ihirupɨ, tëëpɨ noã prearema. Kama përɨamɨ Naɨyawëteri a yai
kuoma. Suwë katehe a yai kuoma. Pë tëë e yai riëhëoma, kamiyë kurenaha
mai! Ɨhɨ suwë katehe yɨpɨ a ha ukëa he ha yatirënɨ, a noã prearema. 

Ɨhɨ a përɨoma, Këyakëya a rë përɨonowei, yai të makui Këyakëya a
Yanomamɨ përɨoma. A titioma kutaenɨ, ɨnaha të tama yaro, yɨpɨ hena
xëpraɨ xi ha wãrironɨ, kama e pë rë kui, e pë no ha prerarunɨ, e pë ha
yarërarɨhenɨ, kama a rë kui a parimi hekura tokua xoarayoma. Ɨhɨ a hëa
xoaa, hekura. 

Kihi hikari a rë kurahari naha, naɨ a pehi pata tihetimamahe ha:

 

--- Ãooo, ãoooo, ãoooo, ãooo! --- Naɨyawëteri e pë kuma. 

Naɨ a waɨhe tëhë, a ukëa hearema. A xëprapë, a wapë. A warema. Yaminɨ a
warema. 

Nakaxi yãhi pë wai ha tanɨ, tanɨ, kama urihiteri pë haikama, pë
topërarema. E yãhi kɨ ha orɨhenɨ. Pë ihirupɨ pë no mapraɨ hearayoma
yaro, hɨtɨtɨwë pë mɨ këa heararema, pë yehire hërɨma, kama yotema e
hamɨ. Ihiru e pë no payeri rë tapraɨwei, pë titire hërɨma. Ɨhɨ pë taɨ
makure, a rë përɨaɨwei, yaɨpɨ rë këprarɨhe, ihirupɨ e rë kui a yure
hërɨma. Temɨ. A rë riëhei. E tikëmare hërɨma. 

Ei a rë yurehe, kiha a kãi kutou tëhë, werehi e kuprarioma. Hëyëha a kãi
rë hare, a kãi rë rahurahumoimati, kihi karexi si rë prarahari naha a
kãi kutou tëhë: 

--- Kuao, kuao, kuao! --- të pë werehi rë kuuwei a no uhutipɨ
kuprarioma, të pë yëɨ ha piyei kunɨ heinaha e ha waroroikunɨ, e kasikɨ
mɨprarema, Këyakëyanɨ, e tikëa mɨɨ ha. Përɨamɨ e ihirupɨ
werehipramarema. A maa ha, e wã kãi yëa mɨ yapakema: 

--- Tãrio, weti ha wa hore piyëkei kuhe? a wãti. Kuato, kuato, kuato!
--- e kuɨ mɨ rërëa mɨ yapakema --- Weti ha a hëprario kuhe? Kuato,
kuato, kuato! 

--- Õiyaoooo! --- e kurayoma. Yanomamɨ ihirupɨ kuoma makui, e kua
topramarema, ɨnaha e kuma. 

Ɨnaha të ã kua, pata të rë kui, ɨnaha të pë kuaama. Ɨhɨ yai tënɨ pë kãi
përɨoma, ɨ̃hɨ tëhë ai të pë rë përɨonowei, ɨ̃hɨ të ha, hei Këyakëya a rë
yarënowehei, Yãimoropɨweiteri pë hirao he paoma. 

Ɨhɨ ei rë pë rë kui, waiha pei rë u kë wãha rë taore hamɨ, pë rë
pakakumare të wãha kupropë. Ɨhɨ rë të he tikë hamɨ, të he tikëatayoa, të
he tikëa kurati, komosi të ã yai. Hei Waika pë rë kui, maa, kama e të
rii taɨhe. Kama pë no patapɨ wãha rii tao. Kamiyë pëma kɨ no patapɨnɨ të
ã rë wëyënowehei ei të ã rii. Të rii ma rë yaitapraruhe. Ɨhɨ kama, ɨ̃hɨ
të rii maxi hamɨ, ɨ̃naha të kuoma, hapa të pë rë përɨonowei. 

Ɨhɨ Këyakëya urihi teri pë rë kuonowei të maxi hamɨ pë rë kuonowei,
Këyakëya urihiteri pë rë kui, ai e pë hëpronomi, kaɨ wakë xinɨ, pë
xëpraremahe. Pë yarëpraɨ haikirayomahe yaro, ai pë temɨ haimi, hei pë rë
kui, Këyakëya kama a rë kui, a hekura ha parimiprarunɨ, a hëtarioma. A nomanomi. Ei a rë kui, hekura pë iha a waroopë, a hëprario kuhe,
Këyakëya hekura. Yanomamɨ pë kuo tëhë, hei pëma kɨ rë kurenaha hekura pë
përɨoma. E warokemahe. Ɨha a kuopë. 

Kama a rë yarënowehei ha a titia xoaami. Hekura yai të pë iha a
warokema. Pei uno kua hëa. Hei pë rë kui, hei pë rë hëpraruhe, hei pë rë
kui, Këyakëya hëyëmɨ e kua rë xoarahai weti naha kuwë të pë hiraopë ha,
Këyakëya a warokema, pë puhi kuu mai! A warokema. 

Naɨyawëteri pë rë hëre, hei pë hëtopë mai! Hei pë mixi rë tuore, pë xĩro
hëprario, ai të pë no patama hëpronomi. Wãisipɨ, ĩsitoripɨ të u wai rë
õkimouwei, pë mixi rë tuowei pë rë kure, hei kë pë. 

Ɨhɨ pë rë kui, te he tikë hamɨ ai pë rii, pei ma pë rë kui, ma pë wãha
rë yehiponowehei, ɨ̃hɨ pë kãi hëprarioma.

 

\chapterspecial{O {surgimento} {das} {cobras}}{}{}
 

\letra{N}{essa época também}, as cobras não rastejavam como rastejam hoje, elas
viviam como os Yanomami. Transformaram"-se onde desceu o Sangue da Lua,
na floresta. Lá, caíram as cobras que picam. Transformaram"-se em cobras
lá em cima, enquanto iam para uma festa. Hoje, quando vocês olham para o
céu, vocês veem o peito daqueles que se transformaram em cobras. Não
havia cobras, nem jiboias, nem sucurijus. Os poraquês não existiam, nem
os peixes. Nós comemos a carne de gente.

Eles se transformaram em cobra, não no xapono, mas nesta floresta mesmo.
Foram chamados e foram lá, Wataperariwë e Jibóia, o irmão mais velho.
Foram lá longe com as Cobras, mas se transformaram na floresta. Eles,
então, não foram dançar. 

Com a cabeça coberta de penas brancas, dessa mesma forma que nós nos
pintamos, cada um pintou seu corpo com listras diferentes. As Cobras
moravam na sua própria região, como gente. Transformaram"-se quando
foram convidadas a dançar. Elas antes viviam como gente. 

Quem eram os dois tuxauas? O irmão mais velho e o irmão mais novo
moravam com as Cobras. Os dois também foram dançar. Watawatariwë e
Jiboia moravam com seu grupo, as Cobras. Jiboia era o irmão do meio.
Watawatariwë era o caçula. Os dois irmãos mais velhos eram esses: Jiboia
e Sucuriju, que nasceu primeiro. Aqueles que se pintaram eram três, pois
havia também Watawatariwë, o caçula, por isso nós nos pintaremos assim. 

Os Parawari também viviam com eles. Por causa deles se metamorfosearam,
porque os Parawari os levaram. 

Todos eles moravam em frente à serra Wãyapoto, que ainda tem esse nome.
Ocupavam essa região ao pé da serra na planície. Eram todos bonitos. É o
nome da região onde moravam os antepassados. É o verdadeiro nome dessa
região. As Cobras bebiam a água do rio Wãyapo, tomavam banho, se lavavam
nesse rio bonito. Tomavam banho e bebiam água. 

Eles nos ensinaram, assim, a dançar, mas, infelizmente, se
metamorfosearam. Eles iriam dançar, mas, infelizmente, se transformaram.
Iriam dançar. Transformaram"-se em cobras imediatamente. Tornaram"-se
cobras. Não foram dançar no xapono de outros. 

Em que xapono iam dançar? No xapono daqueles que se transformaram, que ainda existe na terra plana. Aqueles que se transformaram, apesar de se
pintarem fora do xapono, sofreram a metamorfose, transformaram"-se em
cobras. 

Os que convidaram as Cobras, como se chamavam esses antepassados? Eles
gritavam enquanto cozinhavam o mingau de banana para os visitantes. 

--- Por que estão agindo assim?--- Perguntaram"-se.

Pareciam gritar de propósito. Transformaram"-se perto do xapono dos
Jalouaca. Transformaram"-se perto desse xapono. Transformaram"-se. Os
antepassados se chamavam assim, Jalouaca. Assim se chamava o líder. Eram
espíritos, são nomes de espírito. Eram Yanomami e moravam como os
Yanomami. 

Apesar de morarem, assim, como nós, após a metamorfose em cobra eles não
voltaram à condição de seres comuns. Pintaram"-se fora do xapono dos
Jalouaca, pensando: 

--- Os Yanomami se pintarão assim! 

E se pintaram com listras. Pintaram"-se, na parte superior do braço, com
cor de sangue preto, igual à cor de meu irmão mais novo, como a cor de
seu braço. As cobras \emph{maraxari} se pintaram assim; a cobra coral
também se pintou com manchas vermelhas. 

O segundo grupo do xapono das Cobras se pintou em outro lugar, distante,
para que aquelas do outro grupo, que se achavam bonitas, se zangassem.
Elas tomaram banho no rio Wataperari, cuja água era branca. Ficaram onde
brilhava a luz. Assim era a luz do rio. Perto do xapono dos Jalouaca,
havia o rio, o rio apareceu de repente. 

Um pouco longe do xapono, as outras Cobras se pintavam juntas. 

Pintaram"-se. No segundo grupo havia uma mulher. Os bonitos desse
grupo, eram muito bonitos, chegaram até as outras cobras. Chegaram
também com eles dois Parawari bonitos, todos eram muito bonitos.
Chegaram. A beleza de suas pinturas incomodou os outros, que ficaram com
inveja. Chegaram, enquanto os outros se pintavam com riscas. Aquele,
cujo nome eu dei, apareceu no meio deles, Sucuriju. Ele, o irmão mais
velho, estava ao final dos que chegavam, aquele que tem grandes
desenhos.

--- Hɨ̃hɨ̃! Wĩsa! Wĩsa!--- assobiaram. 

Os do primeiro grupo, ainda se pintando, viraram a cabeça para olhar em
direção das cobras bonitas chegando, e disseram felizes:

--- Acabei de me pintar desse jeito! 

Apesar de não terem dentes como os dos Yanomami, depois de se
transformarem em cobras, depois da metamorfose, os dentes saíram. No
início, não havia cobra, aquelas que picam não andavam no chão, não
havia cobra"-surra, nem coral, nem cobra \emph{maraxa}, nem
cobra \emph{huwëmoxi}. Não havia nenhuma dessas cobras. Lá, onde os
bonitos estavam se transformando em cobras, houve um barulho tão grande
como o de um bando de queixadas, pois as cobras estavam surgindo. As
jararacas, as surucucus, as cobras papagaios e as
cobras \emph{waro} tambémsurgiram. Invadiram toda a floresta. Assim foi. 

Aqueles que haviam convidado as Cobras, os Jalouaca, por causa dos quais
aconteceu a transformação, subiram também ao céu no lugar da
transformação. Os bonitos estavam suspensos. Torurururu! E
trovejou. --- Prohu! --- Chegaram lá. Não estão aqui, nessa terra,
pois andam lá. Queriam viver saudáveis, então estão lá, saudáveis. Não
ficam em baixo. Ficaram em cima. 

Quando as Cobras subiram, o que aconteceu com os amigos delas, os
Jalouacas? Transformaram"-se também em cobra. 

Então, os líderes do primeiro grupo, que se transformaram também
em cobras, ficaram na terra.
\chapterspecial{{Të} {pë} {rë} {oruprarionowei}}{}{}
 

 

Oru pë kãi hunomi, oru Yanomamɨ kurenaha pë kãi përɨoma, pë kãi kuoma.
Ɨhɨ, kihamɨ, Pẽripo ĩyë rë itorati hamɨ, urihi hamɨ, kiha pë xi
wãrihotayoma. Ɨha pë oru rë kerayonowei, të pë si wëyëɨhe, oru pënɨ.
Heaka hamɨ, pë xi rii wãrihipraritayoma. Pë praɨaɨ mɨ ha hurɨnɨ, pë xi
rë wãrihonowei, hei wama parɨkɨ mɨɨ. Oru pë hunomi, hetu pë hunomi,
wãikoya pë kãi kuonomi, yahetipa pë kãi kuonomi, yuri pë kãi kuonomi,
Yanomamɨ wama të pë yãhi kɨ waɨ. 

Urihi ha pë ha nakarehenɨ, pë huɨ ha kuikutunɨ, Wataperariwë, Heturiwë
pata xo Oruri pë kãi hupɨɨ ha kuikutunɨ, urihi ha pë xi wãrihoma, pë
praɨɨ kateheonomi. 

Pei të pë horoimo pë ha, të pë ma rë yãmouwei, pei të pë pata yãprutaaɨ
yaitaama, Oruri pë ha oraora ya të wãha takema, korokoro pë wãha kuami,
pë xi rë wãrihonowei. Ɨhɨ Oruri pë rë kui, kama pë urihipɨ ha, Yanomamɨ
kurenaha kamiyë pëma kɨ rë kurenaha pë përɨoma, pë xi wãrihiprarioma. Pë
xi wãrihopë makui, pë praɨɨ mɨ ayoma, pë ha xorehenɨ, Yanomamɨ pë përɨo
parɨoma. 

Ɨhɨ exi e të përɨamɨ kupɨoma? Pata, pë oxe. Oruri pë kãi rë
përɨpɨonowei. Ɨhɨ pë kãi praɨpɨɨ mɨ kãi rë hurayonowei, Watawatariwënɨ
pë kãi përɨoma, Oruri, Heturiwë xo. Heturiwë pata e wãha yai, pëɨxokɨ
hamɨ ke e. Pata ɨnaha e kɨ kupɨa hei, pë xĩro. Wãikoyariwë pei a haa
xomarayoma. Ɨnaha pë kua. Ai, ai, ai pë kuoma. Wãikoyariwë e kãi kua,
kama pë rë onimonowei, kamiyë pëma kɨ onimopë. Wãikoyariwë, Hetu,
Watawatariwë oxe e wãha, suhe u haikatimɨ. 

Parawari pënɨ pë kãi hiraomahe, Parawari pë xo pë hiraoma, ɨ̃hɨ pënɨ pë
rurure hërɨmahe yaro, Oruri pë xi wãrihamapehe, katehe kama pë xĩro
hirao yarɨtaoma. 

Wãyapoto a parɨkɨ ha pë hiraoma, ɨ̃hɨ Wãyapoto a parɨkɨ ha pë përɨa xoaa.
Ɨhɨ e wãha kua xoahe. Yarɨ ha, ɨ̃hɨ kɨ tëhë pë përɨoma, a urihi pomahe,
kama pë urihipɨ wãha. Pata pë rë përɨanowei, të wãha urihi yai. Oruri
pënɨ u rë koanowehei, Wãyapo u koamahe. Ɨhɨ u yaruamahe. Wãyapo katehe u
yaruamahe. Pë rë yãrɨmonowei, u rë koanowehei. 

 

Kamiyë pëma kɨ praɨpë, ɨ̃hɨ të rë hiranowehei, ɨ̃hɨ pënɨ të praɨɨ
hirapehe, pë hurayoma makui, pë yaitaaɨ tikooma. Pë xi wãrihou tikoopë
makui, pë hurayo hërɨma. Pë praɨɨ mɨ ayo hërɨma. Kama pë oruriprou
xoarayoma. Pë oruriprarioma. Ɨhɨ ai të pë iha pë praɨɨ mɨ hunomi. 

Weti pë iha pë praɨpë pë hurayoma? Ɨhɨ pë xi rë wãrihonowei yarɨ ha,
xapono pata a praa xoaa, yarɨyarɨ të ha. 

Oruri pë rë xoanowehei, pë pata rë hiraonowei ɨ̃hɨ weti naha pë wãha
kuoma? Ɨhɨ pei pë xi yai rë wãrihonowei, sipo ha pë yãmou makure, pë no
rë Oruri preaanowei, pë rë oruriprarionowei. Pë rë yaitaanowei. Të kɨ ã
si pata ma hipikitapiyei, pë kuratapɨ u hariihe ha, të kɨ ã si pata ma
potehetapiyei makui: 

--- Weti naha pë pata kuaaɨ tikoa kupiyei?

Pë nohi kuaama. Ixarowëteri pë iha, pë xaponopɨ ha pë xi wãrihoma.
Ixaropɨwëteri pë përɨoma. Pë xi wãrihoma. Pata pë rë kui, ɨnaha pë wãha
kupramoma, Ixarowëteri. Ɨhɨ përɨamɨ a rë kui ɨ̃naha rë a wãha kuoma.
Hekura pë përɨoma, hekura pë wãha. Yanomamɨ rë pë kuoma. Yanomamɨ
kurenaha pë përɨoma. 

Hei kurenaha pë përɨoma makui, pë poreriprou kõonomi. Ixarowëteri ɨ̃hɨ
oruri pë xi rë wãrihamanowehei pë wãha. Ɨhɨ pë xaponopɨ sipo ha, pë ha
yãmorɨnɨ: 

--- Ɨnaha pë kuaaɨ hëopë tao! 

Pë puhi ha kunɨ, pë tiprutaama. Oxeyë kihi ĩxi kurenaha, wakë poko kɨ
hĩia rë kurenaha, hei ora ĩxi hĩia rë kurenaha, hei koro ĩxi rë
kurenaha, pë yãmou kuaama. Maraxari pë kuaama, huwë moxi pë kuaama, hei
kurenaha pë wakë rukëkoma, yamixano. 

Katehe të rë huxutamarenowei pei pë yãmou hëoma. Ɨhɨ kama Wataperari
kama u ha, pë yãrɨmou hëkema, u wai au, të u xĩi wai praapraamopë ha, pë
hëkema. Heinaha u xĩi kuoma. Ɨhɨ Ixarowëteri pë xaponopɨ ahete ha, e u
kuoma, e u pëtariomahe. 

Hei kamiyë pëma kɨ rë titipiyei hiramorewë nahi ha, kihi Oruri pë yãmou,
ɨ̃naha pë hirao kuoma. 

Pë yãmoma, katehe kɨpɨ yai rë kui, pë pëtarioma, ɨnaha pë pëtou
kurayoma, hei, suwë mahu a, hei Parawari katehe kɨpɨ. Ɨnaha kama pë xĩro
kuoma. Katehe pë yai rë kui. Ɨhɨ pë rë huxutore, pë mɨa kãi no rë
preaare, pë mɨ tikëtikëpraroma e pëtariomahe. Ɨhɨ hapa ya wãha rë
yuprarɨhe e parɨomahe. Noha hamɨ Wãikoyariwë e kuoma. Hɨtɨtɨ, pata e
nohapɨ aimama, pë oni pata rë prei. 

 

 

--- Hɨ̃ɨ! Wɨ̃sa! Wɨ̃sa! Pë husi he ã pë mamo xatipraamapehe. 

--- Hei, ɨnaha ipa ya të taawaikike kuhe! --- pë kuɨ topraroma. 

Yanomamɨ kurenaha pë nakɨ kuonomi makui, pë ha oruriprarunɨ, ɨ̃ha pë xi
ha wãrihiprarunɨ, pë nakɨ hararioma. Hapa oru a kãi hunomi, wa si rë
wëaɨwehei të kãi praonomi, Wãyapotorema pë kãi kuonomi, miomaakahe pë
kãi kuonomi, maraxa pë kãi kuonomi, huwë moxi pë kãi kuonomi, kuonomi.
Ɨha pë xi rë wãrihore, ɨ̃hɨ pë xi rë wãrihorati ha, katehe pë xi wãrihopë
ha, hawë warë kɨ pata hõra kuprarioma. Oru pë kuprou yaro. Karihirima,
pẽreima, arawaomɨ, waro pë kãi pata kurarioma. Ɨha pë rë kuaare ɨ̃ha hei
a urihi rë kui, a haikiremahe. Ɨnaha pë kuprarioma. 

Pë xi rë wãrihamarahei, ɨ̃harë, kama pë xi rë wãrihiprore ha, pë
heakaprario hërɨma. Heaka hamɨ kama pë kurayoma. Kama katehe pë rë kui,
pë pehi sutihiprou yaro Torururururu! Yãru e kurayomahe. Prohu! Kihamɨ
pë kuketayoma. Hëyëmɨ pita ha pë kuami. Kihamɨ pë huɨ. Katehe pë yai
përɨo puhiopë yaro, katehe pë kua kurati. Pë pepiami. Pë heakaketayoma. 

Ɨnaha pë ha kuprarunɨ, ɨ̃hɨ ɨ̃naha pë ha kuprarunɨ, weti naha norimɨ e pë
rii kuaama? Ɨnaha e pë riikuprou mɨ heturayomahe. 

 

\chapterspecial{A onça e a centopeia}{}{}
 

\letra{N}{essa época} as onças não comiam gente, não andavam, não existiam. Não
havia onça na floresta. Daí essa história. Não andava onça por aí para nos
matar e nos comer. Hu, Hu! Hu! A onça não dizia isso. 

Quem encontrou a primeira onça? Sozinha, ela sofria de fome, sequinha,
sua barriga gritava de fome, pois ela não tinha dente. Onça tinha apenas
gengivas, ela não mastigava, ela andava magra no meio dessa região do
Xererei, ela andava sozinha, andarilha, faminta. Como ela não comia quase
nada, ela chorava. Ela chorava por fome de carne. 

Quem a encontrou? Onça chegou onde estava Centopeia, onde morava
sozinha como gente, Onça chegou à casa de Centopeia. Ela apareceu, elas
se encontraram, ela ia de encontro. Com fome, andava como se fosse cego,
sem olhos, sofria mesmo, fazia muito barulho, tropeçava de fome. 

É uma centopeia! Vocês conhecem esse nome? Era gente, aquela que anda
sem fazer barulho. Krihi! Ninguém mais faz esse barulho, andando em cima de um pau.
Foi ela quem ensinou primeiro. 

Ela emprestou seus pés para Onça não fazer mais barulho; ela o ensinou a
andar discretamente. Depois do ensinamento de Centopeia, Onça andou, ela
foi lá, chegou à terra plana e desceu.

 

\chapterspecial{Ɨ{ra} {xo}, {wapororitawë} {xo} kɨ {he} {haap}ɨɨ}{}{}
 

 

Ɨhɨ tëhë, kamiyë ɨra pënɨ pëma kɨ waɨ maopehe, ɨ̃hɨ a kãi hunomi. Ɨra a
hunomi, a kuonomi. Të urihi no ɨrapɨonomi. Te he tikëa. Kamiyë pëma kɨ
ha xëprarɨnɨ, pëma kɨ rë waɨwei, ɨra a hunomi. Hu, hu, hu! Ɨra a kãi
kunomi. 

A hapa he rë harenowei, wetinɨ a he harema? Yami ohi pënɨ a resi no
preaama. Xi kɨ pë kõririwë no preaama. Nakɨ kuami yaro, Ɨra. Tukutuku të
nakɨ pehĩto kua yaro, të pë kãi waxikanomi, maromaro ɨ̃hɨ rë të urihi ha,
ɨ̃hɨ Xererei a urihi mɨ amo ha, ɨra yami a huma. Ohiri hurewë. Ai të wai
waimi yaro, ërëkëwë, a ɨ̃kɨma. Ɨra a naikiri ɨ̃kɨma. 

Wetinɨ a he harema? Wapororitawë a kuopë ha, ɨra a warokema, Yanomamɨ ai
a rii përɨopë ha, yami, Wapororitawë ihaƗrariwë e warokema. A pëtarioma,
a mɨ pamarema. A ohiri rë katitore hamɨ. Ɨhɨ hawë hupëpɨ, hawë mamo kɨ
maa hapa a hõra no preaaɨ kuaama, a kraikraipraotima, a rë yutuhouwei
ohiri. 

Waporomɨ kë kɨ! Ɨhɨ wama të pë wãha yuaɨ? Kutaenɨ ɨ̃hɨ Yanomamɨ a kuoma,
Yanomamɨ të pë mamikɨ hõra wai hĩrio ma rë mai! Krihi! Të kãi kuimi, a
ɨmɨɨ makure kiha. Ɨhɨnɨ a hiraa parɨkema. A huɨ hirakema. 

E mamikɨ mahɨkema, ɨra a kraimou maopë. 

Ɨhɨ a ha hirakɨnɨ, a ha ukuuuuhaparunɨ, a ha yarɨɨɨɨɨhɨ taparunɨ, timɨ
parunɨ. 

\chapterspecial{A {onça} e o {tatu}}{}{}
 

\letra{É}{ um tatu!} Dizemos assim. Tatu estava andando. Yanomami, Tatu, tatu.
Hoje, a onça mata e come tatu. Hoje, os dentes do tatu ficam na boca da
onça. A onça tem dentes de tatu. 

Àquela época, os dentes de Tatu saiam da boca, apesar de ele ter boca
pequena. Ele comia coisas grandes. Onça vai tomar emprestados os dentes
de Tatu, por isso, ele hoje tem dentes pequenos. Primeiro, Tatu
emprestou os dentes a Onça e colocou seus dentes na boca de Onça, seus
próprios dentes. 

A onça nos comerá. Ela não vai me comer!? Não tenham dúvidas! 

Onça e Tatu se encontraram, ela ia como gente. Tëɨ! Tëɨ! Tëɨ! Krai!
Krai! Xiri! Hɨ̃krai! Xiri! Krou! Kopou! Poxo! Rae! Os dois faziam o mesmo
barulho. Tatu ficou parado, ela vinha em sua direção. Quando ela o viu, se
aproximou. Tatu olhou para os dentes de Onça. Os dentes de Tatu saíam da
boca. Hɨ̃ɨa! Originalmente, Tatu tinha os dentes que a onça possui hoje. 

--- Irmão menor! Irmão menor! Com esses dentes, você come sem
problema!--- disse Onça. 

--- Como são seus dentes? Você não tem dentes como os meus?

--- Não tenho! Por isso eu não mastigo quase nada. Eu sofro! 

--- Cadê? --- Quando Tatu perguntou, Onça abriu a boca. 

--- Hɨ̃ɨɨ! Como você vai comer? Quer experimentar os meus? Arranque
os seus! 

Os dois conversavam. Os dentes finos de Onça pareciam frouxos e finos
como agulhas na boca de Tatu.

--- Arrancou! Arrancou! Arrancou! Pronto! 

Os últimos dentes do fundo ficaram grudados, Onça deu os dentes para
Tatu. Os dentes de Tatu se tornaram pequenos.

--- Arrancou! Arrancou! Arrancou! --- disse Tatu.

Os dentes do fundo. 

--- Arrancou! Arrancou! Arrancou! 

Para colocar os dentes, Onça abriu grande a boca. 

--- Krɨtɨ! Krɨtɨ!--- Agora você não passará mais fome. Agora pode logo
comer coisas grandes! Você matará animais, você matará anta!--- Tatu
disse. 

Por isso Onça ficou feliz. Ela o abraçou. 

--- Você mastigará ossos e engolirá ossos mastigados--- Tatu disse a
Onça. 

Imediatamente, Tatu  passou a comer somente minhocas; para comer minhocas, ele
cavava a terra. Ele comerá com esses dentes, eles comem assim. 

Será que vai conseguir quebrar os ossos pequenos? Normalmente, não se
quebram coisas grandes, mas Onça conseguirá quebrar coisas grandes.
Foi assim. 

 

\chapterspecial{Ɨ{ra} {xo}, {opo} {xo} kɨ {he} {haap}ɨɨ}{}{}
 

 

Opo kë a! Pëma kɨ kuɨ. Oporiwë a huma. Yanomamɨ, Oporiwë, opo. 

Kamanɨ a ha xëpranɨ, a wapë makui, ɨ̃hɨ opo, ɨ̃hɨ nakɨ ɨra iha nakɨ kua. 

Ɨhɨ hei opo e nakɨ pata rei pramoma, kahikɨ ihirupɨ makui. A ihirupɨo
tëhë, pata të pë wama. Ɨhɨ oponɨ ɨra nakɨ rë kui opo ɨ̃ha e nakɨ
mahɨkema. Nakɨ ma rë oxei. Ɨhɨ ɨra nakɨ mahɨpou, oponɨɨra nakɨ tɨkema,
kama nakɨ. 

Kamiyë pëma kɨ wapë, ɨranɨ. Ware a waimi! Pë puhi kuu mai! 

A mɨ hetua piyërema, opo. Yanomamɨ kurenaha e huimama. Tëɨ! Tëɨ! Tëɨ!
Krai! Krai! Xiri! Hɨ̃krai! Xiri! Krou! Kopou! Poxo! Rae! Ɨra e kua mɨ
heturayoma. Ɨhɨ Oporiwë e rë kui, e yanɨkɨtarioma, a katitoimaɨ ha, a ha
tararɨnɨ, e u kua katitikema. Ɨra nakɨ mɨma. Opo e nakɨ pata reipramoma.
Hɨ̃ɨa! Opo nakɨ hamɨ, ɨra nakɨ. Ɨranɨ nakɨ rë tapore. 

--- Oxei! Oxei! Mihi kahë wa nakɨ rë kuinɨ, wa iaɨ ha ayaowei --- ɨra e
kuma. 

--- Weti naha kë wa nakɨ kuwë? Hei ya nakɨ rë kupenaha wa nakɨ mata
kupowë! 

--- Kuami! Kuwë yaro, ya të pë wai wãxikipraɨ ha maonɨ, ya no preaaɨ! 

--- Weti hamɨ kë? --- a kuɨ ha, ɨra kahikɨ pata reretarioma.

 

--- Hɨ̃ɨɨ, wa të iaɨ ayao ta yaitakë! --- e ha kunɨ --- Pei! Ipa wa të kɨ
wapaɨ puhio? Mihi ëhë tëkɨ ta ukërari! 

Kɨ noã tapɨyoma. Hawë proreprore ɨ̃hɨ ɨra kama nakɨ wai rë ihirupɨ të kɨ
wai rë kui, hawë unamo të pë wai rë xororoi: 

--- Ukë! Ukë! Ukë! A ɨ̃naharë! 

Hei manakoro ha të kɨ wai xatipɨo tahiopë, ɨ̃hɨ ɨra e nakɨ hipëkema, opo
iha nakɨ oxeprarioma. Kamanɨ: 

--- Ukë! Ukë! Ukë! Ukë! 

Manakoro të kɨ pata. 

--- Ukë! Ukë! Ukë! 

Kamanɨ nakɨ rë tɨaɨwei, ɨra e kahikɨ pata reretaoma. 

 

--- Krɨtɨ!, Krɨtɨ! Pei kuikë wa ohii mai kë të! Ɨhɨ hei kuikë rë wa pata
iaɨ xoao. Yaro wa kãi xëprapë, xama wa xëprapë --- a noã tama. 

Kuwë yaro e puhi topraroma. A hẽkato hãore hërɨma. 

--- Wa të ũ pë wai ha waxikanɨ, wa të pë wai waxikano suhapë --- e
kuma. 

Kama oponɨ horema e xi pë xĩro waɨ xoaoma, kama opo a rë kui, horema xi
pë waɨ ha, a titëtitëmou xoakema. Ɨhɨ nakɨnɨ a iapë, pë ma rë iaɨwei. 

Ɨhɨ ihirupɨ rë të ũ wahatoapë? Hɨ̃ɨ! Ɨnaha kuwë të pë pata wahatomamou ma
mai, të ũ pata wahatopraɨ he yatiopë. Ɨnaha a tama. 

 

 

 

\chapterspecial{A {multiplicação} {das} {onças}}{}{}
 

\letra{E}{m seguida,} segue a história daquele que fez as onças se multiplicarem.
Ele foi à direção certa. Existe nos buracos de pau. Onde havia buraco de
pau, outro tipo de onça existia, a onça \emph{ɨrahena}. 

Não foi obra de ninguém! Eles tinham um xapono como este. Era o mesmo nome
daquele que a tirou do buraco. Aquele que tirou a onça \emph{ɨrahena},
onça parecida com jia, depois de tirá"-la, ele se alegrou com a pele
pintada; depois de ele arrancar as folhas, as onças habitaram toda a
floresta. 

Ele chegou ao xapono. Haviam queimado. Nesse lugar, a roça estava
próxima. Ele plantou as jias no lugar queimado. Ele plantou. Apesar de
ser jia, ela não apodreceu, pois era onça. Onde ele plantou um pouquinho,
ao final do dia, quando a floresta escureceu, da mesma forma que os
capins tem flores, essa flor de onça também desabrochou. 

A onça grande começou a surgir. Onde caíram as sementes, as onças se
levantaram. Os Kaxanawëteri moravam no centro dessa região. São aqueles
que plantaram a onça \emph{ɨrahena}. 

Surgiram as onças suçuaranas, as onças suçuaranas vermelhas, as grandes
onças e as onças pretas. As onças exterminaram os habitantes do xapono
onde haviam plantado as onças, e cujo nome eu dei. Ninguém sobreviveu. 

Elas são famintas de carne, e não foi só uma que andou. Logo comeram os
habitantes. Esses antepassados não tiveram descendentes, pois nenhum
conseguiu fugir. Nenhum sobreviveu. Exterminaram todos. Nenhum. Não foi uma onça só. Em um dia, exterminaram todos. Comeram também aquele que tirou
a onça. Ele morreu também. 

Depois de exterminar todos, a onça continuou a surgir na terra
dos \emph{napë}, apesar de essa terra se estender bem ao sul. Não foi obra
de ninguém. As onças apareceram onde foi plantada a onça. Apesar de ser
jia, a jia não apodreceu. Lá, a onça ficava dentro, a
onça \emph{ɨrahena}. 

O que segue é a história das onças que comeram muita gente. 

Eles moravam perto da serra Yamaro e se chamavam Yamarowëteri. Chamaram
o xapono deles Yamaro. Apesar de eles não terem plantado urucuzeiros,
havia muitos no meio, por isso se chamavam assim. É outro nome para
urucuzeiro. As onças os comiam também nas regiões vizinhas. 

Os vizinhos um pouco mais distantes eram os Sementes"-de"-Urucu. Eles
bebiam a água do rio Ximono. A voz deles era fina. 

Após o xapono deles, havia outro grupo. Eram os vizinhos. Todos tinham
os cabelos vermelhos. Os cabelos deles era de um vermelho bem forte. Os
vizinhos deles eram os Ɨranawëteri. Chamaram o rio, do qual bebiam a
água, Ɨrana; por isso se chamavam Ɨranawëteri. Assim faziam nossos
ancestrais.

\chapterspecial{{Ɨra pë rë pararoyonowei}}{}{}
 

 

Ɨharanɨ, hei të rë kui hamɨ, ɨra pë rë pararoyonowei, a përɨoma. Ɨhɨ te
he tikëa. Pë përɨoma. Ɨhɨ ya pë wãha tokumarema. Hei pë përɨoma. Hei pë
rë kuinɨ, a katitirayo hërɨma. Hii hi pëka ha të pë ka ma rë kuprai.
Ɨnaha te hi ka kuopë ha, ɨra hena titioma. 

Taprano mai të ã. Hei kurenaha pë xapono kuoma. Pei a wãha yai. Pei hena
yai rë ukërenowei. A wãha yaia. Ɨhɨ a wãha kohomowë. Hena pa xërema.
Hena ha ukërënɨ, moka kurenaha, të wai ha ukërënɨ, oni sipo wai oni, e ã
topraroma. Të ha ukërënɨ, ɨra pënɨ urihi a haikiprapehe. 

A kõpema. Kihi naha ĩxino praa. Hikari a ahetea yaro. Ĩxino të ha, hei
moka a keaɨ kure. A kekema. Moka a kuoma makui, a kãi tarenomi, ɨra a
yaro. Të wai ha kekɨnɨ, ɨ̃hɨ mahu të rë keare ha, motoka maprou tëhë, të
urihi mɨ titihiprou tëhë, hei porema hi pë rë kurenaha, porema hi pë
hemoxi rë kurenaha, ɨra e hemoxi kuaama. 

Poroporo pë pata kupro hërɨpë. Ɨhɨ hei të pë hemoxi pata rë prerëre
hamɨ, ɨra pë pata hokëko hërɨma. Kama pë yahipɨ rë mɨ amoonowei ɨ̃hɨ pë
wãha Kaxanawëteri kuoma. Ɨhɨ pë yainɨ ɨra hena kekemahe. 

Ɨra ketihenarimɨ, wakëwërimɨ, poroporokohe pë, hũkumari si pë, të pë
pata kuprarioma. Ɨhɨ ei ya pë yahipɨ wãha rë yuprarɨhe, ɨ̃ha a rë keare
ha,ɨranɨ pë haikirarema. Ai pë hëpranomi. 

Pë naiki yaro, mori mahu të rë hure ha. Ɨhɨ teri pë waa xoaremahe. Ɨha
ai të no hekama rë hëprouwei, patama të pë kãi tokunomi. Të pë
hëpranomi. Pë haikiaremahe. Mori mahu të rë hare. Ɨhɨ mahu të mɨ haru ha
të pë haikia xoaremahe. Kamanɨ a rë ukërenowei, a kãi warema. No payeri
taprarema. 

Ɨhɨ të pë ha haikirahenɨ, ɨ̃hɨ tëhë, hena rë kekɨhe tëhë, ɨra pënɨ napë a
urihi makui, napëpë urihipɨ hamɨ koro hamɨ makui, ɨra a kuprario hërɨma.
Taprano pë mai! Keano pë hamɨ ɨra pë kuprarioma. Moka a makui, moka a
kãi wãrimonomi. Ɨhɨ ɨra a titioma. Ɨra hena. 

Omawë pë kuprarioma, yai të pë. Omawë pë wãha kekema. Omawë yai hena
paxërema. Ɨnaha të pë kuaama. Ya të pë rë hĩrinowei, ya të ã taɨ. Ya toa
hërɨma. Ɨhɨ të waikiwë. Hiriano, wëyëno, pata të pënɨ wamare kɨ noã
tamahe. Ya prao tëhë, ya ha praonɨ, ya të hĩrima. Ɨnaha të kuwë. 

Ɨhɨ te he tikë hamɨ, pë pruka waɨ he rë tikëkonowei,ɨ̃hɨ të kãi tikëa. 

Ɨra henanɨ pë pruka rë wanowei, pë wãha. Yamaro kɨ ha të pë rë
përɨonowei, kama pë wãha Yamarowëteri kuoma. Xapono e rë kuonowehei
Yamaro awãha yupomahe. Nara pruka xi hi pë keanomi makui, ɨ̃hɨ të xihi pë
pata mɨ amo ha pë kua yaro, pë wãha kuoma. Nara xi hi pë wãha Yamaro
kua. Ɨhɨ pë pruka wama. Ɨha pë waɨ he tikëkoma. 

Ɨhɨ te he ĩsitoripɨ tikëre ha, pë pruka yahipɨ he rë tikëkëmonowei,
Ximonowëteri pë hiraoma. Ximonowëteri pë rë hiraonowei, kama Ximono u
koamahe. Pë kahi ã kãi preonomi. 

Ximonowëteri pë yahipɨ he tikëo ha, ai a yahi përɨoma. Ɨnaha të pë
henakɨ wakë kõre kumou xĩrooma. Hei pëma kɨ henakɨ rë kurenaha, pë
henakɨ kuoma? Pë henakɨ wakë kõremoma. Ɨhɨ ei Ximonowëteri pë yahipɨ he
tikë ha, kama Ɨranawëteri pë yahipɨ he tikëoma, Ɨranawëteri. Kutaenɨ
kama pënɨ u rë koanowehei, Ɨrana u wãha yuamahe. Ɨnaha no patama të pë
kuaama. 

 

\chapterspecial{{Minhocão}}{}{}
 

\letra{A}{ história das minhocas.} Quando a floresta existia, mesmo que a terra
existia: 
--- Vou cavar minhocas! --- ninguém dizia isso. 
Não existia minhoca e, como as minhocas não saíam, ninguém saía, ninguém
pescava depois de tirar minhocas. Era assim. Nós não as deixaremos cair
na água, quando estamos com fome, nós cavamos onde há minhocas, nós as
tiramos, muitas surgirão; para que nós fizéssemos assim, ele morou com a
menina. Lá onde surgiu aquela mulher, a filha de Pokoraritawë ensinará
as Yanomami a não gostar do marido;  às vezes elas não gostam dos maridos. Ensinando"-nos, a filha de Pokoraritawë se zangava demais, pois
estava com medo, não queria seu marido. Apesar de ele ser muito bonito,
a mulher não o queria, a filha de Pokoraritawë fez as minhocas surgirem.
A mulher chegou lá com os dois Minhocões, que comiam o esperma deles
mesmos. Aquele que ela desposou, apesar de ser bonito, foi embora caçar,
até que afinal a mãe falou com a filha: 

--- Filhinha querida, teu marido foi de novo! Vai atrás dele! Vai!---
ela disse. 

Ela foi bem devagarzinho atrás dele. 

Ele foi, soprou veneno em cuxiús, matou; ele era muito bom caçador,
Paricá. Ela não gostava dele, de Paricá, era o nome do genro de
Pokoraritawë. Minhocão fez os filhotes se multiplicarem com a esposa de
Paricá. Quando seu marido passou, os dois chegaram aonde Paricá estava.
Ele estava longe, adiante, quando a mulher passou perto dos dois
Minhocões, o mais velho e o mais novo. 

Os dois moravam na terra plana e viviam na condição de Yanomami, pois
não existiam minhocas à época. Os pais das minhocas moravam lá, no
início. Eles farão os filhos se multiplicarem. Passando nesse caminho, lá
em baixo, bem longe, Paricá matava cuxiús. As frutas de Minhocão estavam
grudadas. Naquele caminho, as frutas eram numerosas, para atrair a
mulher. Toso, toso, toso, toso! Faziam os restos. Hõti, hõti, hõti!
Faziam assim também. 

Os dois eram muito bonitos, os pais das minhocas: tinham a testa
enfeitada de rabo de cuxiú, guardando a testa, o rosto dos dois era
enfeitado e bonito. Assim era o rosto dos dois. Os dois Minhocões tinham
barba bonita, para parecer o rosto de Paricá e enganar a mulher. Ela
olhou: 

--- Krai! Rae!--- disseram assim. 

Os dois eram esbranquiçados: 

--- Hɨ̃ɨ! Olhe! Olhe! É você?--- disse a mulher bem bonita, com seios
bonitos.

--- Ô! De quem é essa voz?

Como tinha uma clareira, a mulher ficou em pé no limpo. 

--- Não pergunte quem sou! Sou eu! Você! É você mesmo! --- disse a
mulher.

--- Não, não sou aquele que você pensa, eu sou outro! 

--- É você, é seu rosto mesmo, assim que é o seu rosto! 

Ele pronunciou seu nome: 

--- Eu sou mesmo o Minhocão!

--- Não, você não é outro, é você! 

Enquanto ela insistiu em dizer isso, os dois Minhocões logo contaram a ela quem eram. 

Um deles olhou e disse: 

--- Se você diz assim, tire essa folha nova de arumã, aí, aquela folha
enrolada, você a arranca e a desenrola, e você senta em cima, sente"-se
em cima. Coloca sua bunda em cima--- disseram os dois, de um jeito
cantado. 

Rindo, ela correu para arrancar a folha. Pensando que era Paricá, pois
tinha o mesmo rosto, quando ele disse isso, ela arrancou a folha. Depois
de arrancá"-la e desenrolá"-la em um lugar bonito da clareira, onde não
havia nada, ela se sentou em cima, onde estava limpo. Os dois desceram,
os dois desceram rapidamente e copularam com ela uma vez, não várias
vezes, somente uma vez. Apesar de copular  com ela somente uma vez cada um, os dois copulavam enquanto o marido estava matando todos os cuxiús,
pois era muito bom caçador, acumulando as presas. 

Ela não o alcançou, andava devagar. 

Depois de ter copulado, não foi nos dias seguintes, mas no mesmo dia,
apareceu a barriga que, apesar de uma vez só, já estava crescendo. 

--- Vai! Vai logo! --- disseram os dois Minhocões, que voltaram para a
morada deles. 

O ventre daquela que estava andando sozinha crescia e crescia. 

--- Vai lá, onde teu marido está matando os cuxiús, ouça os gritos!---
disse o Minhocão. 

--- Hõhaaa!--- ela ficou pensando.

Depois de falar isso, ela foi bem devagar à direção onde estava seu
marido. Indo lá, o ventre sempre crescia, porque não havia só um filho.
Apesar de serem pequenos, eles estavam acabando com a carne dela. Ela
ficou em pé, enquanto Paricá estava amarrando os cuxiús, ela ficou em pé
lá longe. 

Ele estava voltando. Ele havia matado todos os cuxiús e estava voltando,
depois de carregá"-los, ele estava voltando. Quando voltava, ele viu o
ventre dela enorme de gravidez. 

--- Nunca mexi nessa mulher, e tem filho nesse ventre!--- ele pensou. 

Ele simplesmente pensou. Ele nunca tinha copulado com ela, pois ela não
gostava dele. Ele passou, voltando. Ela voltou sozinha. Ela estava
voltando rindo. Ela estava voltando atrás, sua barriga cresceu
rapidamente. Ela voltava com esse ventre enorme. 

Depois de um dia, o ventre dela estava gigantesco. Ele olhou atrás e viu
a mulher com a barriga enorme. 

--- Hõãaa! É barriga com criança--- ele pensou, e continuou andando. 

--- Hɨ̃ɨɨ! Será que eu já a sujei?

Xiri! Anoiteceu muito rápido. A noite caiu depressa. O ventre estava
cheio. Olha só o suporte dos bichos. Não havia só um! O ventre estava se
mexendo. 

--- Õa, õa, õa, õa!--- diziam, lá dentro. 

A mulher sofria, sofria passando mal, sofria por causa do que acontecia
dentro dela. Doía muito o ventre dela. O marido dela estava deitado na
sua rede, sem olhar para ela, enquanto o ventre dela doía, pois doía
muito, acariciando sua barba e, enquanto a noite logo ficou densa e
grossa, as minhocas saíram.

Weo! Weo! A placenta se derramava como se fosse água, e saíam filhotes de
minhocão:

--- Ũa! Ũa! Ũa!--- já faziam assim. 

Como parecia voz de criança, ele olhou para as crianças no chão, apesar
de estar deitado na rede, ele olhou. Não havia criança. Ele olhou de
soslaio. Não dava para ver. Embaixo dele:

--- Ũa, ũa, ũa, ũa!--- faziam sem parar. 

Eles nasciam, nasciam, nasciam, nasciam, nasciam, nasciam, nasciam. 

Hɨ̃ɨɨ! Havia tantos montes de minhoca que o fundo da casa sumiu, a vagina
dela estava cheia de minhocas. Depois do trabalho de parto, ela
olhou; ela fez assim. Eles choravam como crianças, chorando de sede
já, eles demonstravam sede: 

--- Sede! Sede!--- diziam, com uma voz de criança --- Estou com sede! ---
diziam rapidamente. 

--- A criança cresceu tão rápido! --- ela pensou assim. 

Como estavam sempre com sede, ela deu o seio. 

--- Tusu! Suku! Tusu! Suku! --- faziam assim enquanto mamavam. Ela
fez assim. Como as minhocas faziam isso, ele ficou esperto. Ele
entendeu: 

--- Hɨ̃ɨ!--- ele pensou. 

A mãe dela chegou correndo. Apesar de olhar, ela não as viu
imediatamente. Apesar de escutar o choro de criança, ela olhou e voltou
a deitar. 

Deitada, a mãe das minhocas as cobriu, cobriu, cobriu, cobriu, cobriu.
Amanheceu. Como a filha estava indo de manhã cedo, ela falou para sua
mãe, enquanto o marido estava ali pensativo.

--- Mãe! Não descubra o que eu cobri no fundo da casa. Não fique olhando
o fundo da minha casa! 

Havia tantas minhocas! Elas se embolavam, zoando, porque estava cheio. 

--- Não olhe o fundo da minha casa. Não descubra o que eu cobri!--- ela
disse, e saiu. 

Xiriririri! E sumiu. Enquanto isso, a mãe levantou da rede. 

--- Por que? Onde está essa criança, que deveria estar no colo,
recém"-nascida? Vai chorar muito, assim!--- ela pensou, e correu até a
casa. 

Ela foi logo. Ela correu e descobriu o que estava onde a filha morava,
aquelas minhocas, todas mexiam a cabeça ao mesmo tempo. 

--- Xiririririri! Sede! Sede! Sede! Avó! Sede!--- eles a chamavam de
avó.--- Avó! Sede! Avó! Sede! Avó! Sede!--- todos diziam. 

--- Hɨ̃ãaaaaë!--- ela gritou logo--- Hɨ̃ãaaaë! Só você para fazer surgir
aquilo! Por isso! Você não trata bem seu marido! É por causa desses
bichos estranhos que você não conseguiu dormir!--- ela disse--- Vai! Meu
genro! Enquanto eles se mexem assim, derruba logo essa lenha, faz um fogo
grande para ela!--- disse a mãe. 

Ela mandou queimar a filha viva! Depois de ela dizer isso, ele desceu da
rede. Ele não demorou: derrubou aquele carapanã-uba.

Kraxi! Kraxi! Kraxi! Krao! Torou! Fazia lenha para cremá"-la. Enquanto
fazia lenha, ela voltou. Ela tinha ido tomar banho sem perceber, ela passou
onde ele estava partindo a lenha. Ele virou as costas, onde ele estava fazendo
lenha. Ele nem olhou. Ela se deitou, encolhida. 

Pou! Pou! Pou! Ele amontoou muita lenha. Pou! Pou! Pou! Ele pegou
brasas para acender o montão de lenha, ele fez aumentar o fogo. Como a
lenha era seca, o fogo pegou logo. 

Weee! Ele fez uma cerca, fez para ela. Depois, ele correu atrás dela.
Ela nem se levantou, ele gritou para pegá"-la, pois queria a cremar viva. 

Weeeee! Ela estava deitada bem reta. Ela nem reagiu, ele correu a carregando
em direção do fogo, e ela chorava: 

--- Ëaë! Ëaë! Mãe! Pai! 

As pernas dela estavam balançando, dando impulso. Ele a jogou no meio
do fogo. 

Pou! Ele pegou outra lenha que estava no chão e amassou, amassou com
força. 

Ëëëaaaëëë! Proto! O fogo queimou, enquanto cremava, a sogra dele correu
em cima dos minhocões para queimar os feios. Ela correu para pegá"-los.
Ela já tinha colocado água em cima do fogo em uma panela de barro para
cozinhá"-los. Ela correu com uma vasilha de água quente em direção das
minhocas cobertas. Ela jogou a palha de coruá que as cobria: 

Weeeo! Os minhocões gritavam:

--- Õiii, õiii, õiii! 

--- Avó! Couro encolhido! Couro encolhido! Couro encolhido! Couro
encolhido!--- diziam atordoados, chamando"-se de pele encolhida. 

--- Avó! Couro encolhido! Avó! Couro encolhido! Avó! Couro encolhido!---
diziam os pedaços, arrebentados. 

Olha só os montões de pedaços! Os pedaços estavam correndo logo, e
ocuparam toda a floresta, os minhocões. Ficaram ocupando a floresta, os
arrebentados, correndo logo pra todas as beiras de rio, entraram depressa
no fundo da terra. 

Depois de acontecer isso: --- São minhocões! --- Dizemos. Foi assim que aconteceu.
Não existiam minhocas. Foi com ela que se multiplicaram. Nós as faremos
cair na água para nós comermos peixes. A minhoca não apareceu do nada. 

Foi depois de os dois Minhocões copularem com ela e multiplicarem seus
filhotes, que foram embora com os pais. Os filhotes não moraram onde
foram cremados, nem ficaram ali perto. Os dois foram logo. Assim foi. Desde
que aconteceu, quando cai a chuva:

--- Tëɨ, tëɨ, tëɨ, tëɨ tëɨ!--- Dizem seus pais, de onde estão. 

Assim foi a história.
\chapterspecial{{Horemariwë}}{}{}
 

 

Ai të ã. Horema të ã. Horema pë rë kui, urihi a kuo tëhë, heinaha xomi
(pita) a kuprou tëhë, a kuo tëhë:

--- Kiha horema ya kɨ tiëaɨ! --- ai të pë kunomi. 

Horema xi pë kãi kuonomi, horema kɨ ha harɨnɨ, ai të hunomi, ai të kɨ
wai ha tiërehenɨ, yuri a kãi rëkaɨ taonomihe, ɨ̃naha të kuoma. Ɨhɨ pëma
të pë keamapë, kamiyë pëma kɨ ohii tëhë, të pë kuopë ha, pëma të pë ha
tiënɨ, pëma të pë pata ukapë, të pë pata kupropë. Tëëpɨ kãi përɨkema.
Ɨhɨ ihamɨ a suwë rë pëtore hamɨ, Pokoraritawë tëëpɨ ihamɨ Yanomamɨ të pë
ɨpɨo hiraɨ ha, të pë ma rë ɨpɨowei; ɨ̃hɨ të hiraɨ ha kamiyë Yanomamɨ pëma
kɨ iha Pokoraritawë tëëpɨ huxuo he parohooma, a kirii yaro, a puhinomi,
a riëhëwë totihiwë makui, suwënɨ a puhinomi, Pokoraritawë tëëpënɨ horema
pë rarakema. Ɨhɨ Horemaritawë kama kɨpɨ iha a suwë ha waroikunɨ, kama
moukɨ kete waoma. A rë ponowei, katehe a makui, a rë ponowei, e hurayo
hërɨma, yakumɨ ɨ̃hɨ ihamɨ pë nɨɨ e ã haɨ heama: 

--- Xõe! Hẽarohë a nohi hua kõrɨhe! A wai huto hërɨɨ! Hurɨ hërɨ! --- e
kuma. 

Opi e hua hërayo hërɨma. E ha hurɨnɨ, wɨxa e kɨ horama, e kɨ niama, a
nɨhɨteo he parohoma, Yakuana a nɨhɨteo he parohoma, ɨ̃hɨ iha a ɨpɨoma,
Yakuana iha, Pokoraritawë siohapɨ wãha, Yakuana hesiopɨ iha
Horemaritawënɨ ihirupɨ pë rarakema. Hẽaropɨ e ha hayuikunɨ, e kɨpɨ nosi
ha wetitarunɨ, a ha kuuuuupohorunɨ, a suwë hayuo ahetou tëhë. 

Horemari kɨpɨ, pata, oxe, kɨpɨ rë përɨpɨa yarɨtaawei ha, Yanomamɨ kɨpɨ
rë kuonowei, horema pë kuami yaro. Ɨhɨ hapa horema hɨɨpɨ kɨpɨ përɨoma.
Ihiru pë raraapë. Hei yo hayua, hei, e ha kuuuu katiiiii tipokɨrɨnɨ,
Yakuananɨ wɨxa kɨ niama. Kama moukɨ të pë pata yëtëpramoma. Hei yo ma rë
kui, të pë pata ximokorepramoma, suwë a rurupëapë. Toso, toso, toso,
toso! Kɨ kanosi kupɨma. Hõti, hõti, hõti! E kupɨma. 

Kɨpɨ riëhëo totihioma, horema pë hɨɨpɨ rë kui, kɨpɨ mohekɨ wëhuhuoma,
hei wɨxa texina si pë rë kui, texina si yohopɨpoma. Kɨpɨ mohekɨ
wëhuhupɨwë totihitaoma. Ɨnaha rë e kɨpɨ mohekɨ kupɨoma, kɨpɨ kaweikɨ kãi
totihitapɨoma. Yakuana mohekɨ kurenaha, suwë a mɨramapɨpë, a mamo
xatitarioma: 

--- Krai! Rae! --- të kutario ha. 

Kɨpɨ pruxixiwë: 

--- Hɨ̃ɨ! Mɨpraa! Mɨpraa! Mɨpraa! Ɨhɨ kahë rë wa? --- e kuɨ pëtarioma,
suwë. E kupɨɨ pëtarioma, xëkëkëwë, suhe puu wai totihitaoma, no xi
aihawë. 

--- Õ! Weti wãwã ta tawë, weti pei wãwã ta tawë?

Heinaha të ka yakëa kua yaro suwë a wawëtowë upratarioma. 

--- Weti mai! Kamiyë kë ya!

--- Kahë rë wa, kahë rë wa nohi kui!

--- Ma, kamiyë ɨ̃hɨ ya tama! Kamiyë yaiwa ya rii.

--- Ɨhɨ rë wa, ɨ̃hɨ wa mohekɨ katitire! Ɨnaha wa mohekɨ kuwë! --- e kuma.
A wãrima. 

A wãha yuprarioma: 

--- Kamiyë Horemari ya rii ta kui!

--- Ma! Wa no yaipimi, ɨ̃hɨ rë wa! 

Ɨnaha e kupɨɨ ka kuaaɨ ha, e kɨpɨ ã hapɨɨ xoaoma. 

E mamo xatitarioma. 

--- Ɨnaha wa kuu kunoi, mihi umoromɨ henakɨ tuku rë tɨrɨrɨre, mihi hena
rë hututure, wa hena kɨpɨ ha ukërënɨ, wa hena ha hapexeprarɨnɨ, ɨ̃hɨ hena
ha wa rokei, wa rokei kë tao! Wa koro pakohekei kë tao --- e kɨ kahiã
kupɨma.

 

E ĩka wã kãi rërëa nokakema. Hawë ɨ̃hɨ a kuwë ha, ɨ̃hɨ Yakuana mohekɨ
kuopë naha, mohekɨ kuo katitioma, a kuɨ ha, e hena kɨ ukërema, e hena kɨ
ha ukëpɨrënɨ, hena kɨ mɨ ha hapexeprarɨnɨ, heinaha të totihitaopë ha, të
ka yakëopë ha, a koro pakohekema, të tãihiopë ha. Kuaaɨ tëhë, e kɨ
itopɨrayoma, a napë itopɨa haɨtarayoma, e kɨ rë itopɨre, na wapɨma,
mahu, hei ai na waɨ, na waɨ, na waɨ të kupronomi, mahu! Ainɨ mahu na
waararei, ainɨ mahu na waararei, ɨnaha mahu makui. Hei na rë wapɨɨ hëre,
kiha hẽaropɨnɨ wɨxa kɨ haikiaɨ kë tëhë, a nɨhɨteo he ha parohoonɨ, kɨ
weyoyamatii kë tëhë. 

Ɨha e waroo mai! 

E opisi huɨ hëo hërɨma, hei na rë wapɨararɨhe ha, ai të henaha e makasi
kãi pëtonomi, mahu makui, ihiru makasi tirehetou waikirayoma, hei a rë
waikare ha. 

--- Pei! Wa hurayou kë tao! --- Kama kɨpɨ përɨopë ha, kɨ kõpɨkema. 

Hei a rë huɨ hëoimati hamɨ, ihiru makasi pataɨ waikio hërɨma. 

--- Hẽarohënɨ kiha wɨxa kɨ hõra rë niayahi ha, ɨ̃hɨ ei rë e kɨ rãawa, ɨ̃ha
wa e waroyei --- e kuma. 

--- Hõaaa! --- e puhi kuɨ hëoma. 

Ha kunɨ, opisi e katitiatarou hëo hërɨma. E rë katitore hamɨ, ihiru a
makasi tirehou waikio hërɨma. Mahu të kuami yaro. Pei të pë ma oxei
makui, yãhi kɨ haikiama. E uprakema, wɨxa pë nanoka hãomaɨ tëhë, e upraa
hëwëpetayoma. 

Kama a kõoimama, kɨ niaa waikirarema yaro, a kõoimama. Pë ha yehirënɨ, a
kõoimama. Kõoimanɨ, suwë a makasi karereoma. Makasi kario tirewë
waikiwë:

--- Kihi rë ya pë napë kuaaɨ taoma mai, kihi ihiru rë makasi ë! --- a
puhi kutarioma. 

Kama a puhi kuɨ pëoma. Kamanɨ na wanomi, e ɨpɨo yaro, a kõo e hayukema.
Yami e kõo hëoimama. E ĩka wa teteo hëoimama. Ɨhɨ pei noha hamɨ e kõo ha
hëoimanɨ, e makasi pataa haɨrayoma. Pataɨ hëoimama. 

Ɨhɨ mahu të mɨ haru ha, makasi pata ihea hërɨma. A mɨ yapatou kõrayoma,
suwë makasi pata kareroma. 

--- Hõaaa! Ihiru rë pesi! --- a puhi kutario hërɨma 

--- Hɨ̃ɨ! Ya no kiriaɨ tao ta yaitakë? --- a ha kunɨ. 

Xiri! Ɨhɨ tëhë e të mɨ titihiprou haɨtarayoma. Rope të mɨ
titihiprarioma. Ihiru të makasi pata ihewë. Hei të pë pesi pata hei!
Mahu kë pë kua yaro! Mahu! Të pë pesi pata uprauprapraroma. 

--- Õa, õa, õa, õa! --- të pë pata kuɨ huxomioma. 

Suwë a no preaama, xi kirihiwë no preaama, huxomi xi kirihiwë no
preaama. E të pë pesi pata niniprarioma. Hẽaropɨ e kutaoma. Ɨhɨ e
hesikakɨ rë rëprapohorohe, ɨ̃hɨ të pë ã pata ma përao tëhë, të makasi
hõra niniaɨ ha haɨtaikunɨ, kuaaɨ ha, huxo huwëtaoma, ɨ̃hɨ tëhë të mɨ titi
supraa hërɨɨ tëhë, heẽteprou hërɨɨ tëhë, të pë pata keama. 

Weo! Weo! Hawë mau u pata rɨpraama, ɨ̃hɨ horemari pë ihirupɨ yono u
pata. 

--- Ũa! Ũa! Ũa! --- hapa e të pata kurayoma. 

Ihiru a yai kuɨ makure, ihiru e rĩya praa ha mɨnɨ, e hesikakɨ ma rëre të
mɨma, kuonomi. Mamo axëoma. Taproimi. Ɨha kama pepi hamɨ: 

--- Ũa! Ũa! Ũa! --- të pë pata kutima. 

Të pë pata keaɨ, keaɨ, keaɨ, keaɨ, keaɨ, keaɨ, keaɨ. 

Hɨ̃ɨɨ! Pei kë të pë he pata poraraprawë xĩka maprarioma, na ka no
nihioma, a nohi kuaama. A ha kuprarunɨ, të pë mɨma. Ɨhɨ të pë pata
ũaũapraroma: 

--- Amixi! --- të pë pata kuɨ haɨtaoma. Pë amixi himou ha --- Amixi!
Amixi! --- të pë pata kuma.

--- Ihiru a wã kãi, wai hõra pataa ropaharayou! --- a kuɨ no mɨhɨtaoma.

--- Tusu! Suku! Tusu! Suku! --- Pë amixi kõo ha, të pë pata tamama.
Ɨnaha të pë pata kuɨ ha, a puhi moyawërayoma. No ɨhɨpɨrema. 

--- Hɨ̃ɨ! --- a puhi kutarioma. 

Ɨhɨ tëhë pë nɨɨ e kua yaro, pë nɨɨ e rërëkema. Të mɨɨ makui, të pë pata
tapraɨ haɨonomi. Ihiru rë a wã makure. Të mɨakema, a përɨa kõkema. 

E ha përɨikunɨ, të pë he pata yohoapotayoma. Yohoaɨ, yohoaɨ, yohoaɨ,
yohoaɨ, të mɨ harurayoma. Harika totihiwë e huɨ yaro, pë nɨɨ iha e ã
hama. Ɨhɨ hesikakɨ kuprao xoao tëhë:

--- Nape! Hei pei ya xĩka rë kui hamɨ të pë he rë yohohore të pë he
karoaɨ heaɨ mai! Ware xĩka mɨɨ heaɨ mai! --- e ku hërɨma. 

Horema pë kua yaro! Të pë mɨ pata puruwë yaro. Të pë pata xiririmoma. 

--- Nape, ware xĩka mɨɨ mai! Hei të pë he rë yohohore të pë he karoaɨ
mai! --- e kuma. E ku hërɨma. E harayo hërɨma. 

Xiriririri! E mato hërɨɨ tëhë, e kuɨ tëhë, pë nɨɨ e waheprarioma. 

--- Exi të ha, exi të ha ihiru weti hamɨ a rë yakapore? Të mɨa hore ma
përamapou --- e puhi ha kunɨ, e rërëkema. 

A rë rërëore, kama a kuopë hamɨ të pë he pata rë yohoawei, të pë he pata
karoprarema, horema të pë pata yai rë prei, ɨ̃hɨ naxomi xĩro, të pë he
pata kuakuaa nokararioma:

--- Xiriririri! Amixi! Amixi! Amixi! Yape! Amixi! --- pë nɨɨ iha të pë
pata yesimoma --- Yape! Amixi! Yape! Amixi! Yape! Amixi! --- të pë pata
pruka kuma. 

--- Hɨ̃ãaaaaë! --- a rarɨa xoarayoma --- Hɨ̃ãaaaë! Ɨnaha wa të pë pata
hore taamaɨ ayao yaro wa të ã hore no kirio ayao nosië! Ɨnaha kë wa të
pë pata xami hore taamatii ayao yaro, wa të ã hore yahoomi ayao no
kuhaë! --- e kuma. --- Pei! Xõe! Ɨnaha pë taamayou tëhë mihi wahë ãxo ha
tuprakɨnɨ, a wakë kata tapɨpa! --- pë nɨɨ e kuma. 

Pë nɨɨnɨ a yaamaɨ puhima. Temɨtemɨ! Kuɨ tëhë e wahetarioma. A no teteheo
mai! Hatoa a rë upraawei. 

Kraxi! Kraxi! Kraxi! Krao! Torou! tarɨkɨ ta ma! Tarɨkɨ tama. Tarɨkɨ taɨ
tëhë, e kõpema. Kihamɨ a rë yarɨmou xomi arui, tarɨkɨ poapë hamɨ e
hayukema. E yaipë rëtakema, kama tarɨkɨ makui ha, e mamo kãi xationomi.
E tasikɨ yoretakema. 

Pou! Pou! Pou! E të pë heru pata taketayoma. Kaɨ wakë ha. Pou! Pou! Pou!
wakë pata ukëa ha piyërënɨ, wakë paramama. Të pata haxitiwë yaro, e wakë
kãi pata waa haɨtakema, arana kɨ tapema, a pehi tapema. A pehi ha
poxokopanɨ, a napë rërëa xoakema:

Weeeee! E kãi hokëpronomi, napë ɨramorayoma, a temɨ yaaɨ yaro. 

Weeeee! E katɨpraoma. Weeeee! E kãi hukëonomi, kaɨ wakë pata hamɨ e kãi
rërëkei ha, e ëaëmorayoma:

--- Ëaë! Ëaë! Napemi! Hapemi! --- e kurayoma. 

E matasikɨ yoayoamoma, e kaxëaɨ ha. Ɨhɨ mɨ amo të wakë pata yai ha: 

Pou! A xëyëkema. Ai ãxo pata rë praawei ãxo pata ha hurihirënɨ, a
patëtëpema. A hɨ̃kɨpema. 

Ëëëaaaëëë! Proto! Kaɨ wakënɨ, a ĩxirayoma, a ĩxii tëhë, pë yesi e
yëkema. Wãriti të pë pata yaprapë. E të pë napë pata rërëa paxikema. Mau
u pata tupoma, hapoka a ha, të pë hete pata rë tuaɨwei. Kaɨ hesi ha e u
kãi pata rërëkema. Të pë he pata rë yohoawei ha. Masiko kɨ pata
maiprarema. 

Weeeeo! Të pë ã pata pëprarioma. 

--- Õiii, õiii, õiiii! --- të pë pata kuɨ pëprarioma.

--- Yape si ãyiki, yape si ãyiki, yape si ãyiki, yape si ãyiki! --- të
pë pata porepɨ kuma. Kama pë si pata ãyikiwë himou ha: 

--- Yape si ãyiki! Yape si ãyiki! Yape si ãyiki! --- të pë pata kuma,
hemata. 

Kihi kë të pë pata hemorokowë yapuruprawë, ɨ̃ha të pë hemata pata rë
rërëoprou xoare, kihi a urihi rë kui, a haikiremahe, horema pënɨ. A urihi haikire hërɨmahe, ɨ̃hɨ rë të pë pata rë hëtɨtɨraruhe, pata u kɨ rë
kutarenaha, hei a pita huxomi hamɨ, të pë pata rërëokema. Rërëo
xoaokema. Ɨnaha pë ha kuoikunɨ: 

Horema kë pë! Pëma kɨ kuɨ, ɨnaha të kuprarioma. Horema pë kuonomi. Ɨhɨ
ihamɨ të pë pata rarokema. Pëma të pë keamapë, yuri pëma pë wapë, kama
horema a xomi kupronomi. 

A ihiru ha rarapɨkɨnɨ, pë hɨɨ ɨha e kɨpɨ rë kupɨonowei, kɨ hupɨrayoma.
Kɨ ihirupɨ yapraɨ aheteopë ha, kɨpɨ kupɨonomi, kɨpɨ përɨpɨonomi, kɨ
ahetepɨonomi. Kɨpɨ hupɨa xoarayoma. Ɨnaha të kuprarioma. Të ha
kuprarunɨ: 

--- Tɨ, tɨ, tɨ! tɨ, tɨ! --- maa a ha kenɨ, ɨ̃hɨ pë kupramopë hamɨ, pruka
pë hɨɨ pë kuɨ. 

Ɨnaha të ã kutaoma.

 

\chapterspecial{O {pássaro} {popomari}}{}{}
 

\letra{H}{á a história para nós,} Yanomami, nos perdermos na mata, ensinou"-se,
ensinou-se a nos perdermos. Eles tinham uma grande roça, e assim faremos.
Aprendemos a nos perder até na roça, de tão grande. Ela se estendia, apesar
de ser roça, e aquele um se perdeu. Tem essa história, também. Foi assim. 

Há os que existiram no início e que se transformaram, aqueles que
existiram; a imagem daquele que gritou existe também na terra
dos \emph{napë}. Ele se perdeu, aquele que se perdeu lá, o eco da sua
voz voou em todas as partes. Ele gritou, o pássaro \emph{popomari} fez
ele se perder. Ela habitou toda a floresta, a voz: Po! Po! Po! Po! Po!
Po! Po! daquele \emph{popomari}, Popomaritawë se perdeu. Aquele que se
perdeu errou de caminho e sua imagem foi embora. Com o eco da sua voz,
para toda a floresta se encher de \emph{popomari}. Nós faremos assim,
pois nos ensinaram. 

Nós também ficamos à deriva em cima dos rios, não voltamos direto. Você se
perderá no rio. Ficamos agindo assim, ele sumiu. Ele gritava, gritava e
ninguém respondia. Não responderam. Ele se perdeu lá longe, no meio da
roça e não responderam. Assim fez, sofreu, por isso, o canto dele se
escuta também na terra dos \emph{napë.}
\chapterspecial{{Popomaritawë}}{}{}
 

 

 

 

Yanomamɨ pëma kɨ mohorupë të kãi kua. Të rë hiranowei. Mohoruu rë
hiranowei. Hikari pata a prapoma. Ɨnaha pëma të pë tapë. Të pata ha
praukurarɨnɨ, hikari a makui ha, a mohoru rë kukenowei, të kãi kua.
Taprano të hamɨ mai! Ai të pë no patama hamɨ mai! Ɨnaha të kuprarioma. 

Hapa, të pë rë kuonowei, ɨ̃hɨ të pë rë kuprarionowei, ɨ̃hɨ napë pë urihipɨ
hamɨ të pë ã no uhutipɨ kãi kuprawë, a rë komɨnowei, a tokurayoma,
kihamɨ a rë tokure hamɨ, a wã no uhutipɨ yëo xoaomopotayoma. A komɨɨ,
popomari pë rë kui, Popomaritawë, a mohorumarema. Ɨhɨ a rë kui: Po, po,
po, po, po, po, po! Të pë rë kutouwei, urihi a haikirema. Ɨhɨ mohoruno a
rii rë yakërayonowei, a no uhutipɨ huokema, ɨ̃hɨ a wã no uhutipɨ; ɨnaha a
urihi no popomaripɨ kuprou haikiopë. Ɨnaha pëma kɨ kuaapë, pëma kɨ
hirama. 

Mau u ha pëma kɨ kãi karëɨ, pëma kɨ kãi kõo katitiomi. Mau u hamɨ wa
mohorurayou. Ɨnaha pëma kɨ kuaaɨ rë hëre, a marayoma. A komɨprarotima, a
komɨpraroma, e të pë ã huonomi. A wã huanomihe. Hikari mɨ amo të pata
hamɨ, a ma mohorurati, a wã hĩrianomihe. Ɨnaha a kuaama, a no preaama,
kutaenɨ, napë a urihi hamɨ të pë ã kãi kuwë. Ɨnaha të pë kuɨ haikiwë.

 

\chapterspecial{O {surgimento} {da} {flecha}}{}{}
 

\letra{A}{ história} da flecha. Aconteceu o seguinte. Tinha o dono. Não foi outro
que depois de abrir um tipo de roça plantou as flechas. Onde morava o
dono, parecia um flechal, essas flechas que eles plantaram em seguida em
todos os xapono\emph{.}

Assim que é, porque ele é o dono mesmo. Aquele que descobriu a flecha se
chamava Xororiakapëwë, é seu flechal, fará atirar as flechas, aquele que
descobriu as flechas, era a imagem das pequenas andorinhas que voam
acima da água. 

Xororiakapëwë descobriu as flechas, fez as flechas \emph{hauya}. Graças
a ele, os Yanomami descobriram a flecha e pegaram"-na. O limite do
flechal fica na boca do rio subindo, é seu flechal, não é de Yanomami.
Eles pegaram as flechas e as espalharam. Ele fez as flechas se
multiplicarem. 

Os Yanomami não tinham flechas, depois de pegarem"-nas e plantarem"-nas,
eles guerrearam. Antes, eram desprovidos, não tinham flechas, eles
flechavam com dalas pequenas de arumãs em penas, aquelas flechas nativas,
ou de caule de planta \emph{tomɨ si}. Ofereciam"-se essas flechas de má
qualidade, pegavam haste de caranarana parecidas com flechas, amarravam
penas na extremidade e flechavam com essas flechas de má qualidade. Não
existiam flechas de verdade. Foi por ele que os Yanomami se flecharam,
pois ele as fez. É o dono mesmo. 

 

\chapterspecial{{Xereka} a {rë} {kuprarionowei}}{}{}
 

 

 

Xereka a rë kuprarionowei. Weti naha të kupronomi! Kama pë teri a kua
yaro, hawë hikari a pata ha tapramarnɨ, xereka si rë kekenowei ai të
kuami. Kama pë teri pënɨ pë kuopë ha, a përɨopë ha, hawë ɨ̃naha si pata
kuoma, e të si pata rë kuprarionowei, ɨ̃hɨ tëhë të rë
piyëmaɨ kukenowehei, e si kuoma. 

Ɨnaha të kua, kama pë teri yai. Si rë taprarenowehei, kama a rë
përɨonowei, ɨ̃hɨ a wãha, Xororiakapëwë e si, ɨ̃hɨnɨ xereka a niaamapë e
si, si rë tapramarenowei, xoro ɨ̃hɨ të pë no uhutipɨ ihirupɨ yëɨ, mau u
hamɨ. 

Xororiakapëwënɨ si xereka taprarema. Hauya si tapramarema. Ɨhɨ iha si he
rë harenowehei, a piyëremahe, kihi ipa u rë para kiri, kihi të si pata
koro, ɨ̃hɨ hei ĩhete rë të si pata yamoo kurayoi, ɨ̃hɨ e si yai,
Yanomamɨ tënɨ mai! Ɨhɨ e si piyëremahe. Yanomamɨ të pënɨ xereka a
ponomihe, pei si piyëremahe, si ha piyërëhenɨ, si ha kearɨhenɨ, të pë
niayorayoma, hapa. 

Hapa të pë hõrimoma, xereka a ponomihe, ruhu masi pë wai xomi niaamahe
tɨrɨtɨrɨ të pë wai xĩro niaamahe mahemahõ, urihi hamɨ të si pë rë
kuprai, tomɨ si poko pë, yãxaamahe, kohere si poko pë hawë xereka pë rë
kure, ɨ̃hɨ të pë he õkawa yãxaaɨ no preomahe. Xereka pë kuami yaro.
Ɨhɨnɨ Yanomamɨ të pë niayopë, si taprarema, ɨ̃hɨ teri a yai.

 


 

 

 

\chapterspecial{{Antes} {do} {surgimento} {do} {terçado}}{}{}
 

\letra{Q}{uando não havia terçado}, quebravam o peito das
tartaruguinhas \emph{pirema}, rachavam pau e amarravam na fenda do pau
aquele peito de tartaruguinha, sofriam com esse tipo de ferramentas com
as quais abriam roças. Assim faziam no início.

Amarravam também peito de jabuti, derrubavam árvores com machados de
pedra, com pedras. Aquelas pretas. Procuravam e juntavam as pedras,
afiavam"-nas e derrubavam as árvores grandes. Com essas pedras amarradas
no pau. Depois de recuperarem todas as pedras, de amarrá"-las bem
fincadas, eles derrubavam as árvores grandes. Por onde eles moravam, por
onde eles habitavam, com a casca do peito das tartaruguinhas, eles
cortavam os esteios das casas. Assim que faziam.

 


\chapterspecial{{Sipara} a {rë} {kuprarionowei}}{}{}
 

 

Sipara a mao tëhë, mixiukëmɨ, misi pë parɨkɨ si ha karoahenɨ, pë
parɨkɨ si hãhopomahe, hãhoa kurenaha të pë hãhoaɨkuo no ha preohenɨ,
ɨ̃hɨ të pë parɨkɨ si ha të pë hikaripɨ taoma. Ɨnaha të pë kuaama, hapa. 

Totori parɨkɨ si hãhopomahe, poo maro pënɨ kayapa hi pë tuyëmahe, maa ma
pë, të pë rë ĩxii, ɨ̃hɨ të pë ha hokahenɨ, të pë namo ha tahenɨ, kayapa
hi pë tuyëmahe. Hãhoa të pënɨ. Pei të pë ha wãkɨahenɨ, të pë posi ha
õkaahenɨ, kayapa hi pë tuyëmahe, ɨ̃hɨ pei të pë përɨapë hamɨ xapono a
tapehe, ɨ̃hɨ misi pë parɨkɨ sinɨ, të pë hãtopɨ nahi pëoma. Ɨhɨ të xĩro
tamahe. 

 

 

\chapterspecial{O {corte} {dos} {cabelos}}{}{}
 

\letra{Q}{uando} não havia \emph{napë}, sofriam de ter o rosto fechado pelos
cabelos que desciam, tinham o rosto como o de mulher por causa dos
cabelos. Ele fez o bambu \emph{sunama} e o
bambu \emph{waharokoma} aparecerem\emph{.} Os Yanomami cortavam os
cabelos com ponta de tacuará. Quando não o encontravam, usavam o
bambu \emph{uhe}. Rasgavam"-no e cortavam os cabelos com isso, faziam o
corte com esses pedaços. Eles se davam esses pedaços de má qualidade,
pois não havia \emph{napë}. As mulheres sofriam com o sangue do corte,
quando faziam assim, cortavam a testa, como faziam assim, eles sofriam.
No início não havia tesoura. 

Qual é o \emph{napë} que apareceria e inventaria aquela tesoura?

No início, se cortavam mutuamente o cabelo com pedaços de tacuará
afiados. Partiam o bambu \emph{sunama}, com o qual se cortavam o cabelo
mutuamente, com o fio da lâmina. Kreti! Kreti! Kreti! Cortavam"-se o
cabelo mutuamente. Assim que faziam entre eles. Também não havia facão. 

Cortavam também a carne com pedaços de tacuará \emph{sunama}, no início.

 

\chapterspecial{{Të} {pë} {hemakasi} {pëyomou} {rë} {hapamonowei}}{}{}
 

 

Hapa napë a mao tëhë, pë mɨ raeke no preaama, të pë henakɨ itoma, të pë
mɨ raeke no preaama, suwë mohekɨ kurenaha, të pë mohekɨ kuaama, pei të
pë henakɨnɨ, ɨ̃hɨ të rë kui, Sunamau he kɨ rë pëtamarenowei, Waharokoma
kɨ rë pëtamarenowei, rahaka pë atahunɨ të pë mɨ pëoma. Ɨhɨ kɨ he hao mao
tëhë, uhe pë wãha yai kua kuhe. Ɨhɨ pënɨ, të pë ha kakahenɨ, të pë
mɨ tayoma, hõra, hanɨma, atahu pënɨ. Ɨhɨ të pë yãxaamahe, napë pë kuami
yaro, suwë të pë mɨ ĩyë no preaama, ɨnaha të pë taɨhe ha, huko si pë
hanɨɨ, ɨnaha të pë pata taɨhe yaro, ɨ̃hɨ të pë ha të pë no preaama. 

Hapa nakɨra pë kuonomi yaro, weti a napë a ha pëtarunɨ, kɨ taprapë?
Mɨ hanɨyou të kuoma hapa. Rahaka namo, rahakaa ãtahu të pë hanɨyoma.
Sunama akasi pë kakaɨ piyëohe, ɨ̃hɨ të pë tutakɨnɨ të pë mɨ hanɨyoma.
Kreti! Kreti! Kreti! Të pë henakɨ tayoma. Ɨnaha pë tayoma. Xokopi pë kãi
kuo mao tëhë, xokopi pë kãi kuonomi. 

Ɨhɨ Sunama akasi pënɨ të pë yaropɨ hanɨoma. Hapa. 

 



\pagebreak
\blankpage
%\blankAteven
\pagestyle{empty}
\begingroup
\fontsize{7}{8}\selectfont

{\large\textsc{coleção «hedra edições»}}

\begin{enumerate}
\setlength\parskip{4.2pt}
\setlength\itemsep{-1.4mm}
%\item \textit{Poemas da cabana montanhesa}, Saigy\=o
\item \textit{A arte da guerra}, Maquiavel
\item \textit{A conjuração de Catilina}, Salústio
\item \textit{A cruzada das crianças/ Vidas imaginárias}, Marcel Schwob
\item \textit{A filosofia na era trágica dos gregos}, Friedrich Nietzsche
\item \textit{A fábrica de robôs}, Karel Tchápek 
\item \textit{A história trágica do Doutor Fausto}, Christopher Marlowe
\item \textit{A metamorfose}, Franz Kafka
\item \textit{A monadologia e outros textos}, Gottfried Leibniz
\item \textit{A morte de Ivan Ilitch}, Lev Tolstói 
\item \textit{A velha Izerguil e outros contos}, Maksim Górki
\item \textit{A vida é sonho}, Calderón de la Barca
\item \textit{A volta do parafuso}, Henry James
\item \textit{A voz dos botequins e outros poemas}, Paul Verlaine 
\item \textit{A vênus das peles}, Leopold von Sacher{}-Masoch
\item \textit{A última folha e outros contos}, O.\,Henry
\item \textit{Americanismo e fordismo}, Antonio Gramsci
\item \textit{Anarquia pela educação}, Élisée Reclus 
\item \textit{Apologia de Galileu}, Tommaso Campanella 
\item \textit{Arcana C\oe lestia} e \textit{Apocalipsis revelata}, Emanuel Swedenborg
\item \textit{As bacantes}, Eurípides
\item \textit{Autobiografia de uma pulga}, [Stanislas de Rhodes]
\item \textit{Ação direta e outros escritos}, Voltairine de Cleyre
\item \textit{Balada dos enforcados e outros poemas}, François Villon
\item \textit{Carmilla, a vampira de Karnstein}, Sheridan Le Fanu
\item \textit{Carta sobre a tolerância}, John Locke
\item \textit{Contos clássicos de vampiro}, L.\,Byron, B.\,Stoker \& outros
\item \textit{Contos de amor, de loucura e de morte}, Horacio Quiroga
\item \textit{Contos indianos}, Stéphane Mallarmé
\item \textit{Cultura estética e liberdade}, Friedrich von Schiller
\item \textit{Cântico dos cânticos}, [Salomão]
\item \textit{Dao De Jing}, Lao Zi
\item \textit{Discursos ímpios}, Marquês de Sade
\item \textit{Dissertação sobre as paixões}, David Hume
\item \textit{Diário de um escritor (1873)}, Fiódor Dostoiévski
\item \textit{Diário parisiense e outros escritos}, Walter Benjamin
\item \textit{Diários de Adão e Eva}, Mark Twain
\item \textit{Don Juan}, Molière
\item \textit{Dos novos sistemas na arte}, Kazimir Maliévitch
\item \textit{Educação e sociologia}, Émile Durkheim
\item \textit{Édipo Rei}, Sófocles
\item \textit{Elogio da loucura}, Erasmo de Rotterdam
\item \textit{Émile e Sophie ou os solitários}, Jean-Jacques Rousseau 
\item \textit{Emília Galotti}, Gotthold Ephraim Lessing
\item \textit{Entre camponeses}, Errico Malatesta
\item \textit{Ernestine ou o nascimento do amor}, Stendhal
\item \textit{Escritos revolucionários}, Errico Malatesta
\item \textit{Escritos sobre arte}, Charles Baudelaire
\item \textit{Escritos sobre literatura}, Sigmund Freud
\item \textit{Eu acuso!}, Zola/\,\textit{O processo do capitão Dreyfus}, Rui Barbosa
\item \textit{Explosão: romance da etnologia}, Hubert Fichte
\item \textit{Fedro}, Platão
\item \textit{Feitiço de amor e outros contos}, Ludwig Tieck
\item \textit{Flossie, a Vênus de quinze anos}, [Swinburne]
\item \textit{Fábula de Polifemo e Galateia e outros poemas}, Góngora
\item \textit{Fé e saber}, Georg W.\,F.\,Hegel
\item \textit{Gente de Hemsö}, August Strindberg 
\item \textit{Hawthorne e seus musgos}, Melville
\item \textit{Hino a Afrodite e outros poemas}, Safo de Lesbos 
\item \textit{História da anarquia (vol.\,\textsc{ii})}, Max Nettlau
\item \textit{História da anarquia (vol.\,\textsc{i})}, Max Nettlau
\item \textit{Imitação de Cristo}, Tomás de Kempis
\item \textit{Incidentes da vida de uma escrava}, Harriet Jacobs
\item \textit{Inferno}, August Strindberg
\item \textit{Investigação sobre o entendimento humano}, David Hume
\item \textit{Jazz rural}, Mário de Andrade
\item \textit{Jerusalém}, William Blake
\item \textit{Joana d'Arc}, Jules Michelet
\item \textit{Lira gregra}, Giuliana Ragusa (org.)
\item \textit{Lisístrata}, Aristófanes 
\item \textit{Ludwig Feuerbach e o fim da filosofia clássica alemã}, Friederich Engels
\item \textit{Manifesto comunista}, Karl Marx e Friederich Engels
\item \textit{Memórias do subsolo}, Fiódor Dostoiévski
\item \textit{Metamorfoses}, Ovídio
\item \textit{Micromegas e outros contos}, Voltaire
\item \textit{Narrativa de William W.\,Brown, escravo fugitivo}, William Wells Brown
\item \textit{Nascidos na escravidão: depoimentos norte-americanos}, \textsc{wpa}
\item \textit{No coração das trevas}, Joseph Conrad
\item \textit{Noites egípcias e outros contos}, Aleksandr Púchkin
\item \textit{O casamento do Céu e do Inferno}, William Blake
\item \textit{O cego e outros contos}, \textsc{d.\,h}.\,Lawrence
\item \textit{O chamado de Cthulhu}, \textsc{h.\,p.}\,lovecraft
\item \textit{O contador de histórias e outros textos}, Walter Benjamin
\item \textit{O corno de si próprio e outros contos}, Marquês de Sade
\item \textit{O destino do erudito}, Johann Fichte
\item \textit{O estranho caso do dr.\,Jekyll e Mr. Hyde}, Robert Louis Stevenson
\item \textit{O fim do ciúme e outros contos}, Marcel Proust
\item \textit{O indivíduo, a sociedade e o Estado, e outros ensaios}, Emma Goldman
\item \textit{O ladrão honesto e outros contos}, Fiódor Dostoiévski
\item \textit{O livro de Monelle}, Marcel Schwob
\item \textit{O mundo ou tratado da luz}, René Descartes
\item \textit{O novo Epicuro: as delícias do sexo}, Edward Sellon
\item \textit{O pequeno Zacarias, chamado Cinábrio}, \textsc{e.\,t.\,a.}\,Hoffmann
\item \textit{O primeiro Hamlet}, William Shakespeare
\item \textit{O princípio anarquista e outros ensaios}, Piotr Kropotkin
\item \textit{O princípio do Estado e outros ensaios}, Mikhail Bakunin
\item \textit{O príncipe}, Maquiavel
\item \textit{O que eu vi, o que nós veremos}, Santos-Dumont
\item \textit{O retrato de Dorian Gray}, Oscar Wilde
\item \textit{O sobrinho de Rameau}, Diderot
\item \textit{Ode ao Vento Oeste e outros poemas}, \textsc{p.\,b.}\,Shelley
\item \textit{Ode sobre a melancolia e outros poemas}, John Keats
\item \textit{Odisseia}, Homero
\item \textit{Oliver Twist}, Charles Dickens
\item \textit{Origem do drama barroco}, Walter Benjamin
\item \textit{Os sofrimentos do jovem Werther}, Goethe
\item \textit{Os sovietes traídos pelos bolcheviques}, Rudolf Rocker
\item \textit{Para serem lidas à noite}, Ion Minulescu
\item \textit{Pensamento político de Maquiavel}, Johann Fichte
\item \textit{Pequeno-burgueses}, Maksim Górki
\item \textit{Pequenos poemas em prosa}, Charles Baudelaire
\item \textit{Perversão: a forma erótica do ódio}, Robert Stoller
\item \textit{Poemas}, Lord Byron
\item \textit{Poesia basca: das origens à Guerra Civil} 
\item \textit{Poesia catalã: das origens à Guerra Civil} 
\item \textit{Poesia espanhola: das origens à Guerra Civil} 
\item \textit{Poesia galega: das origens à Guerra Civil} 
\item \textit{Pr\ae terita}, John Ruskin
\item \textit{Primeiro livro dos Amores}, Ovídio
\item \textit{Rashômon e outros contos}, Ryūnosuke Akutagawa
\item \textit{Revolução e liberdade: cartas de 1845 a 1875}, Mikhail Bakunin
\item \textit{Robinson Crusoé}, Daniel Defoe
\item \textit{Romanceiro cigano}, Federico García Lorca
\item \textit{Sagas}, August Strindberg
\item \textit{Sobre a amizade e outros diálogos}, Jorge Luis Borges e Osvaldo Ferrari
\item \textit{Sobre a filosofia e outros diálogos}, Jorge Luis Borges e Osvaldo Ferrari
\item \textit{Sobre a filosofia e seu método (Parerga e paralipomena)} (v.\textsc{ii}, t.\textsc{i}), Arthur Schopenhauer 
\item \textit{Sobre a liberdade}, Stuart Mill
\item \textit{Sobre a utilidade e a desvantagem da histório para a vida}, Friedrich Nietzsche
\item \textit{Sobre a ética (Parerga e paralipomena)} (v.\textsc{ii}, t.\textsc{ii}), Arthur Schopenhauer 
\item \textit{Sobre anarquismo, sexo e casamento}, Emma Goldman
\item \textit{Sobre o riso e a loucura}, [Hipócrates]
\item \textit{Sobre os sonhos e outros diálogos}, Jorge Luis Borges e Osvaldo Ferrari
\item \textit{Sobre verdade e mentira}, Friedrich Nietzsche
\item \textit{Sonetos}, William Shakespeare
\item \textit{Sátiras, fábulas, aforismos e profecias}, Leonardo da Vinci
\item \textit{Teleny, ou o reverso da medalha}, Oscar Wilde
\item \textit{Teogonia}, Hesíodo
\item \textit{Trabalhos e dias}, Hesíodo
\item \textit{Triunfos}, Petrarca
\item \textit{Um anarquista e outros contos}, Joseph Conrad
\item \textit{Viagem aos Estados Unidos}, Alexis de Tocqueville
\item \textit{Viagem em volta do meu quarto}, Xavier de Maistre 
\item \textit{Viagem sentimental}, Laurence Sterne
\end{enumerate}

{\large\textsc{coleção «metabiblioteca»}}

\begin{enumerate}
\setlength\parskip{4.2pt}
\setlength\itemsep{-1.4mm}
\item \textit{A carteira de meu tio}, Joaquim Manuel de Macedo
\item \textit{A cidade e as serras}, Eça de Queirós
\item \textit{A escrava}, Maria Firmina dos Reis
\item \textit{A família Medeiros}, Júlia Lopes de Almeida 
\item \textit{A pele do lobo e outras peças}, Artur Azevedo
\item \textit{Auto da barca do Inferno}, Gil Vicente
\item \textit{Bom Crioulo}, Adolfo Caminha
\item \textit{Cartas a favor da escravidão}, José de Alencar
\item \textit{Contos e novelas}, Júlia Lopes de Almeida
\item \textit{Crime}, Luiz Gama
\item \textit{Democracia}, Luiz Gama
\item \textit{Direito}, Luiz Gama
\item \textit{Elixir do pajé: poemas de humor, sátira e escatologia}, Bernardo Guimarães
\item \textit{Eu}, Augusto dos Anjos
\item \textit{Farsa de Inês Pereira}, Gil Vicente
\item \textit{Helianto}, Orides Fontela
\item \textit{História da província Santa Cruz}, Gandavo
\item \textit{Iracema}, José de Alencar
\item \textit{Liberdade}, Luiz Gama
\item \textit{Mensagem}, Fernando Pessoa
\item \textit{Nós, os sul-americanos}, Flávio de Carvalho
\item \textit{O Ateneu}, Raul Pompeia
\item \textit{O cortiço}, Aluísio Azevedo
\item \textit{O desertor}, Silva Alvarenga
\item \textit{Oração aos moços}, Rui Barbosa
\item \textit{Pai contra mãe e outros contos}, Machado de Assis
\item \textit{Poemas completos de Alberto Caeiro}, Fernando Pessoa
\item \textit{Teatro de êxtase}, Fernando Pessoa
\item \textit{Transposição}, Orides Fontela
\item \textit{Tratado descritivo do Brasil em 1587}, Gabriel Soares de Sousa
\item \textit{Tratados da terra e gente do Brasil}, Fernão Cardim 
\item \textit{Utopia Brasil}, Darcy Ribeiro
\item \textit{Índice das coisas mais notáveis}, Antônio Vieira
\end{enumerate}

\medskip
{\large\textsc{coleção «que horas são?»}}

\begin{enumerate}
\setlength\parskip{4.2pt}
\setlength\itemsep{-1.4mm}
\item \textit{8/1: A rebelião dos manés}, Pedro Fiori Arantes, Fernando Frias e Maria Luiza Meneses
\item \textit{A linguagem fascista}, Carlos Piovezani \& Emilio Gentile
\item \textit{A sociedade de controle}, J.\,Souza; R.\,Avelino; S.\,Amadeu (orgs.)
\item \textit{Ativismo digital hoje}, R.\,Segurado; C.\,Penteado; S.\,Amadeu (orgs.)
\item \textit{Crédito à morte}, Anselm Jappe
\item \textit{Descobrindo o Islã no Brasil}, Karla Lima
\item \textit{Desinformação e democracia}, Rosemary Segurado
\item \textit{Dilma Rousseff e o ódio político}, Tales Ab'Sáber
\item \textit{Labirintos do fascismo} (v.\textsc{iii}), João Bernardo
\item \textit{Labirintos do fascismo} (v.\textsc{ii}), João Bernardo
\item \textit{Labirintos do fascismo} (v.\textsc{iv}), João Bernardo
\item \textit{Labirintos do fascismo} (v.\textsc{i}), João Bernardo
\item \textit{Labirintos do fascismo} (v.\textsc{vi}), João Bernardo
\item \textit{Labirintos do fascismo} (v.\textsc{v}), João Bernardo
\item \textit{Lugar de negro, lugar de branco?}, Douglas Rodrigues Barros
\item \textit{Lulismo, carisma pop e cultura anticrítica}, Tales Ab'Sáber
\item \textit{Machismo, racismo, capitalismo identitário}, Pablo Polese
\item \textit{Michel Temer e o fascismo comum}, Tales Ab'Sáber
\item \textit{O quarto poder: uma outra história}, Paulo Henrique Amorim
\item \textit{Universidade, cidade e cidadania}, Franklin Leopoldo e Silva
\end{enumerate}

\medskip
{\large\textsc{coleção «mundo indígena»}}

\begin{enumerate}
\setlength\parskip{4.2pt}
\setlength\itemsep{-1.4mm}
\item \textit{A folha divina}, Timóteo Verá Tupã Popygua
\item \textit{A mulher que virou tatu}, Eliane Camargo
\item \textit{A terra uma só}, Timóteo Verá Tupã Popygua
\item \textit{A árvore dos cantos}, Pajés Parahiteri
\item \textit{Cantos dos animais primordiais}, Ava Ñomoandyja Atanásio Teixeira
\item \textit{Crônicas de caça e criação}, Uirá Garcia
\item \textit{Círculos de coca e fumaça}, Danilo Paiva Ramos
\item \textit{Folhas divinas}, Timóteo Verá Tupã Popygua
\item \textit{Nas redes guarani}, Valéria Macedo \& Dominique Tilkin-Gallois
\item \textit{Não havia mais homens}, Luciana Storto
\item \textit{O surgimento da noite}, Pajés Parahiteri
\item \textit{O surgimento dos pássaros}, Pajés Parahiteri
\item \textit{Os Aruaques}, Max Schmidt
\item \textit{Os cantos do homem-sombra}, Patience Epps e Danilo Paiva Ramos
\item \textit{Os comedores de terra}, Pajés Parahiteri
\item \textit{Xamanismos ameríndios}, A.\,Barcelos Neto, L.\,Pérez Gil \& D.\,Paiva Ramos
\end{enumerate}

\medskip
{\large\textsc{coleção «ecopolítica»}}

\begin{enumerate}
\setlength\parskip{4.2pt}
\setlength\itemsep{-1.4mm}
\item \textit{Anarquistas na América do Sul}, E.\,Passetti, S.\,Gallo; A.\,Augusto  (orgs.)
\item \textit{Ecopolítica}, E.\,Passetti; A.\,Augusto; B.\,Carneiro; S.\,Oliveira, T.\,Rodrigues  (orgs.)
\item \textit{Pandemia e anarquia}, E.\,Passetti; J.\,da Mata; J.\,Ferreira  (orgs.)
\end{enumerate}

% \medskip
% {\large\textsc{coleção <<anarc>>}}

% \begin{enumerate}
% \setlength\parskip{4.2pt}
% \setlength\itemsep{-1.4mm}
% \item \textit{Anarquia pela educação}, Élisée Reclus 
% \item \textit{Ação direta e outros escritos}, Voltairine de Cleyre
% \item \textit{Entre camponeses}, Errico Malatesta
% \item \textit{Escritos revolucionários}, Errico Malatesta
% \item \textit{História da anarquia (vol.\,\textsc{ii})}, Max Nettlau
% \item \textit{História da anarquia (vol.\,\textsc{i})}, Max Nettlau
% \item \textit{O indivíduo, a sociedade e o Estado, e outros ensaios}, Emma Goldman
% \item \textit{O princípio anarquista e outros ensaios}, Piotr Kropotkin
% \item \textit{O princípio do Estado e outros ensaios}, Mikhail Bakunin
% \item \textit{Os sovietes traídos pelos bolcheviques}, Rudolf Rocker
% \item \textit{Revolução e liberdade: cartas de 1845 a 1875}, Mikhail Bakunin
% \item \textit{Sobre anarquismo, sexo e casamento}, Emma Goldman
% \end{enumerate}

% \medskip
% {\large\textsc{coleção <<narrativas da escravidão>>}}

% \begin{enumerate}
% \setlength\parskip{4.2pt}
% \setlength\itemsep{-1.4mm}
% \item \textit{Incidentes da vida de uma escrava}, Harriet Jacobs
% \item \textit{Narrativa de William W.\,Brown, escravo fugitivo}, William Wells Brown
% \item \textit{Nascidos na escravidão: depoimentos norte-americanos}, \textsc{wpa}
% \end{enumerate}

\pagebreak	   % [lista de livros publicados]
\blankpage
\pagebreak

\ifodd\thepage\blankpage\fi

\parindent=0pt
\footnotesize\thispagestyle{empty}

% \noindent\textbf{Dados Internacionais de Catalogação na Publicação -- CIP}\\
% \noindent\textbf{(Câmara Brasileira do Livro, SP, Brasil)}\\

% \dotfill\\

% \hspace{20pt}ISBN 978-65-86238-31-0 (Livro do Estudante)

% \hspace{20pt}ISBN 978-65-86238-30-3 (Manual do Professor)\\[6pt]

% \hspace{20pt}\parbox{190pt}{1. Crônicas Brasileira. 2. Contos Brasileiro. 3. Rosa, Alexandre. I. Título.}\\[6pt]

% \hspace{188pt}\textsc{cdd}-B869.8

% \dotfill

% \noindent{}Elaborado por Regina Célia Paiva da Silva CRB -- 1051\\

\mbox{}\vfill
\begin{center}
		\begin{minipage}{.7\textwidth}\tiny\noindent{}
		\centering\tiny
		Adverte-se aos curiosos que se imprimiu este livro na gráfica Eskenazi, na data de \today, em papel Pólen Bold 90, composto em tipologia Minion Pro, em 11 pontos, com diversos sofwares livres,
		dentre eles Lua\LaTeX e git.\\ 
		\ifdef{\RevisionInfo{}}{\par(v.\,\RevisionInfo)}{}\medskip\\\
		\adforn{64}
		\end{minipage}
\end{center}		   % [colofon]


\checkandfixthelayout
\end{document}
